% Latin-script Interslavic source component: Work 44 (Deuring lecture notes).
% Audited editorial provenance: ../00_SESSION_INDEPENDENT_STATE_v019.json and ../NORMALIZATION_DECISIONS_v019.jsonl.
\documentclass[11pt]{article}
\usepackage{fontspec}
\IfFontExistsTF{Times New Roman}{\setmainfont{Times New Roman}}{\setmainfont{TeX Gyre Termes}}
\usepackage{amsmath,amssymb}
\usepackage{geometry}
\usepackage{microtype}
\usepackage{hyperref}
\usepackage{mathrsfs}
\newcommand{\srcfn}[2]{\footnote{#2}}
\providecommand{\tightlist}{\setlength{\itemsep}{0pt}\setlength{\parskip}{0pt}}
\allowdisplaybreaks

\geometry{a4paper,margin=2.1cm}
\setlength{\parindent}{1.1em}
\setlength{\parskip}{0pt}
\emergencystretch=10em
\sloppy
\hbadness=10000
\raggedbottom
\hypersetup{hidelinks}
\newcommand{\foreign}[1]{#1}
\newcommand{\tocline}[2]{\noindent #1\dotfill #2\par}
\newcommand{\tocsec}[3]{\noindent\hspace*{1.2em}\S\,#1.\quad #2\dotfill #3\par}

\begin{document}
\thispagestyle{empty}

\begin{center}
{\LARGE\bfseries Алгебра хиперкомплексных величин\par}
\vspace{3.5em}
{\large Лекција проф. Е. \foreign{Noether}\par}
{\large Зимны семестр 1929/30\par}
\vspace{1.5em}
{\large Обработал проф. M. \foreign{Deuring}\par}
\vspace{4em}
{\Large\bfseries Содржање\par}
\end{center}
\vspace{1em}

\begingroup
\small
\setlength{\parindent}{0pt}
\tocline{\textbf{Увод}}{3}
\medskip
\tocline{\textbf{Глава I: Прєдставјења}}{3}
\tocsec{1}{Дефиниција директного прєдставјења --- дефиниција реципрочного прєдставјења}{3}
\tocsec{2}{Класы прєдставјениј}{4}
\tocsec{3}{Модулы прєдставјениј}{4}
\tocsec{4}{Свез меджу модулами прєдставјениј и прєдставјењами}{5}
\medskip
\tocline{\textbf{Глава II: Галоисова теорија комутативных пољ}}{7}
\tocsec{5}{Разширјење коефициентного поља хиперкомплексных системов}{7}
\tocsec{6}{Иредуцибилне прєдставјења комутативных системов}{8}
\tocsec{7}{Изоморфизмы поља}{8}
\tocsec{8}{Разпадно поље}{9}
\tocsec{9}{Изоморфизмы од $Z_1$ на $Z_i$}{9}
\tocsec{10}{Галоисова група}{9}
\tocsec{11}{Главны теорем Галоисовој теорије}{9}
\tocsec{12}{Формално значење компонентов $e_i$}{11}
\tocsec{13}{Обчи теорем о продолжењу}{12}
\medskip
\tocline{\textbf{Глава III: Абелове групы}}{13}
\tocsec{14}{Групово прстен}{13}
\tocsec{15}{Групове прстены абеловых груп}{14}
\tocsec{16}{Релације карактеров}{15}
\tocsec{17}{Галоисова теорија абеловых груп}{15}
\medskip
\tocline{\textbf{Глава IV: Двустранно просте прстены}}{18}
\tocsec{18}{Помочна лема}{18}
\tocsec{19}{Прєдставјења двустранно простых хиперкомплексных системов в некомутативных тєлах, кторе разширјајут јих коефициентне поља}{19}
\tocsec{20}{Некомутативне тєла}{22}
\tocsec{21}{Галоисова теорија некомутативных тєл}{26}
\clearpage
\tocline{\textbf{Глава V: Факторне системы}}{29}
\tocsec{22}{Група тєл с даным центром}{29}
\tocsec{23}{Факторне системы}{30}
\tocsec{24}{Множење факторных системов}{33}
\tocsec{25}{Нормално прєдставјење $\mathfrak K_r$ с Галоисовыми максималными комутативными подпољами}{39}
\tocsec{26}{Множење скрєженых прєдставјениј}{43}
\medskip
\tocline{\textbf{Глава VI: Теорија скрєженых продуктов}}{44}
\tocsec{27}{Прєдставјења $\mathfrak R_r$ како скрєжене продукты}{44}
\tocsec{28}{Продуктны теорем за факторне системы}{47}
\tocsec{29}{Теорем о главном роду в минималном}{50}
\tocsec{30}{Специализација на цикличне разпадне поља}{50}
\tocsec{31}{Примєњења цикличного специалного случаја}{52}
\endgroup

\clearpage
\section*{Увод}\label{einleitung}

Та лекција обрабја обчу теорију прєдставјениј прстенов, ктора взникла с
једној страны из теорије прєдставјениј конечных груп, а с другој страны
из алгебраичној и аритметичној теорије хиперкомплексных числовых
системов. То обчејше разумєње теорије прєдставјениј имаје слєдне
прєдности над досєгашњим приступом (\foreign{Frobenius},
\foreign{Schur}):

Зато же разгледамо не саму групу, але групово прстен, а обчејше
хиперкомплексны систем или либовољно прстен, можно примєнити
теорију прстенов и идеалов и тым освободити теорију од многых нејасных
изчислениј. То примєњење обчеј теорије прстенов и идеалов не само
упрошчаје теорију, але такоже води ју даље --- напр. в пытању о
одношењах групы с јеј подгрупами --- и отвара јеј нове области
примєњења. Помоцју обчеј теорије прєдставјениј можно дати ново, врло
елегантно основање Галоисовој теорији и --- можда јешче важнєјше ---
установити Галоисову теорију некомутативных тєл.

Основне појмы теорије груп, прстенов и идеалов и обче теоремы из теорије
прєдставјениј можно најдти в дєлу господичны \foreign{Noether}
\foreign{Hyperkomplexe Größen und Darstellungstheorie},
\foreign{Mathematische Zeitschrift} 30, стр.~641. То дєло, цитовано
како \foreign{M.Z.}, трєба разгледати дополњењем к тому изложењу.
Исты материал јест такоже в записах лекције господичны
\foreign{Noether} \foreign{Über Hyperkomplexe Größen und
Gruppencharaktere}, зимны семестр 1927/28.

Појмы прєдставјење и модул прєдставјења сут ту кратко објасњене,
да бы ввести реципрочны модул прєдставјења и реципрочна
прєдставјења, кторе сут нове само формално.

\section*{Глава I. Прєдставјења}\label{kapitel-i.-darstellungen}

$\mathfrak{o}$ нехај буде прстен --- прстен, кторе се прєдставја --- а
$\mathsf T$ друго прстен --- то, в ктором се прєдставја.

% BEGIN BOOK_S01
\subsection*{§ 1. Дефиниција прємого прєдставјења}\label{definition-der-direkten-darstellung}

\emph{Прємо прєдставјење} \(n\)-того ступења од \(\mathfrak{o}\) в \(\mathsf T\) јест подпрстен \(\mathfrak D\) прстена \(\mathsf T_n\) матриц \(n\)-того ступења с елементами из \(\mathsf T\), на кторо јест \(\mathfrak{o}\) образовано прстеновым хомоморфизмом. В знаках

\[
\mathfrak{o}\to\mathfrak D\subseteq\mathsf T_n.
\]

\subsection*{Дефиниција реципрочного прєдставјења}\label{definition-der-reziproken-darstellung}

\emph{Реципрочно прєдставјење} \(n\)-того ступења од \(\mathfrak{o}\) в \(\mathsf T^*\) јест подпрстен \(\mathfrak D^*\) прстена \(\mathsf T_n^*\) матриц \(n\)-того ступења с елементами из \(\mathsf T^*\), на кторо јест \(\mathfrak{o}\) образовано реципрочным прстеновым хомоморфизмом. Реципрочны прстеновы хомоморфизм значи слєдујуче: ако елементам из \(\mathfrak{o}\), \(c,d\), одповєдајут елементи из \(\mathfrak D^*\), \(c^*,d^*\), то елементу из \(\mathfrak{o}\), \(cd\), трєба одповєдати елемент из \(\mathfrak D^*\), \(d^*\cdot c^*\), а елементу из \(\mathfrak{o}\), \(c+d\), елемент из \(\mathfrak D^*\), \(c^*+d^*\). В знаках

\[
\mathfrak{o}\nrightarrow\mathfrak D^*\subseteq\mathsf T_n^*.
\]

Како примєр појма реципрочного хомоморфизма: ако матрицам реда \(n\) над комутативным пољем приредити јих транспоноване матрице, то приреджење јест реципрочны изоморфизм прстенов.

Обче можно к либовољному прстену \(\mathsf T\) конструовати реципрочно изоморфно прстен \(\mathsf T^*\) слєдујучим начином:

Всакому елементу из \(\mathsf T\), \(c\), приредимо симбол \(c^*\) и дефинирујемо сложење тых симболов формулом \(c^*+d^*=(c+d)^*\), а множење формулом

\[
c^*\cdot d^*=(dc)^*.
\]

Множство симболов \(c^*\) јест прстен, реципрочно изоморфно с \(\mathsf T\), и означујемо је \(\mathsf T^*\).
% END BOOK_S01

% BEGIN BOOK_S02
\subsection*{§ 2. Класы прєдставјениј}\label{darstellungsklassen}

Прємо прєдставјење \(\mathfrak D\) од \(\mathfrak{o}\) в \(\mathsf T\) можно прєтворити в изоморфно прєдставјење \(\mathfrak D'\) тако, же се оно по елементах трансформује регуларној матрицоју \(P\), т.~ј. елементу из \(\mathfrak{o}\), \(c\), вмєсто \(C\subseteq\mathfrak D\) приреди се \(P^{-1}CP\subseteq P^{-1}\mathfrak D P\).

Два прєдставјења, кторе можно тако трансформовати једно в друго, называјут се \emph{еквивалентными}.

Вса прєдставјења, еквивалентне једному прєдставјењу, творет \emph{класу прєдставјениј} (к. п.).

Тако само дефинујеме еквивалентност реципрочных прєдставјениј; два реципрочна прєдставјења называјут се еквивалентными, ако једно можно достати из другого трансформацијоју регуларној матрицоју. Вса прєдставјења, еквивалентне једному установјеному реципрочному прєдставјењу, творет \emph{класу реципрочных прєдставјениј}.
% END BOOK_S02

% BEGIN BOOK_S03
\subsection*{\texorpdfstring{\(\S\) 3. Модулы прєдставјениј}{\textbackslash С 3. Модулы прєдставјениј}}\label{s-3.-darstellungsmoduln}

Дефиниција \emph{прємого модула прєдставјења} (п. м. п.).

Прємы модул прєдставјења од \(\mathfrak{o}\) по одношењу к \(\mathsf T\) јест всакы двојны \(\mathfrak{o}\)-\(\mathsf T\)-модул \(\mathfrak M\) (т.~ј. адитивно записана Абелова група, ктора јест лєвым модулом по одношењу к \(\mathfrak{o}\), правым модулом по одношењу к \(\mathsf T\) и изполња смјешаны закон асоциативности: \(c\in\mathfrak{o}\), \(m\in\mathfrak M\), \(\tau\in\mathsf T\); \((cm)\tau=c(m\tau)\)), кторы можно записати како прєму суму конечного числа једночленных \(\mathsf T\)-модулов --- причем \(\tau=0\) слєдује из \(x_i\tau=0\) ---

\[
\mathfrak M=x_1\cdot\mathsf T+x_2\cdot\mathsf T+\cdots+x_n\cdot\mathsf T.
\]

Дефиниција \emph{реципрочного модула прєдставјења} (р. м. п.).

Реципрочны модул прєдставјења од \(\mathfrak{o}\) по одношењу к \(\mathsf T^*\) јест всакы лєвы \(\mathfrak{o}\)-модул и лєвы \(\mathsf T^*\)-модул \(\mathfrak M^*\), кторы изполња закон комутативности (ако \(c\in\mathfrak{o}\), \(m\in\mathfrak M^*\), \(\tau\in\mathsf T^*\), то \(c(\tau^*m^*)=\tau^*(cm)\)) и кторы можно записати како прєму суму конечного числа простых \(\mathsf T^*\)-модулов --- причем \(\tau^*=0\) слєдује из \(\tau^*\cdot x_i^*=0\) ---

\[
\mathfrak M^*=\mathsf T^*\cdot x_1+\cdots+\mathsf T^*\cdot x_n.
\]
% END BOOK_S03

% BEGIN BOOK_S04
\subsection*{§ 4. Свез меджу модулами прєдставјениј и прєдставјењами}\label{der-zusammenhang-zwischen-darstellungsmoduln-und-darstellungen}

\textbf{1.} \emph{Всакы модул прєдставјења једнозначно води к класе прєдставјениј (п. м. п. к прємој класе прєдставјениј, р. м. п. к реципрочној класе прєдставјениј).}

И то слєдујучим начином:

Множење прємого модула прєдставјења \(\mathfrak M\) елементом из \(\mathfrak{o}\), \(c\), значи хомоморфизм од \(\mathfrak M\) на себе самого, кторы заради закона асоциативности \((cm)\tau=c(m\tau)\) допушча елементы из \(\mathsf T\) како операторы. Оно јест зато \emph{линеарна трансформација} од \(\mathfrak M\) в себе, т.~ј. приреджење \(m\mapsto cm=m'\) имаје својство:

\[
\text{Ако } y_i\mapsto y_i'(=cy_i) \text{ јест, то такоже }
\sum y_i\tau_i\mapsto \sum y_i'\tau_i\left(=c\sum y_i\tau_i\right).
\]

Дефинујмо суму \(\mathfrak T_1+\mathfrak T_2\) двох линеарных трансформациј

\[
\mathfrak T_1:\ m\mapsto m'(=c_1m)
\quad\text{и}\quad
\mathfrak T_2:\ m\mapsto m''(=c_2m)
\]

како трансформацију \(m\mapsto m'+m''(=(c_1+c_2)m)\), а јих продукт \(\mathfrak T_1\cdot\mathfrak T_2\) како сложену трансформацију \(m\mapsto (m'')'=c_1c_2m\). Тогда множство различных линеарных трансформациј од \(\mathfrak M\), кторе стварјајут елементи из \(\mathfrak{o}\), јест прстеново хомоморфны образ \(\overline{\mathfrak D}\) од \(\mathfrak{o}\).

Нынє изберимо либовољну базу \(x_1,\ldots,x_n\) модула \(\mathfrak M\). Линеарну трансформацију, створјену елементом \(c\), од \(\mathfrak M\) очевидно јест достаточно задати тым, како се трансформујут базове елементи \(x_i\), то јест познањем матрице \((\gamma_{ij})=C\) система једначениј

\[
x'_j=cx_j=\sum x_i\gamma_{ij}
\]

зато же за друге елементы \(x=\sum x_j\tau_j\) просто јест \(x'=\sum x'_j\tau_j\). Те матрице \(C\) сут једнозначно приреджене различным линеарным трансформацијам (зато же једне опрєдєљајут друге) и, како се види непосрєдње, творет прстен изоморфно с \(\overline{\mathfrak D}\), то јест прстеново хомоморфны образ од \(\mathfrak{o}\), прєдставјење \(n\)-того ступња од \(\mathfrak{o}\).

Изоморфизм меджу \(\overline{\mathfrak D}\) и прєдставјењем содрживан јест в слєдујучих формулах: ако линеарној трансформацији, створјеној елементом \(c\in\mathfrak{o}\), одповєдаје матрица \(C\), то

\[
c(x_1,\ldots,x_n)=(x_1,\ldots,x_n)C
\]

и ако на једнакы начин \(d\in\mathfrak{o}\) и \(D\in\mathsf T_n\) одповєдајут једно другому

\[
d(x_1,\ldots,x_n)=(x_1,\ldots,x_n)D
\]

то очевидно јест

\[
\begin{aligned}
(c+d)(x_1,\ldots,x_n)
&=(cx_1+dx_1,\ldots,cx_n+dx_n)\\
&=(x_1,\ldots,x_n)(C+D)
\end{aligned}
\]

и

\[
\begin{aligned}
cd(x_1,\ldots,x_n)
&=c(dx_1,\ldots,dx_n)=c((x_1,\ldots,x_n)D)\\
&=(cx_1,\ldots,cx_n)D=(x_1,\ldots,x_n)CD.
\end{aligned}
\]

Тако линеарны трансформације, створјене елементами из \(\mathfrak{o}\), од \(\mathfrak M\) в себе при јих реализацији посрєдством установјене базы \(x_\nu\) модула \(\mathfrak M\) давајут прєдставјење \(n\)-того ступења од \(\mathfrak{o}\) в \(\mathsf T\).

Ако \((y_1,\ldots,y_n)\) јест друга база модула \(\mathfrak M\), она настава из \(x_1,\ldots,x_n\) множењем регуларној матрицоју

\[(y_1,\ldots,y_n)=(x_1,\ldots,x_n)P\]

Ако линеарнаја трансформација од \(\mathfrak M\) по одношењу к базе \(x_\nu\) јест реализована матрицоју \(C\), тогда по одношењу к базе \(y_\nu\) она јест реализована матрицоју \(P^{-1}CP\), зато же ако

\[c(x_1,\ldots,x_n)=(x_1,\ldots,x_n)C\]
то имаје мєсто

\[
\begin{aligned}
c(y_1,\ldots,y_n)
&=c(x_1,\ldots,x_n)P=((x_1,\ldots,x_n)C)P\\
&=(x_1,\ldots,x_n)CP=((y_1,\ldots,y_n)P^{-1})CP\\
&=(y_1,\ldots,y_n)P^{-1}CP.
\end{aligned}
\]

Из тога непосрєдње слєдује:

\emph{\(\mathfrak M\) даваје вса прєдставјења једној класы прєдставјениј, ако за реализацију линеарных трансформациј од \(\mathfrak M\), створјеных елементами из \(\mathfrak{o}\), употрєбимо все базы модула \(\mathfrak M\).}

\begin{enumerate}
\def\labelenumi{\arabic{enumi}.}
\setcounter{enumi}{1}
\tightlist
\item
  \emph{Всака класа прєдставјениј јест створјена модулом прєдставјења на начин указаны под 1.} (Поновно се ограничујемо на прєме прєдставјења.)
\end{enumerate}

\emph{Доказ.} Изберемо установјено прєдставјење из класы; нехај оно буде задано хомоморфным приреджењем \(c\mapsto C\). Ако \(n\) јест ступењ прєдставјења, то с \(n\) неизвєстными \(x_1,\ldots,x_n\) образујемо правы \(\mathsf T\)-модул

\[\mathfrak{M} = x_1 \mathsf{T} + \dots + x_n \mathsf{T}\]

(с обычными правилами рачунања \(\sum x_i\tau_i+\sum x_i\overline{\tau}_i=\sum x_i(\tau_i+\overline{\tau}_i)\), \(\sum x_i\tau_i=\sum x_i\overline{\tau}_i\) тогда и само тогда, ако \(\tau_i=\overline{\tau}_i\) за все \(i\), \((\sum x_i\tau_i)\varrho=\sum x_i\tau_i\varrho\)), и чинимо га лєвым \(\mathfrak{o}\)-модулом установјењем \(c\sum x_i\tau_i=\sum x_i'\tau_i\), ако \((x_1',\ldots,x_n')=(x_1,\ldots,x_n)C\), причем под \(C\) разумєјемо матрицу прєдставјења, приреджену елементу из \(\mathfrak{o}\), \(c\).

Тым установјењем \(\mathfrak M\) јест учињено двојным \(\mathfrak{o}\)-\(\mathsf T\)-модулом, зато же прво из релациј хомоморфизма слєдује

\[
c+d\mapsto C+D \quad \text{и} \quad cd\mapsto CD
\]

\[
(c+d)x_i=cx_i+dx_i
\]

\[
cd\cdot x_i=c\cdot dx_i,\quad\text{надаље јест}\quad
c\left(\sum x_i\tau_i+\sum x_i\overline{\tau}_i\right)
=c\sum x_i\tau_i+c\sum x_i\overline{\tau}_i,
\]

то сут својства лєвого модула; друго, очевидно јест

\[c(\sum x_i \, \tau_i \, \varrho) = (c \, \sum x_i \, \tau_i) \cdot \varrho \,,\]

то јест смјешаны закон асоциативности.

\(\mathfrak{M}\) јест зато модул прєдставјења, а употрєбјено прєдставјење јест ним створјено на начин указаны под 1., ако за реализацију његовых линеарных трансформациј употрєбимо базу \(x_1, \ldots, x_n\). Друге базы давајут остале прєдставјења те класы.

За реципрочна прєдставјења и модуле прєдставјениј все иде так само; добыва се, же модул

\[\mathfrak{M}^* = \mathsf{T}^* x_1^* + \dots + \mathsf{T}^* x_n^*\]

стварја реципрочно прєдставјење од \(\mathfrak{o}\) в \(\mathsf T^*\) тако, же елементу из \(\mathfrak{o}\), \(c\), приреджује матрицу \(C\) в једначењу

\[
c\begin{pmatrix}x_1\\ \vdots\\ x_n\end{pmatrix}
=\begin{pmatrix}cx_1\\ \vdots\\ cx_n\end{pmatrix}
=C\begin{pmatrix}x_1\\ \vdots\\ x_n\end{pmatrix}
\]

Различне базы поновно давајут различна прєдставјења једној класы.

\section*{Глава II. Галоисова теорија комутативных пољ}\label{kapitel-ii.-galoissche-theorie-kommutativer-koerper}

Досєга развијена теорија прєдставјениј јест достаточна за основање Галоисовој теорији, ктора имаје прєдност над обычној теоријеју, же не потрєбује специалне системы творечих елементов ни теоремы о полиномах неизвєстных.\srcfn{1}{За слєдујуче сравни M. З. §§ 15 --- 20. M. З. § 21 содрживаје прве теоремы слєдујучего оддєла о Галоисовој теорији. За дефиницију хиперкомплексных системов види M. З. § 6.}
% END BOOK_S04

% BEGIN BOOK_S05
\subsection*{§ 5. Разширјење коефициентов хиперкомплексных системов}\label{koeffizientenerweiterung-hyperkomplexer-systeme}

Формулујемо слєдујучи обчи теорем, чиј доказ оставјамо неизведеным, зато же он јест скоро саморазумны. Нехај \(\mathfrak{o}=c_1\mathsf P+\cdots+c_n\mathsf P\) буде хиперкомплексны систем над комутативным пољем \(\mathsf P\). Нехај \(\Omega\) буде разширјујуче прстен од \(\mathsf P\), чији центар содрживаје \(\mathsf P\). Дефинујемо \emph{разширјены систем} \(\mathfrak{o}_{\Omega}\) како множство всих формалных сум \(\sum c_i\omega_i\), \(\omega_i\in\Omega\), кторо најпрво установјењами

\[
\begin{array}{l}
\sum c_i\omega_i=\sum c_i\omega_i' \quad\text{тогда и само тогда, ако } \omega_i=\omega_i' \text{ за все } i;\\
\sum c_i\omega_i+\sum c_i\omega_i'=\sum c_i(\omega_i+\omega_i');\\
\omega\sum c_i\omega_i=\sum c_i\omega\omega_i;\quad
\sum c_i\omega_i\cdot\omega=\sum c_i\omega_i\omega;\\
\sum c_i\omega_i=\sum \omega_i c_i
\end{array}
\]

поставаје \(\Omega\)-модулом, комутативно свезаным с \(\Omega\), а даљшим установјењем

\[
\sum c_i\omega_i\cdot \sum c_i\omega_i
=\sum_i\sum_j c_ic_j\omega_i\omega_j
\]

--- причем \(c_ic_j\) јест уж дефинованы како елемент из \(\mathfrak{o}\) --- поставаје разширјујучим прстеном од \(\mathfrak{o}\). \(\mathfrak{o}_{\Omega}\) имаје ранг \(n\) над \(\Omega\) и можно га в формє

\[
\mathfrak{o}_{\Omega}=c_1\Omega+\cdots+c_n\Omega
\]

записати.

\emph{\(\mathfrak{o}_{\Omega}\) јест дефинованы независно од базы \(c_1,\ldots,c_n\). Ако \(\Omega\) јест комутативно поље, то \(\mathfrak{o}_{\Omega}\) јест хиперкомплексны систем над \(\Omega\).}
% END BOOK_S05

% BEGIN BOOK_S06
\subsection*{§ 6. Иредуцибилне прєдставјења комутативных системов}\label{die-irreduziblen-darstellungen-kommutativer-systeme}

Нехај \(\mathfrak Z=z_1\mathsf P+\cdots+z_n\mathsf P\) буде комутативны систем над \(\mathsf P\). Хочемо установити вса иредуцибилна прєдставјења од \(\mathfrak Z\) в алгебраично затворјеном разширјеном пољу \(\Omega\) од \(\mathsf P\). Зато же всако прєдставјење од \(\mathfrak Z\) в \(\Omega\) води к прєдставјењу од \(\mathfrak Z_\Omega=z_1\Omega+\cdots+z_n\Omega\) --- зато же прєдставјења базовых елементов, уж содрживаных в \(\mathsf P\), сут достаточне за установјење цєлокупного прєдставјења --- и обратно всако прєдставјење од \(\mathfrak Z_\Omega\) содрживаје прєдставјење од \(\mathfrak Z\), можно пытати такоже за прєдставјења од \(\mathfrak Z_\Omega\). Она сут по M. З. § 20 содрживана в регуларном прєдставјењу од \(\mathfrak Z_\Omega\), т.~ј. в прєдставјењу од \(\mathfrak Z_\Omega\), кторо настава, ако сам \(\mathfrak Z_\Omega\) употрєбимо како модул прєдставјења. Достаточно јест чак ограничити се на \(\mathfrak Z_\Omega/\mathfrak C\) како модул прєдставјења, причем под \(\mathfrak C\) разумєјемо радикал од \(\mathfrak Z_\Omega\). (види M. З. § 19).

Але \(\mathfrak Z_\Omega/\mathfrak C\) јест полно редуцибилны систем, то јест прєма сума пољ конечного ступења над \(\Omega\). Зато же \(\Omega\) јест алгебраично затворјено, те суманды сут изоморфне с \(\Omega\).

\[
\mathfrak Z_\Omega/\mathfrak C=e_1\Omega+\cdots+e_t\Omega.
\]

Елементи \(e_\nu\) сут компоненты јединице. зато имаје мєсто:

\emph{\(\mathfrak Z_\Omega\) имаје точно \(t\) различных клас иредуцибилных прєдставјениј в \(\Omega\), все првого ступења. При том \(t\) јест ранг од \(\mathfrak Z_\Omega/\mathfrak C\), зато највыше равны \(n\).}

Всака из тых клас прєдставјениј содрживаје само једно прєдставјење, зато же трансформација прєдставјења првого ступења в комутативном пољу не измєњаје прєдставјење: \(\pi^{-1}\cdot\alpha\cdot\pi=\alpha\).

Зато обстајајут точно \(t\) различне хомоморфизмы \(\Theta_1,\ldots,\Theta_t\) од \(\mathfrak Z_\Omega\) на подпрстены од \(\Omega\).
% END BOOK_S06

% BEGIN BOOK_S07
\subsection*{§ 7. Изоморфизмы поља}\label{die-isomorphismen-eines-koerpers}

Нынє нехај \(\mathfrak Z\) буде поље. Всакы хомоморфизм \(\Theta_i\) од \(\mathfrak Z_\Omega\) на подпрстен од \(\Omega\) мора изоморфно образовати \(\mathfrak Z\) на подпоље \(Z_i\) од \(\Omega\). Зато же или јест \(Z_i=0\) --- а то ту не јест случај, зато же јединице од \(\mathfrak Z\) одповєдаје јединица од \(\Omega\) --- или \(Z_i\) јест изоморфны образ од \(\mathfrak Z\). Зато:

\emph{Ако \(\mathfrak Z\) јест конечно разширјење \(n\)-того ступења од \(\mathsf P\), то в \(\Omega\) обстајајут толико различне поља, изоморфне с \(\mathfrak Z\), \(Z_i\), коликы јест ранг \(t\) од \(\mathfrak Z_\Omega/\mathfrak C\). Число пољ \(Z_i\) јест зато највыше равно \(n\).}

Нынє называјемо \(\mathfrak Z\) пољем \emph{првого рода} над \(\mathsf P\), ако число пољ \(Z_i\) јест равно \(n\), а пољем \emph{другого рода}, ако то число јест мење од \(n\).

\emph{\(\mathfrak Z\) јест тогда и само тогда првого рода над \(\mathsf P\), ако \(\mathfrak Z_\Omega\) јест полно редуцибилно.}

Ту видимо прєдност тога основања Галоисовој теорији над обычным. В обычној Галоисовој теорији разгледајут се изоморфизмы једного из пољ \(Z_i\) на друге. Ту се избєгаје та несиметрија. За доказ, же поље не имаје выше коњугованых пољ, неж колико указује јего ступењ, обычно се употрєбје примитивны елемент. Ту тој факт просто слєдује из тога, же вса иредуцибилна прєдставјења од \(\mathfrak Z_\Omega\) сут содрживана в регуларном прєдставјењу.
% END BOOK_S07

% BEGIN BOOK_S08
\subsection*{§ 8. Разпадно поље}\label{der-zerfaellungskoerper}

Поља \(Z_i\) сут иредуцибилна прєдставјења од \(\mathfrak Z\) в \(\Omega\). Ако зато поље \(\Gamma\) меджу \(\mathsf P\) и \(\Omega\), \(\mathsf P\subseteq\Gamma\subseteq\Omega\), содрживаје вса \(Z_i\), то уж \(\mathfrak Z_\Gamma/\mathfrak C\) мора се разпадати на \(t\) пољ првого ступења \(\Gamma\). Обратно, да бы \(\mathfrak Z_\Gamma/\mathfrak C\) разпадало се на \(t\) пољ, кторе сут тогда иредуцибилными прєдставјењами првого ступења од \(\mathfrak Z_\Gamma\), поља \(Z_i\) морајут быти содрживане в \(\Gamma\). Ако разпадным пољем од \(\mathfrak Z\) называмо поље \(\Gamma\) с својством, же \(\mathfrak Z_\Gamma/\mathfrak C\) поставаје прємоју сумоју \(t\) пољ, то имајемо:

\emph{Галоисово поље пољ, изоморфных с \(\mathfrak Z\), \(Z_i\), т.~ј. јих композитум, јест минимално разпадно поље од \(\mathfrak Z\).}
% END BOOK_S08

% BEGIN BOOK_S09
\subsection*{§ 9. Изоморфизмы од \(Z_1\) на \(Z_i\)}\label{die-isomorphismen-von-z1-auf-die-zi}

\(\Theta_i\), ако разгледамо само \(\mathfrak Z\), јест изоморфизм од \(\mathfrak Z\) на \(Z_i\). зато \(S_i=\Theta_i\Theta_1^{-1}\) јест изоморфизм од \(Z_1\) на \(Z_i\). Зато же \(\Theta_i\) сут различне, такоже \(S_i\) сут различне. \(S_i\) сут изоморфизмы једного поља \(Z_1\) на јего коњугована поља, кторе се обычно разгледајут в Галоисовој теорији. Другых не имаје, зато же не имаје другых \(\Theta_i\).
% END BOOK_S09

% BEGIN BOOK_S10
\subsection*{§ 10. Галоисова група}\label{die-galoissche-gruppe}

Ако \(\mathfrak Z\) јест Галоисово поље, то јест ако поља \(Z_i\) сбєгајут се како множства елементов, то \(S_i\) творет групу всих изоморфизмов, кторе оставјајут \(\mathsf P\) неподвижным, од \(Z\) на себе. Зато же всакы продукт \(S_iS_j\) јест автоморфизм од \(Z\), а зато же не имаје другых кромє \(S_i\), \(S_iS_j\) јест равны некому \(S_k\). зато \(S_i\) творет групу всих автоморфизмов, кторе оставјајут всакы елемент из \(\mathsf P\) неподвижным, од \(Z\); то јест \emph{Галоисова група од \(Z\) по одношењу к \(\mathsf P\)}.
% END BOOK_S10

% BEGIN BOOK_S11
\subsection*{§ 11. Главны теорем Галоисовој теорије}\label{der-hauptsatz-der-galoisschen-theorie}

Нынє нехај \(\mathfrak Z\) буде Галоисово разширјење првого рода од \(\mathsf P\). Доказујемо главны теорем Галоисовој теорије:

\emph{Поља \(\Gamma\) меджу \(\mathsf P\) и \(\mathfrak Z\) можно једнозначно взајемно приредити подгрупам \(\mathfrak H\) Галоисовој групы \(\mathfrak G\) од \(Z\) по одношењу к \(\mathsf P\) тако, же пољу \(\Gamma\) приреджена група \(\mathfrak H\) состоји се из всих автоморфизмов од \(Z\), кторе оставјајут всакы елемент из \(\Gamma\) неподвижным, а \(\Gamma\) јест множство всих елементов из \(Z\), кторе оставајут неподвижне при всих автоморфизмах, содрживаных в \(\mathfrak H\).}

Кратко то формулујемо тако:

\begin{enumerate}
\item \emph{Ако \(\mathfrak H\) јест инвариантна група од \(\Gamma\), то \(\Gamma\) јест инвариантно поље од \(\mathfrak H\).}
\item \emph{Ако \(\Gamma\) јест инвариантно поље од \(\mathfrak H\), то \(\mathfrak H\) јест инвариантна група од \(\Gamma\).}
\end{enumerate}

\paragraph{1. можно свести на помочну лему:}\label{laesst-sich-zurueckfuehren-auf-den-hilfssatz}

\emph{Ако \(\mathsf P \subseteq \Sigma \subseteq \mathsf T \subseteq \mathsf Z\) и \(\mathsf T\) имаје ступењ \(t\) над \(\Sigma\), то всакы изоморфизм од \(\Sigma\) на неко коњуговано поље (потрєбно содрживано в \(\mathsf Z\)) можно продолжити до \(t\) различных изоморфизмов од \(\mathsf T\) на коњугована поља.}

Ако прєдположимо, же \(\mathfrak H\) јест инвариантна група од \(\Sigma\), а \(\mathsf T\) јест инвариантно поље од \(\mathfrak H\), можно закључити слєдујуче: Елементи из \(\mathfrak H\) сут продолжења на цєло \(\mathsf Z\) идентичного изоморфизма од \(\Sigma\). Ако те продолжења примєнимо само на \(\mathsf T\), из помочној лемы слєдује, же она разпадајут се на \(t\) клас тако, же изоморфизмы исте класы сут идентичне на \(\mathsf T\), а изоморфизмы различных клас сут различне на \(\mathsf T\). Але зато же \(\mathsf T\) јест инвариантно поље од \(\mathfrak H\), мора быти \(t=1\), \(\mathsf T=\Sigma\), тым јест доказано.

За доказ помочној лемы врнемо се к пољам \(\mathfrak S\), \(\mathfrak T\) меджу \(\mathsf P\) и \(\mathsf Z\), кторым сут \(\Sigma\) и \(\mathsf T\) приреджене посрєдством \(\Theta_1\). Тогда трєба доказати:

\emph{Всакы изоморфизм \(H_i\) од \(\mathfrak S\) на неко \(\Sigma_i\) можно на точно \(t\) различных начинов продолжити до изоморфизмов од \(\mathfrak T\) на подпоља \(\mathsf T_i\) од \(\mathsf Z\).}

Заради прєдположења, же \(\mathfrak Z\) јест првого рода над \(\mathsf P\), \(\mathfrak Z_\Omega/\mathfrak C\) јест идентично с \(\mathfrak Z_\Omega\), зато \(\mathfrak S_\Omega\) разпада се на толико пољ, колико указује његов ступењ \(s\)
\[
\mathfrak S_\Omega = E_1\Omega+\cdots+E_s\Omega .
\]
Модулы прєдставјениј \(E_\nu\Omega\) стварјајут \(s\) изоморфизмов од \(\mathfrak S\) на \(s\) подпољ \(\Sigma_1,\ldots,\Sigma_s\). За најдање изоморфизмов од \(\mathfrak T\) потрєбујемо \(\mathfrak T_\Omega\). Зато же \(E_\nu\Omega\) сут идеалы од \(\mathfrak S_\Omega\), имаје мєсто \(E_\nu\Omega=\mathfrak S_\Omega E_\nu\). Надаље очевидно јест
\[
\mathfrak T_\Omega=\mathfrak T\cdot\Omega=\mathfrak T\cdot\mathfrak S_\Omega .
\]
Зато јест
\[
\mathfrak T_\Omega
=E_1\mathfrak S_\Omega\cdot\mathfrak T+\cdots+E_s\mathfrak S_\Omega\cdot\mathfrak T
=E_1\mathfrak T_\Omega+\cdots+E_s\mathfrak T_\Omega .
\]
То јест разклад од \(\mathfrak T_\Omega\) на идеалы. \(E_\nu\cdot\mathfrak T_\Omega\) разпада се на \(t\) простых идеалов. Зато же сут правдиве релације ступњев
\[
[\mathfrak T_\Omega\cdot E_\nu:\mathfrak S_\Omega E_\nu]
\leqq [\mathfrak T_\Omega:\mathfrak S_\Omega]=[\mathfrak T:\mathfrak S]=t,
\]
зато заради \(\mathfrak S_\Omega E_\nu=\Omega E_\nu\simeq\Omega\):
\[
[\mathfrak T_\Omega\cdot E_\nu:\Omega]
=[\mathfrak T_\Omega E_\nu:\mathfrak S_\Omega E_\nu]\leqq t,
\]
и зато же сума всих ступњев \([\mathfrak T_\Omega E_\nu:\Omega]\) мора быти равна \(s\cdot t\):
\[
[\mathfrak T_\Omega\cdot E_\nu:\Omega]=t,\qquad
E_\nu\mathfrak T_\Omega=e_\nu^{(1)}\cdot\Omega+\cdots+e_\nu^{(t)}\Omega .
\]
Всакы \(e_\nu^{(i)}\Omega\) даваје иредуцибилно прєдставјење од \(\mathfrak T_\Omega\). \emph{Покажимо јешче}, же вса та прєдставјења, ако јих примєнимо само на \(\mathfrak S_\Omega\), сбєгајут се с прєдставјењем, створјеным посрєдством \(E_\nu\cdot\Omega\); тогда јесмо готови. То але јест јасно, зато же из
\[
s\cdot E_\nu=E_\nu\sigma_\nu\quad
(s\in\mathfrak S_\Omega\text{ имаје }\sigma_\nu\text{ како прєдставјење})
\]
слєдује
\[
\sum s\cdot e_\nu^{(i)}=\sum e_\nu^{(i)}\sigma_\nu,
\]
зато заради једнозначности записа сумы
\[
s\cdot e_\nu^{(i)}=e_\nu^{(i)}\sigma_\nu,
\]
нове прєдставјења такоже приреджајут елементу \(s\) елемент \(\sigma_\nu\).

\paragraph{2. Нехај \(\mathsf T\) буде инвариантно поље од \(\mathfrak H\).}
Хочемо познати \(\mathfrak H\) како инвариантну групу од \(\mathsf T\); за то \emph{јест достаточно} показати елемент из \(\mathsf T\), кторы не оставаје неподвижным при неком автоморфизму, не належечем к \(\mathfrak H\).

Елементи \(S_i\) Галоисовој групы \(\mathfrak G\) од \(\mathsf Z\) сут автоморфизмы, кторе оставјајут \(\mathsf P\) неподвижным, од \(\mathsf Z\). Чинимо јих автоморфизмами од \(\mathfrak Z_{\mathsf Z}\) установјењем
\[
S_i z=z,\quad \text{ако } z \text{ јест в } \mathfrak Z.
\]
Примєњење од \(S_i\) на
\[
\mathfrak Z_{\mathsf Z}=e_1\mathsf Z+\cdots+e_n\mathsf Z
\]
мора прємєњивати прєме суманды \(e_\nu\mathsf Z\), зато же оне сут једнозначно опрєдєљене; в особливости мора быти
\[
S_i e_\nu=e_{\nu'},\qquad\text{и}\qquad S_i e_\nu\ne S_k e_\nu
\quad\text{за } i\ne k
\]
Ако \(S_1,\ldots,S_h\) сут елементи из \(\mathfrak H\), изберемо неко \(e_\nu\), рецимо \(e'\), и разгледамо все \(e_\nu\), кторе из њега стварјајут \(S_i\), \(i\leqq h\):
\[
e_1=S_1e',\ldots,e_h=S_he'.
\]
Зато же, ако \(i\leqq h, j\leqq h\),
\[
S_i e_j=S_iS_j e'=S_ke'=e_k \quad \text{при}\quad k\leqq h
\]
всакы \(S_i\) из \(\mathfrak H\) прємєњује \(e_1,\ldots,e_h\) меджу собоју,
\[
E_1=e_1+\cdots+e_h
\]
јест зато при всих \(S_i\in\mathfrak H\) инвариантны елемент из \(\mathfrak Z\). Ако прєдставимо \(E_1\) посрєдством базы од \(\mathfrak Z\), јего коефициенты морајут быти при \(\mathfrak H\) инвариантне елементи из \(\mathsf Z\); зато належет к \(\mathsf T\).

\(E_1\) не јест инвариантны при никаком елементу изван \(\mathfrak H\), \(S_i\). Зато же за тако \(S_i\)
\[
S_i E_1=e_i+\cdots
\]
содрживаје меджу компонентами, кторе не наступајут меджу \(e_1,\ldots,e_h\), компоненту \(e_i\), але запис сумы од \(E_1\) јест једнозначны.
% END BOOK_S11

% BEGIN BOOK_S12
\subsection*{§ 12. Формално значење компонентов \(e_i\)}\label{die-formale-bedeutung-der-komponenten-ei}

За основу изберемо установјену базу \(z_1,\ldots,z_n\) од \(\mathfrak Z\).
\[
\mathfrak Z=z_1\mathsf P+\cdots+z_n\mathsf P
\]
(Нынє о \(\mathfrak Z\) прєдпокладамо само, же оно јест првого рода). Изоморфизм \(\Theta_i\) образује систем \(z_1,\ldots,z_n\) на базу \(\zeta_1^{(i)},\ldots,\zeta_n^{(i)}\) од \(\mathsf Z_i\), и то релацијами
\[
z_\nu e_i=e_i\zeta_\nu^{(i)} .
\]

Обратно, компонента јединице \(e_j\) мора быти изразива посрєдством \(z_\nu\) с коефициентами из \(\mathsf Z\)
\[
e_j=\sum_\nu \varrho_\nu^{(j)}z_\nu .
\]
зато имајемо
\[
e_j=e_j^2=\sum \varrho_\nu^{(j)}z_\nu e_j
=\sum \varrho_\nu^{(j)}\cdot\zeta_\nu^{(j)}\cdot e_j
\]
и
\[
0=e_je_k=\sum \varrho_\nu^{(j)}\cdot z_\nu\cdot e_k
=\sum \varrho_\nu^{(j)}\cdot\zeta_\nu^{(k)}\cdot e_k,
\]
зато:
\[
\sum_\nu \varrho_\nu^{(j)}\cdot\zeta_\nu^{(j)}=1;
\qquad
\sum_\nu \varrho_\nu^{(j)}\cdot\zeta_\nu^{(k)}=0,\quad
\text{ако } j\ne k.
\]
То значи: Матрице
\[
\begin{pmatrix}
\zeta_1^{(1)} & \cdots & \zeta_n^{(1)}\\
\cdots & \cdots & \cdots\\
\cdots & \cdots & \cdots\\
\cdots & \cdots & \cdots\\
\zeta_1^{(n)} & \cdots & \zeta_n^{(n)}
\end{pmatrix}
\quad\text{и}\quad
\begin{pmatrix}
\varrho_1^{(1)} & \cdots & \varrho_1^{(n)}\\
\cdots & \cdots & \cdots\\
\cdots & \cdots & \cdots\\
\cdots & \cdots & \cdots\\
\varrho_n^{(1)} & \cdots & \varrho_n^{(n)}
\end{pmatrix}
\]
сут взајемно обратне. То можно такоже изразити слєдујуче: \(\varrho_1^{(i)},\ldots,\varrho_n^{(i)}\) јест база комплементарна к \(\zeta_1^{(i)},\ldots,\zeta_n^{(i)}\). Она имаје ролу в теорији диференты числового поља.
% END BOOK_S12

% BEGIN BOOK_S13
\subsection*{§ 13. Обчи теорем о продолжењу}\label{der-allgemeine-fortsetzungssatz}

Теорем о продолжењу из \(H_i\), доказан на стр.~10, можно обчити.

Нехај \(\mathfrak Z=z_1\mathsf P+\cdots+z_n\mathsf P\) буде комутативны, хиперкомплексны, полно редуцибилны систем. Нехај \(\mathfrak S\) и \(\mathfrak T\) будут -- в тој ситуацији потрєбно такоже полно редуцибилне -- подсистемы од \(\mathfrak Z\):
\[
\mathsf P\subseteq\mathfrak S\subseteq\mathfrak T\subseteq\mathfrak Z .
\]
По M. З. § 13 всакы из трох системов имаје јединицу. Але те три јединице никако не морајут быти једнаке. (Сигурно не сут једнаке, ако \(\mathfrak S\) и \(\mathfrak T\) сут идеалы од \(\mathfrak Z\), различне од \(\mathfrak Z\).) Нехај \(E\) буде јединица од \(\mathfrak S\). Систем \(\mathfrak T\cdot E\) јест модул конечного ранга над \(\mathfrak S\):
\[
E\mathfrak T=a_1\mathfrak S+\cdots+a_t\mathfrak S
\]
(\(\mathfrak T\) обче не јест!).

\emph{Доказ.} \(\mathfrak S\) јест прєма сума пољ: \(\mathfrak S=\mathfrak K_1+\cdots+\mathfrak K_r\). Тому одповєдаје разклад од \(\mathfrak T E\) на прєму суму. Всакы суманд од \(\mathfrak T E\) мора содрживати једно из пољ \(\mathfrak K_i\), иначе бы \(E\) ануловало го, але \(E\) јест јединица од \(\mathfrak T E\). Зато \(\mathfrak T E\) јест прєма сума разширјујучих прстенов од \(\mathfrak K_i\), т.~ј. прєма сума формы
\[
\mathfrak T E=a_{11}\mathfrak K_1+\cdots+a_{1p_1}\mathfrak K_1+\cdots+a_{rp_r}\mathfrak K_r,
\]
или \(E\mathfrak T=a_{11}\mathfrak S+\cdots+a_{rp_r}\mathfrak S\), тым јест доказано.

\begin{center}
Обчи теорем о продолжењу.
\end{center}

\emph{Всако иредуцибилно прєдставјење од \(\mathfrak S\) в \(\Omega\), задано хомоморфизмом \(H_i\), можно на точно \(t\) различных начинов продолжити до прєдставјениј од \(\mathfrak T\).}

Доказ јест аналогны доказу на стр.~10--11. Имајемо
\[
\mathfrak S_\Omega=E_1\Omega+\cdots+E_s\Omega,\qquad
E_\nu\Omega=\mathfrak S_\Omega E_\nu,
\]
\[
\mathfrak T_\Omega\cdot E=\mathfrak T\cdot E\Omega
=\mathfrak T\mathfrak S\Omega=\mathfrak T\mathfrak S_\Omega,
\]
\[
\mathfrak T_\Omega\cdot E
=\mathfrak T(E_1\mathfrak S_\Omega+\cdots+E_s\mathfrak S_\Omega)
=E_1\mathfrak T_\Omega+\cdots+E_s\mathfrak T_\Omega,
\]
\[
[\mathfrak T_\Omega E_\nu:\Omega]
=[\mathfrak T_\Omega E_\nu:\Omega E_\nu]
=[\mathfrak T_\Omega E_\nu:E_\nu\mathfrak S_\Omega]
\leqq[\mathfrak T_\Omega\cdot E:\mathfrak S_\Omega]
=[\mathfrak T\cdot E:\mathfrak S]=t,
\]
\[
\sum[\mathfrak T_\Omega\cdot E_\nu:\Omega]=st
\]
зато
\[
[\mathfrak T_\Omega E_\nu:\Omega]=t.
\]
\(\mathfrak T_\Omega E_\nu\) разпада се зато на \(t\) идеалов, из кторых всакы како суманд од \(\mathfrak T_\Omega E\), зато такоже од \(\mathfrak T_\Omega\), даваје прєдставјење од \(\mathfrak T_\Omega\), кторо, како на стр.~10--11, оказује се продолжењем прєдставјења заданого посрєдством \(H_\nu\). Остала иредуцибилна прєдставјења од \(\mathfrak T_\Omega\) -- створјена сумандами, не содрживаными в \(\mathfrak T_\Omega\cdot E_\nu\), од \(\mathfrak T_\Omega\) -- не могут быти продолжења прєдставјениј, створјеных посрєдством \(H_\nu\), од \(\mathfrak S_\Omega\), зато же те идеалы сут ануловане посрєдством \(\mathfrak T_\Omega\cdot E_\nu\), а тым выше посрєдством \(\mathfrak S_\Omega\).

\section*{Глава III. Абелове групы}\label{kapitel-iii.-abelsche-gruppen}
% END BOOK_S13

% BEGIN BOOK_S14
\subsection*{§ 14. Групово прстен}\label{der-gruppenring}

Нехај \(\mathfrak G\) буде конечна група с елементами \(a_1,\ldots,a_h\), а \(\mathsf P\) комутативно поље. С груповым множењем како множењем образујемо хиперкомплексны систем
\[
\mathfrak o[\mathfrak G]=a_1\mathsf P+\cdots+a_h\mathsf P
\]
\emph{групово прстен од \(\mathfrak G\) в \(\mathsf P\)}. (M. З. § 6, Спеисер, Theorie д. Gruppen енд. Орд. § 48).

Пытајемо за прєдставјења од \(\mathfrak o[\mathfrak G]\) в \(\mathsf P\) или в алгебраично затворјеном пољу \(\Omega\) над \(\mathsf P\). Како уж было споменуто, иредуцибилна прєдставјења сут створјена простыми лєвыми идеалами од \(\mathfrak o[\mathfrak G]/\mathfrak C\). Зато нас -- како в Галоисовој теорији -- интересује пытање, когда \(\mathfrak o[\mathfrak G]\) имаје радикал. О том имаје мєсто обчи теорем:

\begin{center}
\emph{\(\mathfrak o[\mathfrak G]\) јест тогда и само тогда полно редуцибилно, ако ред \(h\) од \(\mathfrak G\) не јест дєливы карактеристикоју \(p\) од \(\mathsf P\).}
\end{center}

Најпрво доказујемо:
\[
\text{Из } h\equiv0(p) \text{ слєдује } \mathfrak C\ne 0.
\]
\[
\mathfrak a=\mathsf P(a_1+\cdots+a_h)
\]
јест двустранны идеал од \(\mathfrak o\). Зато же јест
\[
\mathsf P\cdot(a_1+\cdots+a_h)a
=\mathsf P\cdot(a_1a+\cdots+a_ha)
=\mathsf P(a_1+\cdots+a_h)=\mathfrak a
\]
и
\[
a\cdot\mathsf P(a_1+\cdots+a_h)
=\mathsf P(aa_1+\cdots+aa_h)
=\mathsf P(a_1+\cdots+a_h)=\mathfrak a
\]

заради групового својства елементов \(a_i\). \(\mathfrak a\) очевидно не јест нула, але јест нилпотентны:
\[
\left(\sum_i a_i\right)^2=\sum_i\sum_j a_i a_j
=h\sum_j a_j=0.
\]

Обратно тврджење:
\[
\text{Из } h\not\equiv0(p) \text{ слєдује } \mathfrak C=0,
\]
нынє доказујемо само за Абелове групы. (В M. З. § 26 та чест теорема јест доказана обче.)
% END BOOK_S14

% BEGIN BOOK_S15
\subsection*{§ 15. Групове прстены Абеловых груп}\label{die-gruppenringe-abelscher-gruppen}

Под \(\mathfrak A=\{a_1,\ldots,a_h\}\) разумєјемо Абелову групу. Тогда \(\mathfrak o\) јест комутативно, зато можно примєнити обче теоремы о комутативных хиперкомплексных системах.

\emph{Ако \(h\) не јест дєливы карактеристикоју \(p\) од \(\mathsf P\), то \(\mathfrak o[\mathfrak A]\) имаје точно \(h\) различных иредуцибилных прєдставјениј в \(\Omega\). (Она сут првого ступења, зато же \(\mathfrak o\) јест комутативно.)}

Из тога слєдује, же \(\mathfrak o[\mathfrak A]/\mathfrak C\) имаје тојже ранг како \(\mathfrak o\), а тым самым \(\mathfrak C=0\); то јест втора чест попрєдного, оставјеного без доказа теорема за Абелове групы.

\emph{Доказ.} \(\mathfrak A\) јест прємы продукт цикличных груп \(\mathfrak Z_i\) редов \(h_i\)
\[
\mathfrak A=\mathfrak Z_1\times\cdots\times\mathfrak Z_t,\qquad
h=h_1h_2\cdots h_t .
\]
Ако \(z_i\) јест творечи елемент од \(\mathfrak Z_i\), то јему одговарјајучи елемент иредуцибилного прєдставјења (првого ступења) мора быти \(h_i\)-ты корењ из јединице \(\varepsilon_\nu^{(h_i)}\). С другој страны, ако всакому \(z_i\) приредимо \(h_i\)-ты корењ из јединице \(\varepsilon_\nu^{(h_i)}\), добыва се иредуцибилно прєдставјење од \(\mathfrak o\) в \(\Omega\). Зато же нынє \(h\not\equiv0(p)\), а тым выше \(h_i\not\equiv0(p)\), обстајајут точно \(h_1h_2\cdots h_t=h\) различне прєдставјења од \(\mathfrak o\), тым јест доказано.

(Такоже јест легко видєти: ако \(h=0(p)\), то \(h_i=0(p)\) за принајмење једно \(h_i\), обстајајут највыше \(h_i/p\) различне \(\varepsilon_\nu^{(h_i)}\), зато мење неж \(h\) прєдставјениј, зато \(\mathfrak o_\Omega\) мора содрживати ненуловы радикал. Але из тога безпосрєдње не слєдује ничто о радикалу од \(\mathfrak o\).)

Те \(h\) различне хомоморфизмы од \(\mathfrak o_\Omega\) на подпрстены од \(\Omega\) означимо с \(\Theta_1,\ldots,\Theta_h\); всакому \(\Theta_i\) одповєдаје суманд \(e_i\Omega\) од \(\mathfrak o_\Omega\) како модул прєдставјења. Нєкды примєњујемо \(\Theta_i\) само на \(\mathfrak A\) и тогда говоримо о прєдставјењу од \(\mathfrak A\). Те хомоморфизмы од \(\mathfrak A\) на групы корењев из јединице сут \emph{\(h\) карактеров од \(\mathfrak A\)}. Специалны карактер јест \emph{главны карактер}, кторы всакому елементу из \(\mathfrak A\) приреджаје \(1\). Нехај приналежечи \(\Theta_i\) буде \(\Theta_1\).

Прєдставјење главным карактером можно за всаку групу, не само за Абелове групы. Зато же то прєдставјење уж лежи в \(\mathsf P\), в случају, когда \(\mathfrak o\) јест полно редуцибилно, \(\mathfrak o\) мора одцєпити просты суманд \(E\cdot\mathfrak o\), кторы даваје то прєдставјење. \(E\) можно безпосрєдње израчунати. Мора быти \(E=\sum x_i a_i\). Прво јест
\[
aE=E\cdot 1\quad\text{или}\quad \sum x_i a a_i=\sum x_i a_i
\]
из чего за \(a=a_j a_i^{-1}\)
\[
x_1a_ja_i^{-1}a_1+\cdots+x_i a_j+\cdots+x_ha_ja_i^{-1}a_h
=x_1a_1+\cdots+x_ha_h
\]
зато
\[
x_i=x_j
\]
слєдује.

Друго јест:

\[
E=E^2=x^2\left(\sum a_i\right)^2=x^2h\sum a_i=xhE.
\]

В случају \(h=0(p)\) слєдује \(E=0\), т.~ј.: \(\mathfrak o\) не содрживаје модул прєдставјења за главны карактер, зато не могло јест быти полно редуцибилно. Тым јест поновно доказана једна половица теорема на стр.~13.

Ако але \(h\not\equiv0(p)\), то \(\mathfrak o\) јест полно редуцибилно, непосрєдње слєдује

\[
x=\frac{1}{h},\qquad E=\frac{1}{h}(a_1+\cdots+a_h).
\]
% END BOOK_S15

% BEGIN BOOK_S16
\subsection*{§ 16. Релације карактеров}\label{die-charakterenrelationen}

Факт, же \(E=(1/h)\cdot\sum a_i\) даваје главны карактер, непосрєдње води к релацијам карактеров. Мора быти

\[a \cdot e_i = e_i \cdot \Theta_i a\]

зато

\[e_i \cdot \Theta_i e_j = \begin{cases} e_i & i = j \\ 0 & i \neq j \end{cases}\]

или

\[\Theta_i e_j = \begin{cases} 1 & i = j \\ 0 & i \neq j \end{cases}\]

В особливости јест

\[\Theta_i e_1 = \begin{cases} 1 & i = 1 \\ 0 & i \neq 1 \end{cases}\]

или, зато же \(e_1 = (1/h) \sum a_i\),

\[
\sum_\nu \Theta_i a_\nu=\begin{cases} h & i=1,\\ 0 & i\ne 1, \end{cases}
\]

То сут познате релације карактеров:

\emph{Сума једного карактера по всих елементах групы јест нула, кромє главного карактера, за кторы она јест х.}
% END BOOK_S16

% BEGIN BOOK_S17
\subsection*{§ 17. Галоисова теорија Абеловых груп}\label{die-galoissche-theorie-abelscher-gruppen}

За групове прстены Абеловых груп можно установити теоремы подобне оным, кторе сут доказиване в Галоисовој теорији комутативных пољ. В Галоисовој теорији \(\Theta_i\Theta_1^{-1}\) творили сут групу \(\mathfrak{G}\). Главны теорем изражаје једнозначно взајемно приреджење, основано на установјеных инвариантных својствах, меджу подгрупами од \(\mathfrak{G}\) и подпољами од \(Z\).

Хомоморфизмы \(\Theta_i\) од \(\mathfrak{o}[\mathfrak{A}]\) можно сјединити в групу \(\mathsf A\), чије подгрупы стојет в подобном одношењу к подгрупам од \(\mathfrak{A}\).

\emph{Ако продукт \(\Theta_i \Theta_k = \Theta\) дефинујемо формулом}

\[\Theta a = \Theta_i a \cdot \Theta_k a,\]

\emph{то множство всих \(\Theta_i\) јест с \(\mathfrak A\) изоморфна група \(\mathsf A\).}

\emph{Доказ.} \(\Theta = \Theta_i \cdot \Theta_k\) јест заради

\[\begin{aligned}
\Theta a \cdot \Theta b &= \Theta_i a \cdot \Theta_k a \cdot \Theta_i b \cdot \Theta_k b \\
&= \Theta_i a \cdot \Theta_i b \cdot \Theta_k a \cdot \Theta_k b \\
&= \Theta_i a b \cdot \Theta_k a b = \Theta a b
\end{aligned}\]

хомоморфизм од \(\mathfrak A\) на \(\mathsf A\), а зато же кромє \(\Theta_i\) не имаје другых хомоморфизмов од \(\mathfrak A\), \(\Theta\) јест равны некому \(\Theta_j\); зато \(\Theta_i\) наистино творет групу из \(h\) елементов.

Нынє под \(\varepsilon_{h_\nu}\) разумєјемо примитивны \(h_\nu\)-ты корењ из јединице. Тогда \(\Theta_i\) образује творечи елемент \(z_\nu\) цикличного фактора \(\mathfrak Z_\nu\) од \(\mathfrak A\) на потенцију од \(\varepsilon_{h_\nu}\)

\[
\Theta_i z_\nu=\varepsilon_{h_\nu}^{\varrho_\nu^{(i)}}.
\]

Хомоморфизму \(\Theta_i\) приреджајемо елемент \(\prod_\nu z_\nu^{\varrho_\nu^{(i)}}\) од \(\mathfrak A\). То јест једнозначно взајемно приреджење од \(\mathsf A\) на \(\mathfrak A\), зато же \(\Theta_i\) јест опрєдєљен посрєдством \(\varrho_\nu^{(i)}\), различным \(\Theta_i\) одповєдајут различне \(\prod_\nu z_\nu^{\varrho_\nu^{(i)}}\). Приреджење јест изоморфно. Зато же с једној страны

\[
\Theta_i\Theta_j z_\nu=\Theta_i z_\nu\cdot \Theta_j z_\nu
=\varepsilon_{h_\nu}^{\varrho_\nu^{(i)}+\varrho_\nu^{(j)}},
\]

а с другој страны

\[
\prod_\nu z_\nu^{\varrho_\nu^{(i)}}\cdot
\prod_\nu z_\nu^{\varrho_\nu^{(j)}}=
\prod_\nu z_\nu^{\varrho_\nu^{(i)}+\varrho_\nu^{(j)}},
\quad\text{тым јест доказано.}
\]

Нынє дефинујемо

\emph{Инвариантну област подгрупы \(\mathsf B\) од \(\mathsf A\):}

Подгрупа \(\mathfrak B\) елементов од \(\mathfrak A\), кторе всакы \(\Theta\in\mathsf B\) образује на 1.

\emph{Инвариантну групу подгрупы \(\mathfrak B\) од \(\mathfrak A\):}

Подгрупа \(\mathsf B\) елементов од \(\mathsf A\), кторе образујут всакы \(a\in\mathfrak B\) на 1.

и формулујемо главны теорем:

\emph{Подгрупы \(\mathfrak B\) од \(\mathfrak A\) и подгрупы \(\mathsf B\) од \(\mathsf A\) можно једнозначно взајемно приредити тако, же ако \(\mathfrak B\) и \(\mathsf B\) одповєдајут једна другој, \(\mathsf B\) јест инвариантна група од \(\mathfrak B\), а \(\mathfrak B\) јест инвариантна област од \(\mathsf B\).}

Трєба јест зато показати:

\begin{enumerate}
\item \emph{Ако \(\mathsf B\) јест инвариантна група од \(\mathfrak B\), то \(\mathfrak B\) јест инвариантна област од \(\mathsf B\).}
\item \emph{Ако \(\mathfrak B\) јест инвариантна област од \(\mathsf B\), то \(\mathsf B\) јест инвариантна група од \(\mathfrak B\).}
\end{enumerate}

\paragraph{1.}
слєдује како на стр.~10 из обчего теорема о продолжењу. Прєдпокладамо, же \(\mathsf B\) јест инвариантна група од \(\mathfrak B\), \(\mathfrak B'\) инвариантна област од \(\mathsf B\), а \(\mathfrak B'\) содрживаје \(\mathfrak B\). Елементи од \(\mathsf B\) сут продолжења на цєлу \(\mathfrak A\) једного хомоморфизма од \(\mathfrak B\), кторы образује \(\mathfrak B\) на 1. Ако те продолжења примєнимо само на \(\mathfrak B'\), из обчего теорема о продолжењу слєдује, же она разпадајут се на \(t\) клас тако, же хомоморфизмы исте класы сут идентичне, а хомоморфизмы различных клас сут различне на \(\mathfrak B'\); при том \(t\) јест ступењ од \(\mathfrak o[\mathfrak B']\) над \(\mathfrak o[\mathfrak B]\), то јест индекс од \(\mathfrak B\) в \(\mathfrak B'\), зато же \(\mathfrak o[\mathfrak B]\) и \(\mathfrak o[\mathfrak B']\) имајут тојже елемент јединице, јменно елемент јединице од цєлој \(\mathfrak A\). Але зато же \(\mathfrak B'\) јест инвариантна област од \(\mathsf B\), мора быти \(t=1\), зато \(\mathfrak B'=\mathfrak B\), тым јест доказано.

\textbf{2.} сводимо на 1.

Елементы \(a\) од \(\mathfrak A\) разгледамо хомоморфизмами од \(\mathsf A\). Установјујемо, же \(a\Theta\) јест равно корењу из јединице \(\Theta a\), кратко \(a\Theta=\Theta a\). Тым \(a\) јест наистино опрєдєљен како хомоморфизм од \(\mathsf A\), зато же
\[
a\Theta\cdot a\Theta'=\Theta a\cdot\Theta' a=\Theta\Theta' a=a\Theta\Theta'.
\]

Продукт двох хомоморфизмов \(a,b\) јест обычны груповы продукт:
\[
\begin{array}{cccccc}
ab\Theta&=&a\Theta\cdot b\Theta&=&\Theta a\cdot\Theta b=\Theta ab&=ab\Theta\\
&&\text{продукт хомоморфизмов}&&\text{груповы продукт}&
\end{array}
\]

Елементи \(a,b\), различне како групове елементи, такоже сут различне како хомоморфизмы: из \(a\Theta=b\Theta\) за все \(\Theta\) слєдује \(\Theta a=\Theta b\) за все \(\Theta\), или \(\Theta ab^{-1}=1\) за все \(\Theta\). зато \(\Theta\) сут продолжења на цєлу \(\mathfrak A\) идентичного хомоморфизма посрєдством \(ab^{-1}\) створјене цикличној подгрупы \(\mathfrak Z\) од \(\mathfrak A\). По теорему о продолжењу индекс од \(\mathfrak Z\) в \(\mathfrak A\) јест равны \(h\), т.~ј. \(\mathfrak Z=e,\ a=b\).

Нынє прєводимо тврджења, учињена в новом разумєњу,
\[
\begin{array}{ll}
\text{а)}& \mathfrak B\ \text{јест инвариантна група од }\mathsf B.\\
\text{б)}& \mathsf B\ \text{јест инвариантна област од }\mathfrak B.
\end{array}
\]
в старо разумєње:

К а). \(\mathfrak B\) состоји се из всих \(a\in\mathfrak A\), за кторе \(a\Theta=1\) за все \(\Theta\in\mathsf B\), т.~ј. \(\mathfrak B\) состоји се из всих \(a\in\mathfrak A\), за кторе \(\Theta\cdot a=1\) за все \(\Theta\in\mathsf B\). зато прєвод од а) јест
\[
\text{а')} \quad \mathfrak B\ \text{јест инвариантна област од }\mathsf B.
\]

Одговарјајучи прєвод од б) јест
\[
\text{б')} \quad \mathsf B\ \text{јест инвариантна група од }\mathfrak B.
\]

зато прєвод тврджења уж доказаного под 1.:

Ако \(\mathfrak B\) јест инвариантна група од \(\mathsf B\), то \(\mathsf B\) јест инвариантна област од \(\mathfrak B\),
јест равно тврджењу, кторо трєба доказати:

Ако \(\mathfrak B\) јест инвариантна област од \(\mathsf B\), то \(\mathsf B\) јест инвариантна група од \(\mathfrak B\).

Ако релације карактеров из стр.~15 запишемо за \(\mathsf A\) вмєсто за \(\mathfrak A\), добывајут се једначења
\[
\sum_{i=1}^{h} \Theta_i a_\nu =
\begin{cases}
h & \nu=1,\\
0 & \nu\ne 1.
\end{cases}
\]

Сума всих карактеров над груповым елементом јест нула, кромє елемента јединице, за кторы она јест \(h\).

\section*{Глава IV. Двустранно просте прстены}\label{kapitel-iv.-zweiseitig-einfache-ringe}
% END BOOK_S17

% BEGIN BOOK_S18
\subsection*{§ 18. Помочна лема}\label{ein-hilfssatz}

Под \(\mathsf K\) разумєјемо тєло, а под \(\mathfrak G\) групу автоморфизмов (изоморфных образованиј на себе) од \(\mathsf K\). Множство елементов од \(\mathsf K\), кторе оставајут неизмєњене при всаком автоморфизму належечем к \(\mathfrak G\), јест подтєло \(\mathsf P\) од \(\mathsf K\), \emph{инвариантно тєло од} \(\mathfrak G\).

Ако \(\tau\ne 0\) јест либовољны елемент од \(\mathsf K\), автоморфизм од \(\mathsf K\) добыва се тако, же елементу од \(\mathsf K\), \(\alpha\), приредимо елемент \(\tau^{-1}\alpha\tau\); по одговарјајучем назвању из теорије груп тако автоморфизм называмо \emph{внутрным автоморфизмом}. Инвариантно тєло \(\mathsf P\) групы всих внутрных автоморфизмов тєла \(\mathsf K\) јест његов центар.

Нынє нехај будут задане тєло \(\mathsf K\), група автоморфизмов \(\mathfrak G\) од \(\mathsf K\) и њено инвариантно тєло \(\mathsf P\).

\[
\mathfrak M=x_1\mathsf P+\cdots+x_n\mathsf P
\]

нехај буде \(\mathsf P\)-модул ранга \(n\). Дєјство елемента из \(\mathfrak G\), \(\gamma\), на елементы од \(\mathfrak M\) дефинујемо формулом

\[
\gamma\cdot x_i=x_i,\qquad\text{зато}\qquad
\gamma\sum \varrho_i x_i=\sum \varrho_i x_i.
\]

Послєдица тога јест, же \(\mathfrak G\) јест постала групоју автоморфизмов разширјеного модула

\[
\mathfrak M_{\mathsf K}=x_1\mathsf K+\cdots+x_n\mathsf K
\]

а инвариантна област од \(\mathfrak G\) в \(\mathfrak M_{\mathsf K}\) јест тогда \(\mathfrak M\). При тых прєдположењах имаје мєсто теорем:

\begingroup\itshape
Ако
\[
\mathfrak C=z_1\mathsf K+\cdots+z_r\mathsf K
\]

јест подмодул од \(\mathfrak M_{\mathsf K}\), кторы всакы елемент из \(\mathfrak G\), \(\gamma\), образује в њега самого (не потрєбно по елементах), то \(\mathfrak C\) јест разширјены модул, образованы с \(\mathsf K\), некого подмодула \(\mathfrak c\) од \(\mathfrak M\).

\[
\mathfrak C=\mathfrak c_{\mathsf K}.
\]
\endgroup

\emph{Доказ.} 1. Попрєдна припомнєнка: Мора быти \(\mathfrak c=\mathfrak C\cap\mathfrak M\). Зато же обче имаје мєсто
\[
\mathfrak c=\mathfrak c_{\mathsf K}\cap\mathfrak M.
\]

\(\mathfrak c\subseteq\mathfrak c_{\mathsf K}\cap\mathfrak M\) јест јасно. За доказ од \(\mathfrak c_{\mathsf K}\cap\mathfrak M\subseteq\mathfrak c\) изберемо базу \(y_1,\ldots,y_r\) од \(\mathfrak c\),

\[
\mathfrak c=y_1\mathsf P+\cdots+y_r\mathsf P,
\]

тогда јест

\[
\mathfrak c_{\mathsf K}=y_1\mathsf K+\cdots+y_r\mathsf K.
\]

Всакы елемент из \(\mathfrak c_{\mathsf K}\cap\mathfrak M\), \(c\), допушча базовы запис како елемент од \(\mathfrak c_{\mathsf K}\)
\[
c=\sum_{i=1}^{r}y_i\chi_i
\]
и базовы запис како елемент од \(\mathfrak M\)

\[
c=\sum_{i=1}^{r}y_i\varrho_i+\sum_{j=r+1}^{n}y_j\varrho_j
\]

ако под \(y_{r+1},\ldots,y_n\) разумєјемо систем елементов од \(\mathfrak M\), кторы дополњаје \(y_1,\ldots,y_r\) до базы од \(\mathfrak M\) (M. З. § 6).

Зато же оба записа можно разгледати базовыми записами в \(\mathfrak M_{\mathsf K}\), они сут идентичне, \(\chi_i=\varrho_i\ (i=1,\ldots,r)\), \(\varrho_j=0\ (j=r+1,\ldots,n)\), зато

\[
c=\sum_{i=1}^{r}y_i\varrho_i
\qquad\text{тым јест доказано.}
\]

\textbf{2.} Да бы доказати разгледаны теорем, конструујемо базу \(z_1,\ldots,z_r\) од \(\mathfrak C\), чији елементи \(z_i\) належет к \(\mathfrak M\). Тогда, ако положимо

\[
\mathfrak c=z_1\mathsf P+\cdots+z_r\mathsf P
\]

наистино јест

\[
\mathfrak C=\mathfrak c_{\mathsf K}.
\]

Изберемо либовољну базу \(y_1,\ldots,y_r\) од \(\mathfrak C\) и дополнимо ју до базы од \(\mathfrak M_{\mathsf K}\) одговарјајучими елементами \(x_{r+1},\ldots,x_n\) једној базы \(x_1,\ldots,x_n\) од \(\mathfrak M\), ктора јест такоже база од \(\mathfrak M_{\mathsf K}\). То можно по M. З. § 5.

\[
\mathfrak M_{\mathsf K}=\mathfrak C+x_{r+1}\mathsf K+\cdots+x_n\mathsf K.
\]

Из тога за првых \(r\) елементов \(x_1,\ldots,x_r\) базы од \(\mathfrak M\) слєдујут једначења

\[
x_i=z_i-\sum_{j=r+1}^{n}x_j\chi_j^{(i)},\qquad
z_i\in\mathfrak C,\quad \chi_j^{(i)}\in\mathsf K.
\]

Зато же \(z_i\) заједно с \(x_{r+1},\ldots,x_n\) очевидно творет базу од \(\mathfrak M_{\mathsf K}\), они сут линеарно независне; зато \(z_1,\ldots,z_r\) јест база од \(\mathfrak C\). \(z_1,\ldots,z_r\) јест база исканој формы, чији \(z_i\) належет к \(\mathfrak M\):

Нехај \(\gamma\) буде елемент из \(\mathfrak G\). Тогда с једној страны

\[
\begin{array}{rcl}
(1)\qquad \gamma z_i
&=&\gamma x_i+\displaystyle\sum_{j=r+1}^{n}\gamma x_j\cdot \gamma\chi_j^{(i)}
=x_i+\displaystyle\sum_{j=r+1}^{n}x_j\gamma\chi_j^{(i)}.
\end{array}
\]

а с другој страны
\[
\begin{array}{rcl}
(2)\qquad \gamma z_i&=&\displaystyle\sum_{\mu=1}^{r}z_\mu\sigma_\mu^{(i)},
\qquad \sigma_\mu^{(i)}\in\mathsf K,
\end{array}
\]

зато же \(\mathfrak C\) допушча все \(\gamma\in\mathfrak G\). (2) можно записати јешче в формє

\[
\begin{array}{rcl}
(3)\qquad \gamma z_i
&=&\displaystyle\sum_{\mu=1}^{r}x_\mu\sigma_\mu^{(i)}
+\displaystyle\sum_{j=r+1}^{n}x_j\sum_{\mu=1}^{r}\chi_j^{(\mu)}\sigma_\mu^{(i)}
\end{array}
\]

Сравњење коефициентов в (1) и (3) даваје

\[
\sigma_j^{(i)}=
\begin{cases}
1 & i=j,\\
0 & i\ne j.
\end{cases}
\]

Зато јест

\[
\gamma\cdot z_i=z_i.
\]

зато \(z_i\) належет к инвариантној области од \(\mathfrak G\), т.~ј. к \(\mathfrak M\), тым јест доказано.
% END BOOK_S18

% BEGIN BOOK_S19
\subsection*{§ 19. Прєдставјења двустранно простых хиперкомплексных системов в некомутативных тєлах, кторе разширјајут јих поља коефициентов}\label{darstellungen-zweiseitig-einfacher-hyperkomplexer-systeme-in-nichtkommutativen-erweiterungskoerpern-ihrer-koeffizientenkoerper}

\[
\mathfrak S=x_1\mathsf P+\cdots+x_n\mathsf P
\]

нехај буде двустранно просты, зато полно редуцибилны хиперкомплексны систем, а \(\mathsf K\) тєло с центром \(\mathsf P\).

\textbf{Теорем 1.} \emph{\(\mathfrak S_{\mathsf K}\) јест двустранно просто (зато такоже полно редуцибилно).}

\emph{Доказ} 1. Всакы двустранны \(\mathfrak S_{\mathsf K}\)-идеал \(\mathfrak a\) јест инвариантны по одношењу к групє \(\mathfrak G\) внутрных автоморфизмов од \(\mathsf K\). Зато же ако, под \(\chi\) разумєјучи ненуловы елемент од \(\mathsf K\), положимо \(\chi^{-1}\mathfrak a\chi=\bar{\mathfrak a}\), то \(\bar{\mathfrak a}\subseteq\mathfrak a\) заради својства идеала од \(\mathfrak a\); але зато же \(\bar{\mathfrak a}\), разгледаны како \(\mathsf K\)-модул, имаје над \(\mathsf K\) тојже ранг како \(\mathfrak a\), мора быти \(\bar{\mathfrak a}=\mathfrak a\).

\textbf{2.} По управо установјеној помочној леми двустранны \(\mathfrak S_{\mathsf K}\)-идеал \(\mathfrak a\) мора зато быти разширјены идеал некого подмодула \(\mathfrak A\) од \(\mathfrak S\), \(\mathfrak a=\mathfrak A_{\mathsf K}\). \(\mathfrak A\) јест двустранны идеал од \(\mathfrak S\), зато же из \(\mathfrak A=\mathfrak S\cap\mathfrak a\) добыва се
\[
\mathfrak S\mathfrak A\subseteq\mathfrak S\mathfrak S\cap\mathfrak S\mathfrak a\subseteq\mathfrak S\cap\mathfrak a=\mathfrak A
\]
\[
\mathfrak A\mathfrak S\subseteq\mathfrak S\mathfrak S\cap\mathfrak a\mathfrak S\subseteq\mathfrak S\cap\mathfrak a=\mathfrak A.
\]

Але зато же \(\mathfrak S\) содрживаје само двустранне идеалы 0 и \(\mathfrak S\), в \(\mathfrak S_{\mathsf K}\) такоже обстајајут само два двустранна идеала 0 и \(\mathfrak S_{\mathsf K}\), тым јест доказано.

Нынє разгледајемо иредуцибилна прєдставјења од \(\mathfrak S\) в \(\mathsf K\), и хочемо установити реципрочна прєдставјења; то имаје формалну прєдност. В конечном јест всеједно, избирајемо ли прєма или реципрочна прєдставјења, зато же она сут једнозначно взајемно приреджена.

Нехај \(\mathfrak M\) буде реципрочны модул прєдставјења од \(\mathfrak S\) в \(\mathsf K\), то јест лєвы \(\mathfrak S\)-модул и лєвы \(\mathsf K\)-модул конечного ранга,
\[
\mathfrak M=\mathsf K\cdot y_1+\cdots+\mathsf K\cdot y_r
\]
с законом комутативности
\[
s(xm)=x(sm).
\]

Доказујемо, же \(\mathfrak M\) можно такоже разгледати како лєвы \(\mathfrak S_{\mathsf K}\)-модул. Најпрво трєба објаснити множење елемента од \(\mathfrak S_{\mathsf K}\) елементом од \(\mathfrak M\). Ако \(s\in\mathfrak S,\ x\in\mathsf K,\ m\in\mathfrak M\), установјујемо
\[
sx\cdot m=s\cdot(x\cdot m)
\]
и аналогично за обчи \(\mathfrak S_{\mathsf K}\)-елемент \(\sum s_ix_i\)
\[
\sum s_ix_i\cdot m=\sum s_i\cdot(x_i\cdot m).
\]

Можно увєрити се, же та дефиниција јест једнозначна и не зависи од начина, на кторы елемент од \(\mathfrak S_{\mathsf K}\) записујемо како суму продуктов \(s_ix_i\).

Трєба зато из
\[
\begin{array}{rcl}
(1)\qquad \displaystyle\sum_{i=1}^{p}s_ix_i&=&0
\end{array}
\]
закључити
\[
\sum_{i=1}^{p}s_i(x_i m)=0
\]
Прєдпокладамо, же \(s_1,\ldots,s_q\) сут линеарно независне и же \(s_{q+1},\ldots,s_p\) можно изразити посрєдством \(s_1,\ldots,s_q\)
\[
s_j=\sum_{i=1}^{q}s_i\varrho_i^{(j)},\qquad j=q+1,\ldots,p.
\]

Коефициенты морајут лежати в \(\mathsf P\), зато же \(s_j\) можно изразити само на једен начин, а запис јест уж можны в \(\mathfrak S\). Из (1) добыва се

\[
x_i+\sum_{j=q+1}^{p}\varrho_i^{(j)}x_j=0,\qquad i=1,\ldots,q.
\]

зато јест

\[
\begin{aligned}
\sum_{i=1}^{p}s_i(x_i m)
&=s_1(x_1m)+\cdots+s_q(x_qm)
  +\left(\sum_i s_i\varrho_i^{(q+1)}\right)(x_{q+1}m)\\
&\quad+\cdots+\left(\sum_i s_i\varrho_i^{(p)}\right)(x_pm)\\
&=s_1(x_1m)+\cdots+s_q(x_qm)\\
&\quad+s_1\varrho_1^{(q+1)}(x_{q+1}m)+\cdots+s_q\varrho_q^{(q+1)}(x_{q+1}m)
  \qquad\text{Заради својства модула}\\
&\quad+\cdots\\
&\quad\cdots\\
&\quad+s_1\varrho_1^{(p)}(x_pm)+\cdots+s_q\varrho_q^{(p)}(x_pm)\\
&=s_1(x_1m)+\cdots+s_q(x_qm)\\
&\quad+s_1(\varrho_1^{(q+1)}x_{q+1}m)+\cdots+s_q(\varrho_q^{(q+1)}x_{q+1}m)\\
&\quad+\cdots\\
&\quad\cdots\\
&\quad+s_1(\varrho_1^{(p)}x_pm)+\cdots+s_q(\varrho_q^{(p)}x_pm)
  \qquad\begin{gathered}\text{Заради закона}\\ \text{асоциативности }s'(sm)\\ =s's\cdot m\text{ и}\\ x'(xm)=x'x\cdot m\end{gathered}\\
&=s_1\left(x_1m+\varrho_1^{(q+1)}x_{q+1}\cdot m+\cdots+\varrho_1^{(p)}x_p\cdot m\right)+\cdots\\
&=s_1\left(\left(x_1+\varrho_1^{(q+1)}x_{q+1}+\cdots+\varrho_1^{(p)}x_p\right)m\right)+\cdots\\
&=0.
\end{aligned}
\]


Из својств \(\mathfrak S_{\mathsf K}\)-модула од \(\mathfrak M\) дистрибутивне законы \((S+S')m=Sm+S'm,\ S(m+m')=Sm+Sm'\) сут јасне. Јешче трєба обратити внимање на закон асоциативности \(S(S'm)=SS'\cdot m\); зато же дистрибутивне законы сут правдиве, можно ограничити се на случај

\[
sx\cdot(s'x' m)=sx(s'x'\cdot m)
\]

С једној страны, по дефиницији од \(sx\cdot m\), јест

\[
sx\cdot(s'x'\cdot m)=s\cdot(x\cdot(s'\cdot(x'\cdot m)))
\]

а с другој страны

\[
sx s'x'\cdot m=ss'xx'\cdot m=ss'\cdot(xx'\cdot m)
=s\cdot(s'\cdot(x\cdot(x'\cdot m)))
\]

заради двох релациј комутативности \(=s\cdot(x\cdot(s'\cdot(x'\cdot m)))\)

\[
x\cdot s'=s'\cdot x
\]
\[
s'\cdot(x\cdot m)=x\cdot(s'\cdot m)
\]

зато \(\mathfrak M\) наистино јест лєвы \(\mathfrak S_{\mathsf K}\)-модул, очевидно конечны. Обратно, всакы конечны лєвы \(\mathfrak S_{\mathsf K}\)-модул јест реципрочны модул прєдставјења од \(\mathfrak S\) в \(\mathsf K\), зато же очевидно јест конечны \(\mathsf K\)-модул и заради

\[
s(x\cdot m)=sx\cdot m=xs\cdot m=x(s\cdot m)
\]

такоже изполња условје комутативности. (Сравни такоже простєјши доказ в § 24).

Нынє по M.З. § 18 всакы конечны \(\mathfrak S_{\mathsf K}\)-модул јест сума простых модулов, кторе сут операторно изоморфне с простыми лєвыми идеалами од \(\mathfrak S_{\mathsf K}\) по одношењу к операторам из \(\mathfrak S_{\mathsf K}\). Обратно, всакы лєвы идеал од \(\mathfrak S_{\mathsf K}\) јест модул прєдставјења од \(\mathfrak S\) в \(\mathsf K\), а зато же вси лєви идеалы од \(\mathfrak S_{\mathsf K}\) сут операторно изоморфне, имаје мєсто

\textbf{Теорем 2.} \emph{Једина класа иредуцибилных реципрочных прєдставјениј од \(\mathfrak S\) в \(\mathsf K\) јест она, ктора јест створјена простым лєвым идеалом од \(\mathfrak S_{\mathsf K}\).}

Ако то иредуцибилно прєдставјење имаје ступењ \(r\), то ступењ либовољного прєдставјења јест кратны од \(r\); очевидно можно надовезивањем иредуцибилных прєдставјениј за заданы кратны \(rs\) од \(r\) конструовати прєдставјење \(rs\)-того ступења. Тому одповєдаје образовање модула прєдставјења, кторы јест прєма сума од \(s\) простых модулов.

Нынє уж знајемо једно прєдставјење од \(\mathfrak S\) в \(\mathsf K\), јменно оно, кторо јест створјено самым \(\mathfrak S\) како модулом прєдставјења (главно прєдставјење, M. З. § 20); оно уж лежи в \(\mathsf P\). Његов ступењ јест \(n\), зато \(n\) јест дєливы посрєдством \(r\):

\[
n=rt.
\]

Квоциент \(t\) јест легко опрєдєлити: зато же \(r\) јест ранг једного лєвого идеала над \(\mathsf K\), а \(n\) ранг цєлокупного \(\mathfrak S_{\mathsf K}\), \(t\) јест число лєвых идеалов од \(\mathfrak S_{\mathsf K}\), зато \(t^2\) јест ранг од \(\mathfrak S_{\mathsf K}\) над тєлом автоморфизмов његовых простых идеалов.

Специализујемо наше разгледање на случај, когда \(\mathfrak S\) јест комутативно поље.

\textbf{Теорем 3.} \emph{Ако \(\mathfrak S\) јест комутативно поље, то \(\mathfrak S\) јест центар од \(\mathfrak S_{\mathsf K}\).}

\emph{Доказ.} Зато же всакы елемент од \(\mathfrak S\) јест комутативен с всакым елементом од \(\mathfrak S\) и с всакым елементом од \(\mathsf K\), зато такоже с всакым елементом од \(\mathfrak S_{\mathsf K}\), \(\mathfrak S\) належи к центру од \(\mathfrak S_{\mathsf K}\). Елемент од \(\mathfrak S_{\mathsf K}\) јест комутативен с всакым елементом од \(\mathfrak S_{\mathsf K}\) само тогда, ако јест комутативен принајмење с всакым елементом од \(\mathsf K\), то јест ако његови коефициенты належет к центру од \(\mathsf K\), кторы јест \(\mathsf P\), то јест ако он лежи в \(\mathfrak S\), тым јест доказано.

Зато же при комутативных прстенах не имаје разлике меджу реципрочным и прємым изоморфизмом, всако реципрочно прєдставјење можно разгледати такоже прємым. Иредуцибилно прєдставјење од \(\mathfrak S\) в \(\mathsf K\) мора зато такоже быти дано прємым модулом прєдставјења. Тој прємы модул прєдставјења јест просты правы идеал од \(\mathfrak S_{\mathsf K}\). Зато же \(\mathfrak r\) јест конечны правы \(\mathsf K\)-модул ранга \(r\) и заради \(\mathfrak r\mathfrak S=\mathfrak S\mathfrak r\) такоже лєвы \(\mathfrak S\)-модул (саморазумно с смјешаным законом асоциативности), зато прємы модул прєдставјења ранга \(r\); а зато же обстаја само једна класа иредуцибилных прєдставјениј, он мора стварјати дано прєдставјење.
% END BOOK_S19

% BEGIN BOOK_S20
\subsection*{§ 20. Некомутативне тєла}\label{nichtkommutative-koerper}

Нынє под

\[
\mathfrak K=y_1\mathsf P+\cdots+y_m\mathsf P
\]

разумєјемо некомутативно тєло с центром \(\mathsf P\). За \(\mathfrak K\) развијемо ред теоремов, аналогных теоремам из теорије комутативных конечных алгебраичных разширјениј пољ.

Ако \(Z=\xi_1\mathsf P+\cdots+\xi_n\mathsf P\) означаје хиперкомплексны систем над \(\mathsf P\), два разширјена система \(\mathfrak K_Z\) и \(Z_{\mathfrak K}\) сут идентичне, зато же за \(\mathfrak K_Z\) можно записати

\[
\mathfrak K_Z=y_1\xi_1\mathsf P+\cdots+y_m\xi_n\mathsf P
\]

а за \(Z_{\mathfrak K}\)

\[
Z_{\mathfrak K}=\xi_1y_1\mathsf P+\cdots+\xi_ny_m\mathsf P
\]

але оба израза сут идентичне заради \(y_i\xi_j=\xi_jy_i\).

В слєдујучем \(Z\) всечасно буде комутативно поље, зато можно примєнити теоремы из § 19. По теорему 1 из § 19 \(\mathfrak K_Z\) јест всечасно двустранно просто.

\textbf{Теорем 1.} \emph{Нехај \(\Omega\) буде алгебраично затворјено разширјено поље од \(\mathsf P\). Тогда центар од \(\mathfrak K_\Omega\) јест равны \(\Omega\).}

\emph{Доказ.} Ослањамо се на одговарјајучи теорем 3 из § 19. Јасно јест, же \(\Omega\) лежи в центру од \(\mathfrak K_\Omega\). Ако \(\alpha\) јест либовољны елемент центра од \(\mathfrak K_\Omega\), његови коефициенты належет к некому конечному алгебраичному разширјеному пољу \(Z\) од \(\mathsf P\) (јменно к пољу, кторо они стварјајут). Саморазумно \(\alpha\) такоже мора належати к центру од \(\mathfrak K_Z\). Але тој центар јест \(Z\), зато \(\alpha\) належи к \(Z\), а зато такоже к \(\Omega\), тым јест доказано.

На совсєм подобны начин доказујемо

\textbf{Теорем 2.} \emph{\(\mathfrak K_\Omega\) јест двустранно просто.}

\emph{Доказ.} Нехај \(\mathfrak a=z_1\Omega+\cdots+z_r\Omega\) буде двустранны идеал. Коефициенты од \(z_1,\ldots,z_r\) належет к некому конечному разширјеному пољу \(Z\) од \(\mathsf P\); в \(\mathfrak K_Z\) јест \(\mathfrak a'=z_1Z+\cdots+z_rZ\) двустранны идеал, а зато же \(\mathfrak K_Z\) јест двустранно просто, \(\mathfrak a'\) јест или нула или цєло \(\mathfrak K_Z\); зато \(\mathfrak a\) јест или 0 или цєло \(\mathfrak K_\Omega\), т.~ј. \(\mathfrak K_\Omega\) јест двустранно просто, тым јест доказано.

По том теорему \(\mathfrak K_\Omega\) јест матрицно прстен над некым разширјеным пољем својего центра \(\Omega\). То поље како подпрстен цєлого \(\mathfrak K_\Omega\) мора имєти конечны ступењ над \(\Omega\), зато мора быти равно \(\Omega\):

\textbf{Теорем 3.} \emph{\(\mathfrak K_\Omega\) јест матрицно прстен над \(\Omega\),}

\[
\mathfrak K_\Omega=\sum\Omega c_{ik}
\]

\emph{зато ступењ од \(\mathfrak K\) над његовым центром \(\mathsf P\) јест квадратно число \(m=t^2\).}

Число \(t\) правых (или лєвых) идеалов, на кторе разпада се \(\mathfrak K_\Omega\), называјмо кратко абсолутным числом компонентов од \(\mathfrak K\).

Својство од \(\Omega\), же \(\mathfrak K_\Omega\) разпада се на \(t\) простых идеалов, очевидно имајут уж конечна разширјења \(Z\) од \(\mathsf P\), напр. така \(Z\), кторе се добывајут приложењем коефициентов од \(c_{ik}\) к \(\mathsf P\). Нынє разгледајемо та разпадна поља.

\textbf{Дефиниција.} Конечно алгебраично разширјено поље \(Z\) од \(\mathsf P\) называје се \emph{разпадным пољем} од \(\mathfrak K\), ако число простых правых идеалов, на кторе разпада се \(\mathfrak K_Z\), уж јест равно абсолутному числу компонентов \(t\)\srcfn{1}{Сравни појем разпадного поља на стр.~9}.

\textbf{Теорем 4.} \emph{Ако \(Z\) јест разпадно поље од \(\mathfrak K\), то \(Z\) јест тєло автоморфизмов правых идеалов од \(\mathfrak K_Z\); обратно, ако тєло автоморфизмов правых идеалов од \(\mathfrak K_Z\) јест равно \(Z\), то \(Z\) јест разпадно поље.}

\emph{Доказ.} 1. Нехај \(Z\) буде разпадно поље, а \(\mathfrak T\) тєло автоморфизмов правых идеалов од \(\mathfrak K_Z\), то јест

\[
\mathfrak K_Z=\sum\mathfrak T c_{ik}
\]

Из тога слєдује

\[
\mathfrak K_\Omega=\sum\mathfrak T_\Omega c_{ik}
\]

ранг од \(\mathfrak K_\Omega\) над \(\Omega\) јест зато продукт од \(t^2\) с рангом од \(\mathfrak T_\Omega\); с другој страны он јест равны \(t^2\), зато ранг од \(\mathfrak T_\Omega\), а такоже од \(\mathfrak T\), јест равны 1, \(\mathfrak T=Z\), тым јест доказано.

\textbf{2.} Ако \(Z\) јест тєло автоморфизмов простых идеалов од \(\mathfrak K_Z\), то јест

\[
\mathfrak K_Z=\sum Zc_{ik}
\]

зато

\[
\mathfrak K_\Omega=\sum\Omega c_{ik}
\]

число елементов \(c_{ik}\) мора зато быти равно \(t^2\), тым јест доказано.

\textbf{Теорем 5.} \emph{Иредуцибилно прєдставјење \(\mathfrak Z\) разпадного поља \(Z\) од \(\mathfrak K\) в \(\mathfrak K\) -- по теорему 2 на стр.~22 обстаја само једна класа иредуцибилных прєдставјениј -- јест максимално комутативно подпоље матрицного прстена \(\mathfrak K_r\) (\(r\) јест ступењ иредуцибилного прєдставјења од \(Z\)).}

\emph{Ако обратно иредуцибилно прєдставјење \(\mathfrak Z\) комутативного разширјеного поља \(Z\) од \(\mathsf P\) јест максимално комутативно подпоље од \(\mathfrak K_r\), то \(Z\) јест разпадно}

\emph{поље.}

\emph{Доказ.} 1. Ако \(n\) јест ступењ разпадного поља \(Z\) над \(\mathsf P\), то ступењ \(r\) иредуцибилного прєдставјења \(\mathfrak Z\) од \(Z\) јест равны \(n/t\) (стр.~22). Нехај нынє \(\mathfrak Z\subseteq\mathfrak Z^*\subseteq\mathfrak K_r\), а \(\mathfrak Z^*\) буде максимално комутативно поље ступења \(n^*\geq n\) над \(\mathsf P\). Трєба доказати \(\mathfrak Z^*=\mathfrak Z\). Вмєсто тога достаточно јест доказати \(n^*=n\). Ако \(r^*\) означаје ступењ иредуцибилного прєдставјења с \(\mathfrak Z^*\) изоморфного разширјеного поља \(Z^*\) од \(Z\), поновно јест \(n^*=r^*t\); с другој страны \(r\) јест дєливы посрєдством \(r^*\), зато же в \(\mathfrak Z^*\) јест дано прєдставјење \(r\)-того ступења од \(Z^*\), зато \(r^*\leq r\). Из

\[
n\leq n^*=r^*t\leq rt=n \qquad\text{слєдује}\qquad n^*=n.
\]

2. Нехај иредуцибилно прєдставјење \(\mathfrak Z\) \(r\)-того ступења комутативного разширјеного поља \(Z\) од \(\mathsf P\) в \(\mathfrak K\) буде максимално комутативно подпоље од \(\mathfrak K_r\). Нехај тєло автоморфизмов простых идеалов од \(\mathfrak K_Z\) буде \(\mathfrak T\), то јест

\[
\mathfrak K_Z=\sum_{1}^{t'}\mathfrak T c_{ik}.
\]

\(\mathfrak T\) содрживаје \(Z\). Трєба доказати \(\mathfrak T=Z\). Нынє просты правы идеал \(\mathfrak r\) од \(\mathfrak K_r\) не јест само модул прєдставјења од \(Z\) в \(\mathfrak K\), але такоже модул прєдставјења од \(\mathfrak T\), зато же јест

\[
\mathfrak r=\mathfrak T c_{i1}+\cdots+\mathfrak T c_{ik},
\]

зато \(\mathfrak T\mathfrak r=\mathfrak r\); \(\mathfrak r\) јест лєвы \(\mathfrak T\)-модул и очевидно правы \(\mathfrak K\)-модул.

Обстаја зато с \(\mathfrak T\) изоморфно подтєло \(\overline{\mathfrak T}\) од \(\mathfrak K_r\), кторо содрживаје \(\mathfrak Z\). Ако бы \(\mathfrak T\) было различно од \(Z\), то \(\overline{\mathfrak T}\) различно од \(\mathfrak Z\), приложење не належечего к \(\mathfrak Z\) елемента од \(\overline{\mathfrak T}\), \(\alpha\), к \(\mathfrak Z\) дало бы право комутативно разширјено поље од \(\mathfrak Z\) в \(\mathfrak K_r\), а то противрєчи максималности од \(\mathfrak Z\), тым јест доказано.

Из теорема 5 безпосрєдње слєдује

\textbf{Теорем 6.} \emph{Максимална комутативна подпоља од \(\mathfrak K\) имајут ступењ \(t\) над \(\mathsf P\) и сут разпадна поља\srcfn{1}{Прєдположење, же всако разпадно поље такоже содрживаје неко максимално комутативно подпоље од \(\mathfrak K\), не јест правилно; могут обстајати разпадна поља без тога својства, а одговарјајучи \(r\) сигурно јест выше од 1. Види: \foreign{R. Brauer} и \foreign{E. Noether}, \foreign{Über minimale Zerfällungskörper irreduzibler Darstellungen; Sitzungsberichte d. preuss. Akademie der Wisssch. 1927 XXXII, str.~221.}}.}

Нынє развијемо теорем о изоморфизму подпрстенов једного \(\mathfrak K_r\), на ктором основана јест подробнејша знајемост структуры Галоисовој групы од \(\mathfrak K\) и с тым свезаных фактов.

\textbf{Теорем 7.} \emph{Нехај \(\mathfrak S_1\) и \(\mathfrak S_2\) будут двустранно просте, \(\mathsf P\) содрживајуче подпрстены матрицного прстена \(\mathfrak K_r\) над \(\mathfrak K\), и нехај меджу ними обстаја изоморфизм, оставјајучи \(\mathsf P\) по елементах неподвижным, \(\mathfrak S_1\cong\mathfrak S_2\). Тогда тој изоморфизм јест трансформација регуларным \(\mathfrak K_r\)-елементом \(\tau\).}

\emph{Ако \(\sigma_1\in\mathfrak S_1\) и \(\sigma_2\in\mathfrak S_2\) сут взајемно приреджене, то јест}

\[
\tau^{-1}\sigma_1\tau=\sigma_2.
\]

\emph{Доказ.} Конструујемо с \(\mathfrak S_1\) и \(\mathfrak S_2\) реципрочно изоморфно прстен \(\Sigma\), чије подпоље изоморфно с \(\mathsf P\) идентификујемо с \(\mathsf P\). \(\mathfrak S_1,\mathfrak S_2\) сут тогда два реципрочна прєдставјења једнакого ступења \(r\) од \(\Sigma\) в \(\mathfrak K\). Зато же по теорему 2 на стр.~22 обстаја само једна класа иредуцибилных реципрочных прєдставјениј, зато такоже само једна класа прєдставјениј заданого ступења од \(\Sigma\) в \(\mathfrak K\) (кратны од иредуцибилној класы), \(\mathfrak S_1\) и \(\mathfrak S_2\) морајут належати к истој класе; зато једно настава из другого трансформацијоју, тым јест доказано.

Из теорема 7 непосрєдње слєдује

\textbf{Теорем 8.} \emph{Група автоморфизмов, кторе оставјајут \(\mathsf P\) по елементах неподвижным, од \(\mathfrak K\) -- т.~ј. Галоисова група \(\mathfrak G\) од \(\mathfrak K\) по одношењу к његовому центру -- јест група внутрных автоморфизмов од \(\mathfrak K\).}

В теорему 7 трєба само положити \(r=1\) и \(\mathfrak S_1=\mathfrak S_2=\mathfrak K\).

Слєдујуча два теорема показывајут употрєбљивост досєга развитој обчеј теорије за разгледање специалных пытаниј.

\textbf{Теорем 9.} \emph{Једино некомутативно тєло конечного ступења, чији центар јест поље \(R\) реалных числ, јест тєло кватернионов.}

\emph{Доказ.} Нехај \(\mathfrak K\) буде некомутативно тєло с центром \(R\) (конечного ранга над \(R\)). По теорему 6 ступењ од \(\mathfrak K\) над \(R\) јест квадрат ступења \(t\) некого максималного комутативного подпоља од \(\mathfrak K\); а зато же оно може быти само поље изоморфно с пољем комплексных числ, јест \(t=2\), а \(\mathfrak K\) имаје ступењ 4 над \(R\). Нехај \(\mathfrak Z_1\) буде максимално комутативно подпоље од \(\mathfrak K\); тогда можно положити \(\mathfrak Z_1=R(i)\), кде \(i^2=-1\). Але тогда такоже \(\mathfrak Z_2=R(-i)\) јест -- како множство сбєгајуче се с \(\mathfrak Z_1\) -- поље изоморфно с \(\mathfrak Z_1\), а изо-

морфизм јест заданы посрєдством \(i\mapsto -i\). По теорему 7 обстаја зато елемент \(j^*\in\mathfrak K\) с

\[
j^{*-1}ij^*=-i.
\]

Ранг од \(\mathfrak Z_1(j^*)\) над \(\mathfrak Z_1\) јест нынє принајмење 2, зато же \(j^*\) не може належати к \(\mathfrak Z_1\). зато \(\mathfrak Z_1(j^*)\) имаје ранг 4 над \(R\), јест равно \(\mathfrak K\), и можно положити:

\[
\mathfrak K=R+Ri+Rj^*+Rij^*.
\]

Заради

\[
\begin{aligned}
j^{*-2}1j^{*2}&=1,\\
j^{*-2}ij^{*2}&=i,\\
j^{*-2}j^*j^{*2}&=j^*,\\
j^{*-2}ij^*j^{*2}&=ij^*
\end{aligned}
\]

\(j^{*2}\) јест комутативен с всими елементами од \(\mathfrak K\), належи к центру \(R\), а \(j^{*2}\) јест реално число \(r\). \(r\) не може быти позитивно, зато же иначе полином \(x^2-r\) в комутативном пољу \(R(j^*)\) имел бы четыри корења \(+\sqrt r,-\sqrt r,+j^*\) и \(-j^*\). зато \(j^{*2}=-m^2,\ m\in R\). Положимо \(j=j^*m^{-1}\); тогда \(j^2=-1\), а \(1,i,j,k=ij\) сут јединице кватернионов, тым јест доказано.

Доказ очевидно имаје мєсто за всако алгебраично реално затворјено поље вмєсто \(R\).

\textbf{Теорем 10 (теорем од Веддербурна).} \emph{Всако тєло из конечного числа елементов јест комутативно.}

\emph{Доказ.} Нехај \(\mathfrak K\) буде тєло с конечным числом елементов, а \(\mathsf P\) јего центар. Нехај \(t^2\) буде ступењ од \(\mathfrak K\) над \(\mathsf P\). Вса максимална комутативна подпоља од \(\mathfrak K\) имајут ступењ \(t\) над \(\mathsf P\), зато имајут равнако много елементов; по познатом теорему из теорије Галоисовых пољ (теорем од Моореа) она сут изоморфне, а те изоморфизмы, зато же всако Галоисово поље јест Галоисово разширјење својего простого поља, можно избрати тако, же всакы елемент од \(\mathsf P\) одповєдаје самому себе. По теорему 7 вса максимална комутативна подпоља од \(\mathfrak K\) наставајут зато из једного \(\mathfrak Z\) трансформацијоју. Але \(\mathfrak K\) јест множство сјединења всих својих максималных комутативных подпољ. (Приложење \emph{једного} елемента к \(\mathsf P\) стварјаје комутативно поље). Ако нынє разгледамо само мултипликативну групу \(\mathfrak K^*\) од \(\mathfrak K\) (\(\mathfrak K\) без нулы), достали јесмо, же \(\mathfrak K^*\) јест множство сјединења подгруп коњугованых подгрупє \(\mathfrak Z^*\). Але в случају \(t>1\) то даваје противрєчност. Зато же ако \(\mathfrak K^*\) содрживаје \(M\) елементов, \(\mathfrak Z^*\) \(N\) елементов, а \(L\) јест индекс од \(\mathfrak Z^*\) в \(\mathfrak K^*\), то имаје мєсто \(M=NL\), а число различных коњугованых подгруп јест највыше \(L\). Чак ако то число было бы равно \(L\), да бы все коњуговане подгрупы заједно изчрпале цєлу групу, никаке двє из њих не смєли бы имєти обчи елемент. Але все имајут обчи елемент јединице.
% END BOOK_S20

% BEGIN BOOK_S21
\subsection*{§ 21. Галоисова теорија некомутативных тєл}\label{die-galoissche-theorie-der-nichtkommutativen-koerper}

Под \(G\) разумєјемо Галоисову групу од \(\mathfrak K\) по одношењу к његовому центру \(\mathsf P\). По теорему 8 \(G\) јест изоморфна с \(\overline G=\mathfrak K^*/\mathsf P^*\); разгледамо тој изоморфизм установјеным. Подгрупу \(H\) од \(G\) называмо \emph{затворјеној}, ако посрєдством изоморфизма \(G\simeq\mathfrak K^*/\mathsf P^*\) јеј приреджена подгрупа \(\mathfrak H^*\) од \(\mathfrak K^*\) по додању \emph{нулы} поставаје тєлом.


\textbf{Теорем 11. Главны теорем Галоисовој теорије.} \emph{Затворјене подгрупы \(H\) од \(G\) и тєла \(\mathfrak S\) меджу \(\mathsf P\) и \(\mathfrak K\) можно једнозначно взајемно приредити тако, же ако \(H\) и \(\mathfrak S\) одповєдајут једно другому, \(H\) јест инвариантна група од \(\mathfrak S\), а \(\mathfrak S\) инвариантно тєло од \(H\).}

\emph{Доказ.} Раздєљамо наш теорем на три честне теоремы
\begin{enumerate}
\item Инвариантна група \(H\) тєла \(\mathfrak S\) јест затворјена.
\item Ако \(\mathfrak S\) имаје инвариантну групу \(H\), то \(\mathfrak S\) јест инвариантно тєло од \(H\).
\item Ако \(H\) имаје инвариантно тєло \(\mathfrak S\), то \(H\) јест инвариантна група од \(\mathfrak S\).
\end{enumerate}

\begin{enumerate}
\item Инвариантној групє \(H\) тєла \(\mathfrak S\) посрєдством изоморфизма \(G\simeq \mathfrak K^*/\mathsf P^*\) приреджено јест множство \(\mathfrak H^*\) всих ненуловых елементов од \(\mathfrak K\), кторе сут по елементах комутативне с \(\mathfrak S\). Заједно с 0 она очевидно творет тєло, т.~ј. \(H\) јест затворјена.

\item Трети честны теорем сводимо на вторы. Тєлу \(\mathfrak S\) приреджена јест подгрупа \(S\) од \(G\), јменно подгрупа, приналежеча к \(\mathfrak S^*/\mathsf P^*\), од \(G\). Так само групє \(H\) приналежи подтєло \(\mathfrak H\) од \(\mathfrak K\), и јест \(H\simeq \mathfrak H^*/\mathsf P^*\). Ако \(\mathfrak S\) јест инвариантно тєло од \(H\), то \(\mathfrak S\) состоји се из всих с \(\mathfrak H\) по елементах комутативных елементов од \(\mathfrak K\); можно то јест такоже казати: \(S\) јест инвариантна група од \(\mathfrak H\). Аналогично: Ако \(H\) јест инвариантна група од \(\mathfrak S\), то \(\mathfrak H\) јест инвариантно тєло од \(S\). Зато трети честны теорем можно формуловати такоже овако:

\item[\(3'\).] Ако \(\mathfrak H\) имаје инвариантну групу \(S\), то \(\mathfrak H\) јест инвариантно тєло од \(S\), а то јест равно второму честному теорему, кторы јешче трєба доказати.

\item[3.] Тој доказ тече, кромє нєколико малых разлика, како в комутативном случају:
\end{enumerate}

Конструујемо с \(\mathfrak K\) реципрочно изоморфно тєло \(K\), чији центар идентификујемо с \(\mathsf P\). Вмєсто автоморфизмов од \(\mathfrak K\) можно разгледати реципрочне изоморфизмы од \(\mathfrak S\) с \(K\).

Поновно хочемо доказати помочну лему, аналогну теорему о продолжењу, употрєбјеному в комутативном случају:

Нехај \(\mathfrak S\) и \(\mathfrak T\) будут тєла меджу \(\mathsf P\) и \(\mathfrak K\),
\[
\mathsf P \subseteq \mathfrak S \subseteq \mathfrak T \subseteq \mathfrak K.
\]
Тврдимо, же всакы реципрочны изоморфизм од \(\mathfrak S\) на подтєло од \(K\) (реципрочно прєдставјење првого ступења од \(\mathfrak S\) в \(K\)) можно на принајмење толико начинов продолжити до реципрочных изоморфизмов од \(\mathfrak T\), колико указује ступењ \(s\) од \(\mathfrak T\) над \(\mathfrak S\).

Тєла \(\mathsf P\), \(\mathfrak S\), \(\mathfrak T\), \(\mathfrak K\) разширјајут се с \(K\)
\[
K\subseteq \mathfrak S_K\subseteq \mathfrak T_K\subseteq \mathfrak K_K.
\]
По теорему 1 на стр.~20 всако \(\mathfrak S_K\) јест двустранно просто и полно редуцибилно, а прости лєви идеалы -- меджу собоју операторно изоморфне -- сут модулы прєдставјења за једину класу иредуцибилных прєдставјениј од \(\mathfrak S\) в \(K\). Зато же та класа прєдставјениј имаје ступењ 1 -- зато же по конструкцији од \(K\) в \(K\) обстајајут тєла реципрочно изоморфне с \(\mathfrak S\) -- вси прости лєви идеалы имајут ранг 1, т.~ј. систем јест сума толико лєвых идеалов, колико указује његов ступењ над \(\mathsf P\)
\[
\mathfrak S_K=K E_1+\cdots+K E_p,\qquad K E_\nu=\mathfrak S_K E_\nu.
\]

Тогда јест
\[
\begin{aligned}
\mathfrak T_K&=\mathfrak T\cdot K=\mathfrak T\cdot \mathfrak S_K,\\
\mathfrak T_K&=\mathfrak T\mathfrak S_K E_1+\cdots+\mathfrak T\mathfrak S_K E_p
              =\mathfrak T_K E_1+\cdots+\mathfrak T_K E_p.
\end{aligned}
\]
Всакы идеал \(\mathfrak T_K\cdot E_\nu\) разпада се на \(s\) простых идеалов, зато же
\[
[\mathfrak T_K E_\nu:K]=[\mathfrak T_K E_\nu:K E_\nu]=[\mathfrak T_K E_\nu:\mathfrak S_K E_\nu]\leqq[\mathfrak T_K:\mathfrak S_K]=[\mathfrak T:\mathfrak S]=s
\]
и
\[
\sum [\mathfrak T_K E_\nu:K]=sp
\]
зато
\[
[\mathfrak T_K E_\nu:K]=s,
\qquad
\mathfrak T_K E_\nu=K e_1^{(\nu)}+\cdots+K e_s^{(\nu)}.
\]

Прєдставјење, дано посрєдством \(K e_i^{(\nu)}\), од \(\mathfrak T\) јест продолжење прєдставјења, заданого посрєдством \(K E_\nu\), од \(\mathfrak S\), зато же ако \(\alpha\in \mathfrak S\), то јест
\[
\alpha E_\nu=\sigma E_\nu,\qquad
\sum \alpha e_i^{(\nu)}=\sigma\sum e_i^{(\nu)},\qquad
\alpha e_i^{(\nu)}=\sigma e_i^{(\nu)}.
\]

Нынє трєба доказати само, же прєдставјења дане посрєдством \(K e_i^{(\nu)}\) сут все различне. В комутативном случају то было јасно, зато же она належали к различным класам прєдставјениј. Ту она, насупрот, все належет к истој класе прєдставјениј, зато трєба разгледати саме прєдставјења. За то разгледајемо дане посрєдством \(K e_i^{(\nu)}\) \(K\)-хомоморфизмы цєлого \(\mathfrak T_K\); они сут уж цєлком задане, ако знајемо прєдставјења од \(\mathfrak T\), зато же за конструкцију цєлого хомоморфизма достаточно јест знати, како сут прєдставјене базове елементи. Можно зато ограничити се на доказ, же посрєдством \(K e_i^{(\nu)}\) дане хомоморфизмы од \(\mathfrak T_K\) сут все различне. Ако \(e_i^{(\nu)}\) једен раз прєдставимо посрєдством модула прєдставјења \(K e_i^{(\nu)}\), а други раз посрєдством \(K e_j^{(\nu)}\), тогда заради
\[
e_i^{(\nu)2}=e_i^{(\nu)}
\]
првы раз најдемо 1, а заради
\[
e_i^{(\nu)}e_j^{(\nu)}=0
\]
други раз 0 како прєдставјајучи елемент од \(e_i^{(\nu)}\); оба прєдставјења сут зато наистино различне. (Точно теже закључкы можно примєнити такоже в комутативном случају, але там сут непотрєбне).

Од нынє закључујемо точно како на стр.~10 и стр.~16. Нехај \(H\) буде инвариантна група од \(\mathfrak S\), а \(\mathfrak T\) инвариантно тєло од \(H\). Елементи од \(H\) сут вса продолжења идентичного автоморфизма од \(\mathfrak S\) на цєло \(\mathfrak K\). Зато же по управо доказаној помочној леми меджу ними обстајајут принајмење \(s\) различне на \(\mathfrak T\), мора быти \(s=1\), \(\mathfrak T=\mathfrak S\), тым јест доказано.

\section*{Глава V. Факторне системы}\label{kapitel-v.-faktorensysteme}
% END BOOK_S21

% BEGIN BOOK_S22
\subsection*{§ 22. Група тєл с даным центром}\label{die-gruppe-der-koerper-mit-gegebenem-zentrum}

За основу нашего разгледања изберемо установјено комутативно поље \(\mathsf P\). Множство всих матрицных прстенов \(\mathfrak K_r\), чија тєла автоморфизмов \(\mathfrak K\) сут тєла конечного ранга над \(\mathsf P\) с самым \(\mathsf P\) како центром, јест мултипликативно затворјены систем по одношењу к ,,образовању прємого продукта``. Прємы продукт двох такых прстенов \(\mathfrak K_r\) и \(\mathfrak L_s\) јест просто прстен \(\mathfrak K_r\cdot\mathfrak L_s\), кторо настава из \(\mathfrak K_r\) разширјењем основного поља до \(\mathfrak L_s\):
\[
\mathfrak K_r\times\mathfrak L_s=\mathfrak K_r\cdot\mathfrak L_s.
\]
По припомнєнке на стр.~22 то образовање прємого продукта јест комутативно:
\[
\mathfrak K_r\times\mathfrak L_s=\mathfrak L_s\times\mathfrak K_r.
\]
Тврджење, же систем всих \(\mathfrak K_r\) јест мултипликативно затворјены, не значи ничто друго неж, же \(\mathfrak K_r\times\mathfrak L_s=\mathfrak K_r\mathfrak L_s\) јест всечасно поновно матрицно прстен, чије тєло автоморфизмов \(\mathfrak M\) имаје центар \(\mathsf P\) -- конечны ранг по одношењу к \(\mathsf P\) разумєје се сам собоју. По теорему 1 на стр.~20 \(\mathfrak K_r\times\mathfrak L_s\) јест двустранно просто и полно редуцибилно, зато матрицно прстен над некым разширјеным тєлом \(\mathfrak M\) од \(\mathsf P\). Факт, же центар од \(\mathfrak M\) јест равны \(\mathsf P\), непосрєдње слєдује из тога, же прво \(\mathsf P\) сигурно јест содрживано в центру, а друго оно мора быти содрживано в центру \(\mathsf P\) од \(\mathfrak K\) (или од \(\mathfrak L\)). Множење \(\mathfrak K_r\times\mathfrak L_s\) очевидно јест асоциативно. \(\mathfrak K_r\) творет мултипликативно затворјены систем, але не групу, зато же ако \(\mathfrak K_r\times\mathfrak L_s=\mathfrak M_m\), то сигурно \(m\geq rs\); зато не буде можно најдти \(\mathfrak L_s\), кторо изполња једначење \(\mathfrak K_r\times\mathfrak L_s=\mathfrak M_m\), ако \(r>m\).

Нынє сдружујемо все \(\mathfrak K_r\) с једнакым \(\mathfrak K\) в класу \(\{\mathfrak K\}\). Тогда имаје мєсто: Ако \(\mathfrak K_r\) и \(\mathfrak K_{r'}\) належет к истој класе \(\{\mathfrak K\}\), а \(\mathfrak L_s\) и \(\mathfrak L_{s'}\) к истој класе \(\{\mathfrak L\}\), то такоже \(\mathfrak K_r\times\mathfrak L_s\) и \(\mathfrak K_{r'}\times\mathfrak L_{s'}\) належет к истој класе. Достаточно јест показати, же \(\mathfrak K_r\times\mathfrak L_s\) належи к истој класе како \(\mathfrak K\times\mathfrak L\), то јест же ако \(\mathfrak K\times\mathfrak L=\mathfrak M_m\), то такоже \(\mathfrak K_r\times\mathfrak L_s\) јест матрицно прстен над \(\mathfrak M\). То јест легко видєти: Јест
\[
\mathfrak K_r\times\mathfrak L
=\left(\sum_1^r\mathfrak K c_{ik}\right)\mathfrak L
=\sum_1^r\mathfrak K_{\mathfrak L}c_{ik}
=\sum_1^r\mathfrak M_m c_{ik}.
\]
Матрицно прстен \(r\)-того ступења над \(\mathfrak M_m\), а то очевидно јест матрицно прстен \(mr\)-того ступења над \(\mathfrak M\). Тогда продолжујемо точно тако
\[
\mathfrak K_r\times\mathfrak L_s
=\left(\sum_1^s\mathfrak L d_{ik}\right)_{\mathfrak K_r}
=\sum_1^s\mathfrak L_{\mathfrak K_r}d_{ik}
=\sum_1^s\mathfrak M_{mr}d_{ik}
=\mathfrak M_{mrs}
\]
чим наше тврджење уж јест доказано. Запомнимо, же из \(\mathfrak L\times\mathfrak K=\mathfrak M_m\) слєдује \(\mathfrak K_r\times\mathfrak L_s=\mathfrak M_{mrs}\).

Из доказаного слєдује: ако продукт двох клас \(\{\mathfrak K\}\) и \(\{\mathfrak L\}\) дефинујемо како класу продукта једного репрезентанта од \(\{\mathfrak K\}\) с једним репрезентантом од \(\{\mathfrak L\}\), та дефиниција не зависи од избора репрезентантов.

Приреджење \(\mathfrak K_r\mapsto\{\mathfrak K\}\) јест зато хомоморфизм система всих \(\mathfrak K_r\) на множство всих клас \(\{\mathfrak K\}\); множство \(\mathscr K\) клас \(\{\mathfrak K\}\) јест зато мултипликативно затворјены систем. Доказујемо, же \(\mathscr K\) јест група. Како се непосрєдње види, елемент јединице јест класа \(\{\mathsf P\}\), \(\mathfrak K_r\times\mathsf P=\mathfrak K_r\).

Остаје само доказати за всаку класу \(\{\mathfrak K\}\) обстајање класы \(\{\mathfrak K\}^{-1}\). Показује се, же трєба положити \(\{\mathfrak K\}^{-1}=\{\overline{\mathfrak K}\}\), кде \(\{\overline{\mathfrak K}\}\) означаје тєло реципрочно изоморфно с \(\mathfrak K\). Другыми словами: \(\mathfrak K\times\overline{\mathfrak K}\) мора быти матрицно прстен над \(\mathsf P\), или ако \(\mathfrak K\) имаје ранг \(t^2\), зато \(\mathfrak K\times\overline{\mathfrak K}\) ранг \(t^4\) над \(\mathsf P\), то \(\mathfrak K\times\overline{\mathfrak K}\) мора разпадати се на \(t^2\) простых лєвых идеалов. То слєдује из тога, же просты лєвы идеал од \(\mathfrak K\times\overline{\mathfrak K}\), како модул прєдставјења некого иредуцибилного реципрочного прєдставјења од \(\overline{\mathfrak K}\) в \(\mathfrak K\) -- кторо јест првого ступења -- имаје ранг 1 над \(\mathfrak K\), то јест ранг \(t^2\) над \(\mathsf P\).

\(\mathfrak K\) хочемо называти \emph{простым тєлом}, ако меджу \(\mathsf P\) и \(\mathfrak K\) не имаје тєла, чији центар такоже јест \(\mathsf P\).

\textbf{Теорем.} \emph{Всако тєло \(\mathfrak K\) јест прємы продукт простых подтєл \(\mathfrak K^{(i)}\),}
\[
\mathfrak K=\mathfrak K^{(1)}\times\cdots\times\mathfrak K^{(p)}.
\]

\emph{Доказ.} Достаточно јест показати, же ако \(\mathfrak K^{(1)}\) јест подтєло од \(\mathfrak K\) с центром \(\mathsf P\), то обстаја подтєло \(\mathfrak S\) од \(\mathfrak K\) с центром \(\mathsf P\), кторо изполња једначење
\[
\mathfrak K^{(1)}\times\mathfrak S=\mathfrak K
\]
В всаком случају обстаја класа \(\{\mathfrak S\}\) с својством
\[
\{\mathfrak K^{(1)}\}\times\{\mathfrak S\}=\{\mathfrak K\}
\]
јменно
\[
\{\mathfrak S\}=\{\overline{\mathfrak K^{(1)}}\}\times\{\mathfrak K\}.
\qquad
\text{Нехај буде}\qquad
\overline{\mathfrak K^{(1)}}\times\mathfrak K=\mathfrak S_\lambda.
\]
Тогда јест
\[
\mathfrak K^{(1)}\times\mathfrak S_\lambda
=\mathfrak K^{(1)}\times\overline{\mathfrak K^{(1)}}\times\mathfrak K
=\mathsf P_{t^2}\times\mathfrak K
=\mathfrak K_{t^2},
\qquad
t^2=\text{ранг од }\mathfrak K^{(1)}.
\]
Але \(\mathfrak K^{(1)}\times\mathfrak S_\lambda\) јест матрицно прстен ступења принајмење равного \(\lambda\), зато \(\lambda\leq t^2\). \(\mathfrak S_\lambda=\overline{\mathfrak K^{(1)}}\times\mathfrak K\) мора разпадати се на просте идеалы ранга 1 над \(\mathfrak K\), зато же \(\overline{\mathfrak K^{(1)}}\) допушча реципрочна прєдставјења првого ступења в \(\mathfrak K\); зато \(\lambda\geq t^2\), \(\lambda=t^2\),
\[
\mathfrak K_{t^2}=\mathfrak K^{(1)}\times\mathfrak S_{t^2}.
\]
зато само \(\mathfrak K^{(1)}\times\mathfrak S\) мора быти тєло, и
\[
\mathfrak K^{(1)}\times\mathfrak S_{t^2}
=\bigl(\mathfrak K^{(1)}\times\mathfrak S\bigr)_{t^2}
=\mathfrak K_{t^2}.
\]
\(\mathfrak K^{(1)}\times\mathfrak S\) јест тєло автоморфизмов простых лєвых идеалов в једном разкладу од \(\mathfrak K_{t^2}\), а \(\mathfrak K\) тєло автоморфизмов простых лєвых идеалов в другом разкладу од \(\mathfrak K_{t^2}\). Але зато же всака два различна лєва идеала од \(\mathfrak K_{t^2}\) можно вмєстити в једен разклад, \(\mathfrak K\) и \(\mathfrak K^{(1)}\times\mathfrak S\) сут изоморфне, а тој изоморфизм можно избрати тако, же \(\mathsf P\) оставаје по елементах неподвижно. Зато по теорему 7 на стр.~25 тој изоморфизм јест створјен внутрным автоморфизмом од \(\mathfrak K_{t^2}\).
\[
\mathfrak K
=\lambda^{-1}\bigl(\mathfrak K^{(1)}\times\mathfrak S\bigr)\lambda
=\lambda^{-1}\mathfrak K^{(1)}\lambda\times\lambda^{-1}\mathfrak S\lambda,
\qquad
\lambda\in\mathfrak K_{t^2}.
\]
\(\lambda^{-1}\mathfrak K^{(1)}\lambda\) јест с \(\mathfrak K^{(1)}\) изоморфно подтєло од \(\mathfrak K\); зато по теорему 7 на стр.~25 с одговарјајучим \(\mu\) из \(\mathfrak K\):
\[
\mu^{-1}\lambda^{-1}\mathfrak K^{(1)}\lambda\mu=\mathfrak K^{(1)}
\]
зато
\[
\mu^{-1}\mathfrak K\mu
=\mathfrak K
=\mathfrak K^{(1)}\times\mu^{-1}\lambda^{-1}\mathfrak S\lambda\mu,
\]
чим јест тврджење доказано.
% END BOOK_S22

% BEGIN BOOK_S23
\subsection*{§ 23. Факторне системы}\label{faktorensysteme}

В слєдујучих разгледањах прєдпокладамо, же основно поље \(\mathsf P\) јест досконало; вмєсто тога можно учинити обчејше прєдположење, же разгледајут се само така некомутативна \(\mathfrak K\) с центром \(\mathsf P\), кторе имајут својство, же всако \(\mathfrak K_r\) содрживаје максимално комутативно подпоље првого рода над \(\mathsf P\). (Како меджутым показал \foreign{Köthe}, то имаје мєсто за всако \(\mathfrak K\).)

Под \(Z\) разумєјемо разпадно поље некомутативного тєла \(\mathfrak K\) с центром \(\mathsf P\); нехај иредуцибилно прєдставјење \(\mathfrak Z\) од \(Z\) в \(\mathfrak K\) буде ступења \(r\), тогда оно јест максимално комутативно подпоље од \(\mathfrak K_r\). Нехај \(\Gamma\) буде Галоисово поље од \(Z\), зато
\[
\mathfrak Z_\Gamma=e_1\Gamma+\cdots+e_n\Gamma
\qquad\text{(стр.~9)}
\]
\(n\) јест ступењ од \(Z\) над \(\mathsf P\). За с \(\Gamma\) разширјено прстен \(\mathfrak K_r\) добыва се такоже разклад
\[
\mathfrak K_{r\Gamma}=\mathfrak K_{r\Gamma}e_1+\cdots+\mathfrak K_{r\Gamma}e_n
\]
\[
\mathfrak K_{r\Gamma}=\mathfrak l_1+\cdots+\mathfrak l_n,\qquad
\mathfrak l_i=\mathfrak K_{r\Gamma}e_i.
\]
\(\mathfrak l_i\) очевидно сут лєви идеалы. Они сут чак прости лєви идеалы, зато же \(Z\) јест разпадно поље, \(\mathfrak K_{r\Gamma}\) јест матрицно прстен ступења \(r\cdot t=n\) над \(\Gamma\). (Теорем 4, стр.~23, \(n=rt\) на стр.~22.) Тако достаны разклад од \(\mathfrak K_{r\Gamma}\) на просте лєве идеалы особливо јест прилагоджены алгебраичным својствам од \(\mathfrak K\) и \(Z\). Ако, како на стр.~18, дєјство субституције \(S\) из Галоисовој групы \(\mathfrak G\) од \(\Gamma/\mathsf P\) на все елементы од \(\mathfrak K_{r\Gamma}\) дефинујемо тако, же за всакы \(\mathfrak K_r\)-елемент \(z\) положимо
\[
S(z)=z
\]
то идеалы \(\mathfrak l_i\) сут коњуговане, т.~ј. примєњење једного \(S\) на неко \(\mathfrak l_j\) даваје друго
\[
S(\mathfrak l_j)=\mathfrak l_i,
\]
зато же всака неразложива компонента јединице \(e_j\) мора прєходити в неку другу \(e_i\). (Види такоже стр.~11.)

Модул прєдставјења \(e_i\Gamma\) образује \(\mathfrak Z\) на подпоље \(Z_i\) од \(\Gamma\). \(e_i\) уж лежи в \(\mathfrak Z_{Z_i}\), зато же прєдставјење \(Z_i\) од \(\mathfrak Z\) уж јест содрживано в \(Z_i\), зато \(\mathfrak Z_{Z_i}\) мора одцєпити модул прєдставјења \(e_iZ_i\). Нехај \(\{Z_i,Z_k\}\) буде композитум од \(Z_i\) и \(Z_k\). \(\mathfrak K_{r\{Z_i,Z_k\}}\) одцєпја два лєва идеала
\[
\bar{\mathfrak l}_i=\mathfrak K_{r\{Z_i,Z_k\}}e_i,\qquad
\bar{\mathfrak l}_k=\mathfrak K_{r\{Z_i,Z_k\}}e_k
\]
кторе сут просте, зато же имаје мєсто \(\mathfrak l_i=\mathfrak K_{r\Gamma}\bar{\mathfrak l}_i\). зато јест
\[
\bar{\mathfrak l}_i=\sum_{\lambda=1}^m\{Z_i,Z_k\}c_{\lambda i}^{(ik)},\qquad
\bar{\mathfrak l}_k=\sum_{\lambda}\{Z_i,Z_k\}c_{\lambda k}^{(ik)}.
\]
\(c_{\lambda\rho}^{(ik)}\) јест систем матрицных јединиц од \(\mathfrak K_{r\{Z_i,Z_k\}}\).

Всакы изоморфизм \(a_i\mapsto a_k\) од \(\bar{\mathfrak l}_i\) на \(\bar{\mathfrak l}_k\) с \(\mathfrak K_{r\{Z_i,Z_k\}}\) како областју операторов опрєдєљен јест елементом \(p_{ik}\) с \(e_i\mapsto p_{ik}\), зато же тогда јест \(a_i=a_ie_i\mapsto a_ip_{ik}\), зато \(\bar{\mathfrak l}_i\cdot p_{ik}=\bar{\mathfrak l}_k\). Вси можне таки \(p_{ik}\) очевидно сут
\[
p_{ik}=c_{ik}^{(ik)}\gamma_{ik},\qquad \gamma_{ik}\ne0\quad\text{из}\quad \{Z_i,Z_k\}.
\]
Посрєдством \(a_i\mapsto a_ip_{ik}\) тогда јест заданы такоже операторны изоморфизм од \(\mathfrak l_i\) на \(\mathfrak l_k\). Нынє избирајемо установјене системы такых \(p_{ik}\), налагајучи јим условје коњугованости.

Дєјство једного \(S\) на индекс \(i\) нехај буде дефиновано тако
\[
S(i)=k,\qquad\text{ако}\qquad S(Z_i)=Z_k.
\]
За всакы пар индексов \((\nu,\mu)\) обстаја принајмење једен коњугованы пар \(S(\nu,\mu)\) формы \(S(\nu,\mu)=(1,k)\). Из всакој класы коњугованых паров индексов избираје се установјены пар \((1,k)\), а тогда за тој \((1,k)\):
\[
p_{1k}=c_{1k}^{(1k)}\gamma_{1k},\qquad \gamma_{1k}\ne0\quad\text{из}\quad \{Z_1,Z_k\},
\]
иначе се положи либовољно, а тогда, ако \(S(1,k)=(\nu,\mu)\),
\[
p_{\nu\mu}=S(p_{1k}).
\]
\(p_{\nu\mu}\) тогда наистино имаје форму
\[
p_{\nu\mu}=c_{\nu\mu}^{(\nu\mu)}\gamma_{\nu\mu},\qquad
\gamma_{\nu\mu}\ne0\quad\text{из}\quad \{Z_\nu,Z_\mu\}.
\]
Зато же дєјство од \(S\) на изоморфизм
\[
a_1\mapsto a_1p_{1k},\qquad\text{специално}\qquad e_1\mapsto p_{1k}
\quad\text{од }\bar{\mathfrak l}_1\text{ на }\bar{\mathfrak l}_k
\]
даваје изоморфизм
\[
a_\nu\mapsto a_\nu p_{\nu\mu},\qquad\text{специално}\qquad e_\nu\mapsto p_{\nu\mu}
\quad\text{од }\bar{\mathfrak l}_\nu\text{ на }\bar{\mathfrak l}_\mu.
\]

Ако под \(c_{\lambda\rho}\) разумєјемо систем матрицных јединиц од \(\mathfrak K_{r\Gamma}\), јасно јест, же \(p_{\nu\mu}\) можно такоже записати в формє
\[
p_{\nu\mu}=c_{\nu\mu}\bar\gamma_{\nu\mu}
\]
релације
\[
c_{ij}c_{1k}=0,\quad\text{ако}\quad j\ne1,
\qquad
c_{ij}c_{jk}=c_{ik}
\]
имајут зато како послєдицу релације
\[
\begin{aligned}
p_{ij}p_{1k}&=0,\quad &&\text{ако }j\ne1,\\
p_{ij}p_{jk}&=\alpha_{ik}^{(j)}p_{ik},\quad
&&\alpha_{ik}^{(j)}\text{ из }\Gamma\text{ (точнеје из }\{Z_i,Z_j,Z_k\}\text{)}
\end{aligned}
\]
Те \(n^3\) величины \(\alpha\) творет \emph{факторны систем од \(\mathfrak K\) при избрању од \(Z\)}, кторы приналежи к систему \(p_{ik}\) ,,\emph{псеудоматрицных јединиц}``.

\(\alpha\) сут коњуговане: Из \(S(i,k,j)=(\nu,\mu,\tau)\) слєдује
\[
S\bigl(\alpha_{ik}^{(j)}\bigr)=\alpha_{\nu\mu}^{\tau}.
\]

Всакы систем \(p_{ik}^*\) псеудоматрицных јединиц настава из некого система \(p_{ik}\) тако, же почетне \(p_{ik}\) множимо либовољными \(\delta_{ik}\ne0\) из \(\{Z_i,Z_k\}\):
\[
p_{1k}^*=p_{1k}\delta_{1k}
\]
и тогда поновно, ако \(S(1,k)=(\nu,\mu)\), положимо
\[
p_{\nu\mu}^*=S(p_{1k}^*)
\]
зато јест
\[
p_{\nu\mu}^*=p_{\nu\mu}\delta_{\nu\mu}
\]
с \(\delta_{\nu\mu}\) из \(\{Z_\nu,Z_\mu\}\); \(\delta_{\nu\mu}\) сут коњуговане, из \(S(i,k)=(\nu,\mu)\) слєдује \(S(\delta_{ik})=\delta_{\nu\mu}\). Ако \(\alpha^*\) јест факторны систем приналежечи к \(p_{ik}^*\), то јест
\[
\alpha_{ik}^{(j)*}=\frac{\delta_{ij}\delta_{jk}}{\delta_{ik}}\alpha_{ik}^{(j)}.
\]

Факторне системы \(\alpha^*\) и \(\alpha\) называјут се \emph{асоциоваными} (так само \(p_{ik}^*\) и \(p_{ik}\)); вси факторне системы од \(\mathfrak K\) при избрању од \(Z\) творет класу \(\{\alpha\}\) асоциованых факторных системов.

Класа \(\{\alpha\}\) не зависи од тога, кторо прєдставјење \(\mathfrak Z\) од \(Z\) в \(\mathfrak K_r\) изберемо за основу јеј дефиниције. Зато же прєход од једного \(\mathfrak Z\) к другому можно учинити трансформацијоју елементом \(\lambda\) из \(\mathfrak K_r\), кторы јест зато \(\mathfrak G\)-инвариантны; зато из система \(p_{ik}\), приналежечего к \(\mathfrak Z\), настава систем
\[
p'_{ik}=\lambda^{-1}p_{ik}\lambda,
\]
кторы приналежи к \(\mathfrak Z'\); факторне системы од \(p_{ik}\) и \(p'_{ik}\) але сут идентичне.

Факторны систем не јест ничто друго неж таблица множења од \(\mathfrak K_{n\Gamma}\), прилагоджена особливым алгебраичным својствам од \(\mathfrak K\).
% END BOOK_S23

% BEGIN BOOK_S24
\subsection*{§ 24. Множење факторных системов}\label{multiplikation-von-faktorensystemen}

Вмєсто о разпадном пољу некомутативного тєла \(\mathfrak K\) говоримо такоже о разпадном пољу од \(\{\mathfrak K\}\); то имаје смысл, зато же сполу с \(\mathfrak K_Z\) такоже всако прстен \(\mathfrak K_{rZ}\) разпада се на абсолутно иредуцибилне компоненты, и обратно.

Всака два тєла \(\mathfrak K\), \(\mathfrak S\) с центром \(\mathsf P\) имајут обча разпадна поља, напр. композитум разпадного поља од \(\mathfrak K\) и разпадного поља од \(\mathfrak S\).

\textbf{Теорем 1.} \emph{Класы \(\{\mathfrak K\}\), кторе имајут поље \(Z\) како разпадно поље, творет подгрупу \(\mathscr K_Z\) групы \(\mathscr K\) всих \(\{\mathfrak K\}\).}

\emph{Доказ.} Ако \(\mathfrak K\) и \(\mathfrak S\) имајут \(Z\) како разпадно поље, то \(\mathfrak K_Z\) и \(\mathfrak S_Z\), зато такоже \((\mathfrak K\times\mathfrak S)_Z\), разпадајут се на абсолутно иредуцибилне компоненты, т.~ј. \(Z\) јест разпадно поље од \(\{\mathfrak K\times\mathfrak S\}\) (теорем 4 на стр.~23).

Реципрочно изоморфне тєла имајут вса разпадна поља обче.

Помочно разгледање.
\[
\mathfrak K_r=\mathfrak L_1+\cdots+\mathfrak L_r,\qquad
\mathsf P_m=\mathfrak z_1+\cdots+\mathfrak z_m
\]
нехај будут разклады од \(\mathfrak K_r\), \(\mathsf P_m\) на просте лєве идеалы; тогда
\[
\mathfrak K_{rm}=\mathfrak K_r\times\mathsf P_m
=\mathfrak L_1\mathfrak z_1+\cdots+\mathfrak L_i\mathfrak z_j+\cdots+\mathfrak L_r\mathfrak z_m
\]
јест разклад од \(\mathfrak K_{rm}\) на \(rm\) ненуловых, зато простых лєвых идеалов. Зато всакы \(r\)-членны лєвы идеал \(A\) од \(\mathfrak K_{rm}\) јест свезаны с \(r\)-членным лєвым идеалом
\[
\mathfrak L_1\mathfrak z_1+\cdots+\mathfrak L_r\mathfrak z_1=\mathfrak K_r\mathfrak z_1
\]
посрєдством внутрного автоморфизма од \(\mathfrak K_r\)
\[
A=\tau^{-1}\mathfrak K_r\tau\cdot\tau^{-1}\mathfrak z_1\tau
=\mathfrak K'_r\cdot\mathfrak z'_1.
\]

Ако \(\mathfrak K_{r\Gamma}=\mathfrak l_1+\cdots+\mathfrak l_n\) јест разклад од \(\mathfrak K_{r\Gamma}\) на коњуговане лєве идеалы, приналежечи к некому \(\mathfrak Z\) в \(\mathfrak K_r\), \(\mathfrak l_i=\mathfrak K_{r\Gamma}e_i\), то заради
\[
S(\tau^{-1}\mathfrak l_i\tau\cdot\mathfrak z'_1)
=S(\tau^{-1}\mathfrak l_i\tau)\,\mathfrak z'_1
=\tau^{-1}(S\mathfrak l_i)\tau\cdot\mathfrak z'_1
=\tau^{-1}\mathfrak l_k\tau\cdot\mathfrak z'_1
\]
в
\[
\tag{1}
A_\Gamma=\mathfrak l'_1\mathfrak z'_1+\cdots+\mathfrak l'_n\mathfrak z'_1
\]
имајемо разклад од \(A_\Gamma\) на коњуговане лєве идеалы. Јест
\[
\mathfrak l'_i\mathfrak z'_1=\mathfrak K_{rm\Gamma}\cdot e'_iE',
\]
ако положимо \(\tau^{-1}e_i\tau=e'_i\), а \(E'\) означаје јединицу од \(\mathfrak z'_1\).

Всакы другы разклад
\[
\tag{2}
A_\Gamma=\mathfrak q_1+\cdots+\mathfrak q_n,\qquad
S\mathfrak q_i=\mathfrak q_k\quad\text{ако}\quad
S\mathfrak l'_i\mathfrak z'_1=\mathfrak l'_k\mathfrak z'_1
\]
на коњуговане лєве идеалы настава из (1) трансформацијоју с \(\mathfrak K_{rm}\)-елементом \(\lambda\)
\[
\mathfrak q_i=\lambda^{-1}\mathfrak l'_i\mathfrak z'_1\lambda,
\qquad
\text{то јест, ако \(e_i^*\) јест јединица од \(\mathfrak q_i\),}\quad
e_i^*=\lambda^{-1}e'_iE'\lambda.
\]

Доказ јест непосрєдња послєдица слєдујучеј помочној лемы:

\textbf{Помочна лема.} \emph{Нехај \(A\) буде лєвы идеал из \(\mathfrak K_n\), и нехај}
\[
A=\mathfrak l_1+\cdots+\mathfrak l_n
=\mathfrak K_{n\Gamma}e_1+\cdots+\mathfrak K_{n\Gamma}e_n
\]
\emph{буде -- по прєдположењу обстајајучи -- разклад на коњуговане лєве идеалы; при том можно -- без ограничења обчости -- прєдпокладати, же \(\mathfrak l_i\) сут коњуговане.}

\emph{Нехај \(\mathfrak z\) буде операторно изоморфны с \(A\); и нехај будут изполњене теже прєдположења:}
\[
\mathfrak z=\mathfrak q_1+\cdots+\mathfrak q_n
=\mathfrak K_{n\Gamma}e_1^*+\cdots+\mathfrak K_{n\Gamma}e_n^*
\]
\emph{(зато \(\mathfrak q_i\) и \(e_i^*\) сут коњуговане). Тогда обстаја регуларны елемент \(\lambda\) из \(\mathfrak K_n\), тако же \(\mathfrak q_i=\lambda^{-1}\mathfrak l_i\lambda\).}

\emph{Доказ. 1. Факт, же \(e_i\) можно прєдпокладати коњуговаными, слєдује тако:} Ако \(e_1\) јест либовољны творечи идемпотент од \(\mathfrak l_1\), то \(\mathfrak l_1=\mathfrak K_{n\Gamma}e_1\), примєњењем автоморфизма на \(\Gamma\) слєдује \(\mathfrak l_i=\mathfrak K_{n\Gamma}e_i\), кде \(e_i\) јест коњугованы с \(e_1\) и поновно творечи идемпотент.

2. \emph{Всако \(\mathfrak l\) јест изоморфно с всакым \(\mathfrak q\);} зато же \(\mathfrak l\), како коњуговане, все имајут једнаку (абсолутну) должину, числа од \(\mathfrak l\) и \(\mathfrak q\) сут једнака, а \(\mathfrak z\) и \(A\) сут по прєдположењу изоморфне.

3. Приреджењем
\[
\tag{3}
e_1\mapsto s_1=e_1s_1;\qquad
e_1^*\mapsto t_1=e_1^*t_1;\qquad
s_1t_1=e_1;\qquad
t_1s_1=e_1^*
\]
нехај будут установјене изоморфизм меджу \(\mathfrak l_1\) и \(\mathfrak q_1\) и његов обрат. Примєњењем всих автоморфизмов (3) прєходи в
\[
e_i\mapsto s_i=e_is_i;\qquad
e_i^*\mapsto t_i=e_i^*t_i;\qquad
s_it_i=e_i;\qquad
t_is_i=e_i^*,
\]
чим сут зато установјене изоморфизмы \(\mathfrak q_i=\mathfrak l_is_i\) и \(\mathfrak l_i=\mathfrak q_it_i\).

4. Нынє положимо:
\[
s=s_1+\cdots+s_n;\qquad t=t_1+\cdots+t_n;
\]
тогда добыва се:
\[
\mathfrak q_i=\mathfrak l_is;\qquad
\mathfrak l_i=\mathfrak q_it;\qquad
st=e_1+\cdots+e_n=\xi;\qquad
ts=e_1^*+\cdots+e_n^*=\xi^*,
\]
причем \(\xi\) и \(\xi^*\) лежет в \(\mathfrak K_n\); добыва се
\[
A\xi=A,\qquad \mathfrak z\xi^*=\mathfrak z\quad
(\text{заради }\mathfrak l_i\xi=\mathfrak l_i\mathfrak l_i=\mathfrak l_i),
\]
\[
\xi^2=\xi,\qquad \xi^{*2}=\xi^*.
\]
Ако нынє положимо: \(e=\overline E+E=\overline E^*+E^*\) и с тым:
\[
\mathfrak K_n=\overline A+A
=\mathfrak K_n\overline\xi+\mathfrak K_n\xi
=\overline{\mathfrak z}+\mathfrak z
=\mathfrak K_n\overline\xi^*+\mathfrak K_n\xi^*,
\]
то \(\overline A\) и \(\overline{\mathfrak z}\) сут такоже изоморфне; ако зато
\[
\overline E\mapsto\overline s=\overline E\,\overline s
\qquad\text{и}\qquad
\overline E^*\mapsto\overline t=\overline E^*\overline t
\]
сут обратне изоморфизмы, и ако положимо:
\[
\lambda=s+\overline s,\qquad \mu=t+\overline t,
\]
добыва се:
\[
\lambda\mu=\mu\lambda=e,\qquad
\mathfrak l_i\lambda=\mathfrak q_i,\qquad
\lambda^{-1}\mathfrak q_i=\mathfrak q_i
\quad
(\lambda^{-1}\mathfrak q_i\subseteq\mathfrak q_i
=\lambda\cdot\lambda^{-1}\mathfrak q_i\subseteq\lambda^{-1}\mathfrak q_i)
\]
и зато
\[
\lambda^{-1}\mathfrak l_i\lambda=\mathfrak q_i,\qquad \text{тым јест доказано.}
\]

К разкладу (1) можно поновно опрєдєлити елементы
\[
r_{ik}=c_{ik}^{\prime(ik)}\gamma_{ik},\qquad \gamma_{ik}\in\{Z_i,Z_k\},
\]
кторе посрєдством \(e'_iE'\mapsto r_{ik}\) давајут операторне изоморфизмы меджу идеалами \(\mathfrak l'_i\mathfrak z'_1\) и \(\mathfrak l'_k\mathfrak z'_1\), и такоже изполњајут условје коњугованости \(S(r_{ik})=r_{\nu\mu}\), ако \(S(i,k)=(\nu,\mu)\). Они посрєдством
\[
r_{ij}r_{jk}=\alpha_{ik}^{(j)}r_{ik}
\]
опрєдєљајут факторны систем \(\alpha\) од \(A_\Gamma\). Нынє јест важно, же за опрєдєљење факторных системов од \(A_\Gamma\) можно такоже употрєбити либовољны разклад (2). Зато же ако систем \(r_{ik}\), приналежечи к (1), замєнимо с \(r'_{ik}=\lambda^{-1}r_{ik}\lambda\), то \(e_i^*\mapsto r'_{ik}\) даваје операторны изоморфизм од \(\mathfrak q_i\) на \(\mathfrak q_k\); \(r'_{ik}\) изполњајут условје коњугованости, зато же \(\lambda\) јест \(\mathfrak q\)-инвариантны, а \(r'_{ik}\) имајут тојже факторны систем како \(r_{ik}\).

Ако изберемо систем псеудоматрицных јединиц \(p_{ik}\) од \(\mathfrak K\), то
\[
r_{ik}=\tau^{-1}p_{ik}\tau\cdot E'
\]
јест \(r_{ik}\)-систем од \(A_\Gamma\) с једнакым факторным системом. зато имајемо помочну лему:

\textbf{Теорем.} \emph{Факторне системы од \(A_\Gamma\), образоване с помочју либовољного разклада}
\[
A_\Gamma=\mathfrak q_1+\cdots+\mathfrak q_n
\]
\emph{од \(A_\Gamma\) на коњуговане лєве идеалы, сут једнака факторным системам од \(\mathfrak K_r\).}

Слєдујучи теорем показује, же ако \(\alpha_{ik}^{(j)}\), \(\beta_{ik}^{(j)}\) сут факторне системы од \(\{\mathfrak K\}\), \(\{\mathfrak S\}\) при избрању од \(Z\), то \(\varepsilon_{ik}^{(j)}=\alpha_{ik}^{(j)}\cdot\beta_{ik}^{(j)}\) јест факторны систем од \(\{\mathfrak K\times\mathfrak S\}\) при избрању од \(Z\). Зато имаје смысл дефинирати \emph{продукт \(\{\alpha\}\times\{\beta\}\) двох клас асоциованых факторных системов} \(\{\alpha\}\), \(\{\beta\}\) како класу \(\{\alpha_{ik}^{(j)}\beta_{ik}^{(j)}\}\).

\textbf{Теорем 2.} \emph{Нехај \(\{\mathfrak K\}\) при избрању од \(Z\) имаје класу \(\{\alpha\}\) асоциованых факторных системов. Класы \(\{\alpha\}\) всих \(\{\mathfrak K\}\) с разпадным пољем \(Z\) творет групу, ктора јест посрєдством \(\{\mathfrak K\}\mapsto\{\alpha\}\) изоморфно однесена на \(\mathscr K_Z\).}

\textbf{За доказ} трєба показати двоје:

1. Ако \(\{\mathfrak K\}\) имаје факторны систем \(\alpha_{ik}^{(j)}\), а \(\{\mathfrak S\}\) \(\beta_{ik}^{(j)}\), то \(\varepsilon_{ik}^{(j)}=\alpha_{ik}^{(j)}\beta_{ik}^{(j)}\) јест факторны систем од \(\{\mathfrak K\times\mathfrak S\}\).

2. Ако \(\{\mathfrak K\}\) имаје факторны систем \(\alpha_{ik}^{(j)}=1\), то \(\mathfrak K=\mathsf P\).

1. Нехај ступењи иредуцибилных прєдставјениј од \(Z\) в \(\mathfrak K\), \(\mathfrak S\) будут \(r,s\)
\[
\mathfrak K_{r\Gamma}=\mathfrak l_1+\cdots+\mathfrak l_n,\qquad
\mathfrak S_{s\Gamma}=\mathfrak m_1+\cdots+\mathfrak m_n
\]
разклады на коњуговане идеалы, причем \(\mathfrak l_i\) и \(\mathfrak m_i\) трєба да будут взајемно одговарјајучи идеалы (прве изведене из одговарјајучих компонентов од \(\mathfrak Z_r\), а друге из \(\mathfrak Z'_r\)). Тогда јест
\[
(\mathfrak K_r\times\mathfrak S_s)_\Gamma
=\mathfrak l_1\mathfrak m_1+\cdots+\mathfrak l_i\mathfrak m_j+\cdots+\mathfrak l_n\mathfrak m_n
\]

разклад на \(n^2\) ненуловых, зато простых идеалов. Зато јест
\[
A_1=\mathfrak l_1\mathfrak m_1+\mathfrak l_2\mathfrak m_2+\cdots+\mathfrak l_n\mathfrak m_n
\]
идеал абсолутној должины \(n\). \(A_1\) јест инвариантны при всих \(S\), зато же \(\mathfrak m_i\) трансформујут се тако како \(\mathfrak l_i\); зато обстаја \(\mathfrak K_r\times\mathfrak S_s\)-идеал \(A\) с
\[
A_1=A_\Gamma
\]
(помочна лема, стр.~18).

Ако \(u\) јест ступењ иредуцибилного прєдставјења од \(\mathfrak Z\) в тєлу автоморфизмов \(\mathfrak M\) простых идеалов од \(\mathfrak K\times\mathfrak S\), то \(A\) разпада се на \(u\) простых идеалов (зато же абсолутна должина јест \(n\)). Наша помочна лема даваје зато, же факторне системы од \(A\) сут такоже факторне системы од \(\mathfrak M\).

Але разклад од \(A_\Gamma\) на коњуговане идеалы, употрєбјены за дефиницију, јест:
\[
\tag{3}
A_\Gamma=\mathfrak l_1\mathfrak m_1+\mathfrak l_2\mathfrak m_2+\cdots+\mathfrak l_n\mathfrak m_n.
\]
Ако \(\alpha_{ik}^{(j)}\), \(\beta_{ik}^{(j)}\) приналежет к псеудоматрицным јединицам \(p_{ik}\), \(q_{ik}\) од \(\mathfrak K\), \(\mathfrak S\), то очевидно \(r_{ik}=p_{ik}q_{ik}\) јест систем псеудоматрицных јединиц од \(A_\Gamma\) по одношењу к разкладу (3); к тому приналежечи факторны систем \(\varepsilon_{ik}^{(j)}=\alpha_{ik}^{(j)}\beta_{ik}^{(j)}\) јест тогда такоже једен из факторных системов од \(\mathfrak M\).

2. Нехај \(\{\mathfrak K\}\) имаје факторны систем \(\alpha_{ik}^{(j)}=1\) с \(p_{ik}\) како приналежечими псеудоматрицными јединицами и разкладом
\[
\mathfrak K_{r\Gamma}=\mathfrak l_1+\cdots+\mathfrak l_n
\]
на коњуговане идеалы. Трєба доказати \(\mathfrak K=\mathsf P\), или \(\mathfrak K_r=\mathsf P_r\), или обстајање \(n\)-членного лєвого идеала \(\mathfrak l\) од \(\mathfrak K_r\) -- зато же просты лєвы идеал од \(\mathfrak K_r\) имаје ранг \(t^2r\), зато \(n=t^2rh\) с цєлым \(h\), или мора быти \(t=1\).

Идеал од \(\mathfrak K_{r\Gamma}\)
\[
\Omega=\mathfrak l_1(p_{11}+\cdots+p_{1n})
\]
не јест нула -- \(e_1\sum_i p_{1i}\ne0\) -- и јест просты, зато ранга \(n\). \(\Omega\) јест \(\mathfrak G\)-инвариантны,
\[
S\Omega
=\mathfrak l_i\sum_\nu S p_{1\nu}
=\mathfrak l_i\sum_\mu p_{i\mu}
=\mathfrak l_i\sum_\mu p_{i1}p_{1\mu}
\]
зато же \(\alpha_{ik}^{(j)}=1\). Зато
\[
S\Omega=\mathfrak l_ip_{i1}\sum_\mu p_{1\mu}
=\mathfrak l_1\sum_\mu p_{1\mu}
=\Omega.
\]
зато \(\Omega=\mathfrak l_\Gamma\), \emph{кде \(\mathfrak l\) означаје \(\mathfrak K_r\)-идеал \(n\)-того ранга.}

\textbf{Теорем 3.} \emph{Всакы елемент \(\{\mathfrak K\}\) групы \(\mathscr K\) имаје конечны ред \(\lambda\), кторы јест дєлитељ абсолутного числа компонентов \(t\) од \(\mathfrak K\).}

\emph{Првы доказ (Brauer).} Изберемо два разклады од \(\mathfrak K_{r\Gamma}\) на коњуговане лєве идеалы
\[
\mathfrak K_{r\Gamma}=\mathfrak l_1+\cdots+\mathfrak l_n=\mathfrak z_1+\cdots+\mathfrak z_n,\qquad
\mathfrak l_i=\mathfrak K_{r\Gamma}e_i,\quad
\mathfrak z_i=\mathfrak K_{r\Gamma}e_i^*.
\]
Нехај посрєдством \(a_i\mapsto a_ir_{ik}\), специално \(e_i\mapsto r_{ik}\), буде заданы операторны изоморфизм од \(\mathfrak l_i\) на \(\mathfrak z_k\). При његовом обрату \(e_k^*\) прєходи в неко \(s_{ki}\) из \(\mathfrak l_i\), \(s_{ki}\cdot r_{ik}=e_k^*\). Ако изберемо \(r_{ik}\) коњуговаными
\[
S(r_{ik})=r_{\nu\mu}\quad\text{ако}\quad S(i,k)=(\nu,\mu),
\]
то \(s_{ik}\) такоже сут коњуговане:
\[
S(s_{ki})S(r_{ik})=S(s_{ki}r_{ik})=S(e_k^*)=e_\mu^*
=s_{\mu\nu}r_{\nu\mu}=s_{\mu\nu}S(r_{ik}),
\]
зато \(S(s_{ki})=s_{\mu\nu}\). Зато же посрєдством \(e_i\mapsto e_ir_{ij}s_{jk}\) јест заданы операторны изоморфизм од \(\mathfrak l_i\) на \(\mathfrak l_k\), мора быти
\[
r_{ij}s_{jk}=p_{ik}\xi_{ik}^{(j)},\qquad
\xi_{ik}^{(j)}\ne0\in\Gamma
\]

Зато же \(r\)-елементи, \(s\)-елементи и \(p\)-елементи сут коњуговане, \(\xi\) такоже сут коњуговане: Из \(S(i,j,k)=(\nu,\tau,\mu)\) слєдује
\[
S(\xi_{ik}^{(j)})=\xi_{\nu\mu}^{\tau}.
\]

Из
\[
r_{i\lambda}s_{\lambda j}r_{j\lambda}s_{\lambda k}
=r_{i\lambda}e_\lambda^*s_{\lambda k}
=r_{i\lambda}s_{\lambda k}
\]
слєдује
\[
p_{ij}\xi_{ij}^{(\lambda)}\cdot p_{jk}\xi_{jk}^{(\lambda)}
=p_{ik}\xi_{ik}^{(\lambda)}
\]
или, зато же \(p_{ij}p_{jk}=p_{ik}\alpha_{ik}^{(j)}\):
\[
\alpha_{ik}^{(j)}
=\frac{\xi_{ik}^{(\lambda)}}{\xi_{ij}^{(\lambda)}\xi_{jk}^{(\lambda)}}.
\]
Ако положимо \(\delta_{\nu\mu}=\prod_{\lambda}\xi_{\nu\mu}^{(\lambda)}\), то \(\delta_{\nu\mu}\) јест инвариантны при всих \(S\), кторе оставјајут \(\nu\) и \(\mu\) неподвижными, лежи зато в \(\{Z_\nu,Z_\mu\}\); надто \(\delta_{\nu\mu}\) сут коњуговане, зато же \(\xi_{\nu\mu}^{(j)}\) сут коњуговане. Зато же нынє
\[
\alpha_{ik}^{(j)n}
=\frac{\prod_{\lambda}\xi_{ik}^{(\lambda)}}
       {\prod_{\lambda}\xi_{ij}^{(\lambda)}\prod_{\lambda}\xi_{jk}^{(\lambda)}}
=\frac{\delta_{ik}}{\delta_{ij}\delta_{jk}}
\]
из формулы за асоциоване факторне системы (стр.~32) слєдује, же \(\alpha_{ik}^{(j)n}\) јест асоциованы с системом \(\beta_{ik}^{(j)}=1\); зато, зато же ред од \(\{\mathfrak K\}\) јест равны реду од \(\{\alpha\}\), ред од \(\{\mathfrak K\}\) јест дєлитељ од \(n\). Ако особливо изберемо \(\mathfrak Z\) како максимално комутативно подпоље од \(\mathfrak K\), то јест \(n=t\), то \(\lambda\) јест дєлитељ од \(t\).

\emph{Другы доказ (Schur).} Тој доказ основан јест на прєдставјењу \(p_{ik}\) посрєдством регуларных матриц реда \(n\), \(P_{ik}\), в \(\Gamma\), кторе изполњајут једнаке условја коњугованости како \(p_{ik}\). \(P_{ik}\) лежи в \(\{Z_i,Z_k\}\); ако \(p_{ik}\) приналежи к \(P_{ik}\), то \(S(p_{ik})\) приналежи к \(S(P_{ik})\). То прєдставјење јест хомоморфно; т.~ј. ако \(P_{ik}\) приналежи к \(p_{ik}\), а \(P_{kj}\) к \(p_{kj}\), то \(P_{ik}P_{kj}\) приналежи к \(p_{ik}p_{kj}\), или
\[
P_{ik}P_{kj}=\alpha_{ij}^{(k)}P_{ij}
\]
\((P_{ik_1}P_{k_2j}=0\) за \(k_1\neq k_2\) не може быти правдиво). Ако имајемо то прєдставјење од \(p_{ik}\), закључујемо тако: Прєход к детерминантам даваје:
\[
|P_{ik}|\cdot |P_{kj}|=\alpha_{ij}^{(k)n}\cdot |P_{ij}|,
\]
тым самым, зато же матрице \(P_{ik}\) сут регуларне, \(|P_{\nu\mu}|\ne0\),
\[
\alpha_{ij}^{(k)n}=\frac{\delta_{ik}\delta_{kj}}{\delta_{ik}},
\]
кде јест положено \(|P_{\nu\mu}|=\delta_{\nu\mu}\). И потом даље како в Брауеровом доказу.

\emph{Прєдставјење од \(p_{ik}\) посрєдством матриц \(P_{ik}\):}

Идеалы \(\mathfrak l_i\) разклада
\[
\mathfrak K_{r\Gamma}=\mathfrak l_1+\cdots+\mathfrak l_n
\]
сут реципрочне модулы прєдставјења од \(\mathfrak K_r\) в \(\Gamma\) (теорем 2 на стр.~32). Ако за основу изберемо установјену базу \(k_1,\ldots,k_n\) од \(\mathfrak K_r\) по одношењу к \(\mathfrak Z\), то јест
\[
\mathfrak K_{r\Gamma}=k_1\mathfrak Z_r+\cdots+k_n\mathfrak Z_r,
\]
зато
\[
\mathfrak l_i=\mathfrak K_{r\Gamma}e_i
=k_1\mathfrak Z_r e_i+\cdots+k_n\mathfrak Z_r e_i
=k_1e_i\Gamma+\cdots+k_ne_i\Gamma.
\]
Из тога слєдује
\[
\mathfrak l_k=\mathfrak l_i p_{ik}
=k_1p_{ik}\Gamma+\cdots+k_np_{ik}\Gamma.
\]
Двє базы \((k_1e_i,\ldots,k_ne_i)\) и \((k_1p_{ik},\ldots,k_np_{ik})\) од \(\mathfrak l_k\) можно прєвести једну в другу множењем регуларној матрицоју \(P_{ik}\):
\[
P_{ik}\cdot
\begin{pmatrix}
k_1e_i\\ \vdots\\ k_ne_i
\end{pmatrix}
=
\begin{pmatrix}
k_1p_{ik}\\ \vdots\\ k_np_{ik}
\end{pmatrix},
\qquad
\text{кратко}\quad
P_{ik}(k_\varrho e_i)=(k_\varrho p_{ik}).
\]
Из тога слєдује коњугованост од \(P_{ik}\):
\[
S(P_{ik})(k_\varrho S(e_k))=(k_\varrho S(p_{ik}))
\quad\text{или}\quad
S(P_{ik})(k_\varrho e_\mu)=(k_\varrho p_{\nu\mu})
=P_{\nu\mu}(k_\varrho e_\mu),
\]
т.~ј.
\[
S(P_{ik})=P_{\nu\mu}\quad\text{ако}\quad S(i,k)=(\nu,\mu).
\]
Из тога поновно слєдује, же \(P_{ik}\) лежи в \(\{Z_i,Z_k\}\) -- јменно в инвариантном пољу тых \(S\), кторе оставјајут \(i\) и \(k\) неподвижными. Хомоморфност такоже јест легко видєти:
\[
\begin{split}
P_{ik}P_{kj}(k_\varrho e_j)
&=P_{ik}(k_\varrho p_{kj})
=(k_\varrho p_{ik}p_{kj})\\
&=\alpha_{ij}^{(k)}(k_\varrho p_{ij})
=\alpha_{ij}P_{ij}(k_\varrho e_j),
\quad\text{зато}\\
&\cdot\, P_{ik}P_{kj}=\alpha_{ij}^{(k)}P_{ij},
\quad\text{тым јест доказано.}
\end{split}
\]

Теорем 3 можно посилити:

\textbf{Теорем 4.} \emph{Ако \(p\) јест просты фактор абсолутного числа компонентов \(t\) од \(\mathfrak K\), то \(p\) входи такоже в ред \(\lambda\) од \(\{\mathfrak K\}\). (Brauer)}

\emph{Доказ.} Нехај \(p^\sigma\) буде највышша потенција од \(p\), ктора се налази в ступењу \(n\) разпадного поља \(Z\), \(\mathfrak H\) подгрупа реда \(p^\sigma\) од \(\mathfrak G\) (она обстаја по \emph{теорему од Сылова}, Спеисер 1. изд. § 19, стр.~42, теорем 43), а \(T\) инвариантно поље од \(\mathfrak H\).

Тогда \([\Gamma:T]=p^\sigma\) и \([T:\mathsf P]=n/p^\sigma\) не јест дєливы посрєдством \(p\), а тым выше не посрєдством \(t\). \(T\) зато не јест разпадно поље; тєло автоморфизмов \(\mathfrak S\) од \(\mathfrak K_T=\mathfrak S_m\) јест различно од \(\mathsf P\). Але зато же \(\mathfrak S\) имаје разпадно поље \(\Gamma\), чији ступењ над центром \(T\) од \(\mathfrak S\) јест равны \(p^\sigma\), абсолутно число компонентов од \(\mathfrak S\) јест потенција \(p^\beta>1\), зато по теорему 3 ред од \(\{\mathfrak S\}\) јест потенција \(p^\alpha>1\) од \(p\). Зато же из
\[
\{\mathfrak K\}^{\lambda}=\{\mathsf P\},
\qquad
\{\mathfrak K_T\}^{\lambda}=\{T\}
\]
слєдује, то \(\lambda\equiv0\pmod {p^\alpha}\), тым јест доказано. (Brauer дал примєры, же не мора быти \(\lambda=t\)\srcfn{1}{\foreign{Brauer: Untersuchungen über die arithmetischen Eigenschaften von Gruppen linearer Substitutionen II. Math. Zeitschr. 31, str.~733 § 5, (1930).}}).

\textbf{Теорем 5.} \emph{Нехај \(t\), разложено на просте факторы, буде \(t=\prod_i p_i^{\varrho_i}\), \(\varrho_i\ge1\), \(p_i\ne p_j\), ако \(i\ne j\). Тогда јест \(\mathfrak K=\mathfrak K^{(1)}\times\mathfrak K^{(2)}\times\cdots\), кде \(\mathfrak K^{(i)}\) јест тєло с абсолутным числом компонентов \(p_i^{\varrho_i}\).}

\emph{Доказ.} По теоремах 3 и 4 ред \(\lambda\) од \(\{\mathfrak K\}\) јест
\[
\lambda=\prod_i p_i^{\alpha_i},\qquad 1\le\alpha_i\le\varrho_i.
\]
Посрєдством \(\{\mathfrak K\}\) створјена циклична подгрупа \(\mathfrak Z\) од \(\mathscr K\) јест прємы продукт подгруп \(\mathfrak Z_i\) редов \(p_i^{\alpha_i}\), то јест
\[
\{\mathfrak K\}=\prod_i\{\mathfrak K^{(i)}\},
\]

из чего добыва се, же всако разпадно поље \(Z\) од \(\mathfrak K\) јест такоже разпадно поље од \(\mathfrak K^{(i)}\), зато абсолутно число компонентов \(p_i^{\sigma_i}\) од \(\mathfrak K^{(i)}\) дєли \(t\), \(\sigma_i\le\varrho_i\). Нехај \(\prod_i \mathfrak K^{(i)}=\mathfrak K_m\), то јест
\[
\prod_i p_i^{\varrho_i}=\prod_i p_i^{\sigma_i}\cdot m,
\qquad
\sigma_i\ge\varrho_i,
\]
зато
\[
\sigma_i=\varrho_i,\qquad m=1,\qquad
\mathfrak K=\prod_i \mathfrak K^{(i)},
\quad\text{тым јест доказано.}
\]
% END BOOK_S24

% BEGIN BOOK_S25
\subsection*{§ 25. Нормално прєдставјење од \(\mathfrak K_r\) с Галоисовым максималным комутативным подпољем\protect\footnotemark}
\label{normaldarstellung-von-r_r-mit-galoisschem-maximalem-kommutativem-teilkoerper-1}
\footnotetext{Простєјше основање од § 27 даље.}

Нехај максимално комутативно подпоље \(\mathfrak Z\) од \(\mathfrak K_r\) буде Галоисово. \(n\) автоморфизмов \(H_i\) од \(\mathfrak Z/\mathsf P\) сут по теорему 7 на стр.~25 внутрне автоморфизмы од \(\mathfrak K_r\), т.~ј. обстајајут елементи \(u_i\) с
\[
H_i z=u_i^{-1}zu_i\quad\text{за все }z\text{ из }\mathfrak Z.
\]
Всакы \(u_i\) јест опрєдєљен до ненулового фактора \(y\) из \(\mathfrak Z\). Зато же ако \(u^{-1}zu=u'^{-1}zu'\) за все \(z\) из \(\mathfrak Z\), то слєдује \((u^{-1}u')^{-1}z(u^{-1}u')=z\) за все \(z\in\mathfrak Z\). Приложење од \(u^{-1}u'\) к \(\mathfrak Z\) даваје зато комутативно подпоље од \(\mathfrak K_r\), кторо заради максималности од \(\mathfrak Z\) сбєгаје се с \(\mathfrak Z\). зато \(u^{-1}u'\) лежи в \(\mathfrak Z\).

Доказујемо слєдујучи теорем:
\[
\mathfrak K_r=u_1\mathfrak Z+\cdots+u_n\mathfrak Z.
\]
За то очевидно јест потрєбна само линеарна независност од \(u_i\) по одношењу к \(\mathfrak Z\). Доказ се изводи непосрєдним задањем од \(u_i\). Нехај буде
\[
\mathfrak Z_Z=e_1Z+\cdots+e_nZ,\qquad
\mathfrak K_{rZ}=\mathfrak l_1+\cdots+\mathfrak l_n,\quad
\mathfrak l_i=\mathfrak K_{rZ}e_i
\]
и нехај \(p_{ik}\) буде систем псеудоматрицных јединиц по одношењу к тому разкладу. На стр.~8 означили јесмо с \(\Theta_i\) изоморфизм од \(\mathfrak Z\) на \(Z\), даны посрєдством модула прєдставјења \(e_iZ\), тако же
\[
S_i=\Theta_i\Theta_1^{-1}
\]
сут \(n\) автоморфизмов од \(Z\) (але само колико се они примєњујут на само \(Z\), не выше за ту всечасно прєдпокладано дєјство од \(S_i\) на цєло \(\mathfrak K_{rZ}\)!). Заради простоты хочемо вмєсто \(\Theta_i\) в \(S_i=\Theta_i\Theta_1^{-1}\) писати \(\Theta_{S_i}\), а тогда такоже при \(p_{ik}\) замєнити индексы \(i,k\) приналежечими автоморфизмами \(S\)
\[
p_{ik}=p_{S_iS_k}.
\]

\(n\) автоморфизмов од \(\mathfrak Z\) сут (\(E\) јест идентична субституција)
\[
H=\Theta_E^{-1}\Theta.
\]
Нумеровање од \(H\) нынє избирајемо тако:
\[
H_S=\Theta_E^{-1}\Theta_{S^{-1}}.
\]

При тых условјах можно положити

\[
u_T=\sum_S p_{S,ST}=\sum_S S(p_{E,T})
\]



1. \(u_T^{-1}\) обстаја:
\[
u_T^{-1}=\sum_S p_{ST,S}
\frac{1}{\alpha_{S,S}^{ST}\overline{\gamma}_{S,ST}\overline{\gamma}_{ST,S}}
\]

(Значење од \(\bar{\gamma}\) види на стр.~32).

2. \(u_T\) јест елемент од \(\mathfrak K_r\), зато же јест \(\mathfrak G\)-инвариантны:
\[
R(u_T)=\sum_S RS(p_{E,T})=\sum_{S'} S'(p_{E,T})=u_T.
\]

3. \(u_T\) стварјаје \(H_T\), т.~ј. ако \(z\in\mathfrak Z\), то \(u_T^{-1}zu_T=H_T(z)\), или
\[
z u_T=u_T H_T(z).
\]

Имамо

\[
z=\sum_S\Theta_S(z)e_S,\qquad
u_T=\sum_S S(e_Ep_{E,T})=\sum_S e_S S(p_{E,T})
\]

зато

\[
zu_T=\sum_{S,S'}\Theta_S(z)e_Se_{S'}S(p_{E,T})
=\sum_S\Theta_S(z)S(p_{E,T}).
\]

С другој страны имајемо

\[
u_T=\sum_S S(p_{E,T}e_T)=\sum_S S(p_{E,T})e_{ST}
\]

и имаје мєсто слєдујуче једначење за дєјство на елементы од \(\mathfrak Z\)

\[
\Theta_S H_T=\Theta_S\Theta_E^{-1}\Theta_{T^{-1}}
=S\Theta_{T^{-1}}=ST^{-1}\Theta_E=\Theta_{ST^{-1}}
\]

зато

\[
H_T(z)=\sum_S e_S\Theta_S H_T(z)=\sum_S e_S\Theta_{ST^{-1}}(z)
\]

т.~ј.

\[
\begin{aligned}
u_TH_T(z)
&=\sum_{S,S'} S(p_{E,T})e_{ST}\cdot e_{S'}\Theta_{S'T^{-1}}(z)\\
&=\sum_S S(p_{E,T})\Theta_S(z)
\end{aligned}
\]

зато

\[
z\cdot u_T=u_T\cdot H_T(z).
\]

4. \(u_S\) сут линеарно независне по одношењу к \(\mathfrak Z\).

Из
\[
\sum_T z_Tu_T=0
\]
заради \(z=\sum_S\Theta_S(z)e_S\) слєдује

\[
\sum_{T,S,L}\Theta_S(z_T)e_Sp_{L,LT}
=\sum_{S,T}\Theta_S(z_T)p_{S,ST}\quad\text{зато}\quad
\Theta_S(z_T)=0,
\]

зато: \(z_T=0\).

Од нынє вмєсто \(H_S(z)\) пишемо \(z^S\), ,,симболична потенција``. За \(u_S\) морајут быти изполњене релације слєдујучего вида

\[
u_Ru_T=u_{RT}a_{R,T},\qquad a_{R,T}\ne0\quad\text{из }\mathfrak Z.
\]

Зато же

\[
\begin{aligned}
u_Ru_T
&=\sum_{S,S'} S(p_{E,R})S'(p_{E,T})
=\sum_S p_{S,SR}p_{SR,SRT}\\
&=\sum_S p_{S,SRT}\alpha_{S,SRT}^{(SR)}
=\sum_S p_{S,SRT}\alpha_{S,SRT}^{(SR)}e_{SRT}
\end{aligned}
\]

и

\[a_{R,\,T} = \sum\limits_{S} e_{SRT}\, \Theta_{SRT} \left(a_{R,\,T}\right),\]

зато

\[u_{RT} a_{R,T} = \sum\limits_{S} p_{S,SRT} \Theta_{SRT} (a_{R,T})\]

слєдује

\[\Theta_{SRT}(a_{R,\,T}) = \alpha_{S,SRT}^{(SR)}, ~~ a_{R,\,T} = \sum\limits_{S} e_{SRT} \alpha_{S,SRT}^{(SR)} = \sum\limits_{L} e_{L} \alpha_{LT^{-1}R^{-1},L}^{(LT^{-1})}.\]

Систем од \(a_{R,\,T}\) хочемо называти малым факторным системом од \(\mathfrak K_r\); велики факторны систем \(\alpha_{BC}^{(A)}\) и мали факторны систем \(a_{R,\,T}\) опрєдєљајут се взајемно. К продукту \(\varepsilon_{BC}^{(A)} = \alpha_{BC}^{(A)} \beta_{BC}^{(A)}\) двох великых факторных системов приналежи продукт \(f_{R,\,T} = a_{R,\,T}b_{R,\,T}\) малых факторных системов \(a_{R,\,T}\), \(b_{R,\,T}\), приредженых к \(\alpha_{BC}^{(A)}\), \(\beta_{BC}^{(A)}\); то непосрєдње слєдује из горњих формул. К великому јединчному факторному систему \(\alpha_{BC}^{(A)} = 1\) приналежи мали јединчны факторны систем \(a_{R,\,T} = 1\), и обратно.

Можно замєнити \(u_S\) и \(u_S'=u_Sr_S\), кде \(r_S\neq 0\) јест избрано либовољно из \(\mathfrak Z\). К \(u_S'\) приналежи мали факторны систем \(a_{R,T}'\)

\[
\text{(1)}\qquad u_R'u_T'=u_{RT}'a_{R,T}'
\]

\(a_{R,T}^{\prime}\) называјемо асоциованым с \(a_{R,T}\). Јест

\[u_R r_R u_T r_T = u_{RT} r_{RT} a'_{R,T} = u_R u_T r_R^T r_T = u_{RT} a_{R,T} r_R^T r_T\]

из чего

\[
\text{(2)}\qquad a'_{R,T}=a_{R,T}\frac{r_R^T r_T}{r_{RT}}
\]

слєдује.

К \(u_S' = u_S r_S\) приналежи слєдујучи систем псеудоматрицных јединиц

\[p_{S,ST}^{\prime}=p_{S,ST}r_{T}\]

мали факторны систем \(a'_{R,\,T}\) приналежи зато к великому факторному систему \(\alpha^{(A)\prime}_{BC}\), кторы јест асоциованы с факторным системом \(\alpha^{(A)}_{BC}\), приналежечим к \(a_{R,\,T}\); асоциоване мале факторне системы одповєдајут асоциованым великым факторным системам и обратно. Приреджење клас \(\{\alpha\}\) асоциованых великых факторных системов к приналежечим класам асоциованых малых факторных системов, \(\{a\}\), јест изоморфизм груп.

Нынє прєходимо к прєдставјењу од \(u_S\) посрєдством матриц реда н в \(\mathfrak{Z}\), кторе се сучствено разликује од досєга разгледаных прєдставјениј.

Ако \(\mathfrak Z\)-базу

\[
\text{(I)}\qquad k_{S_1}, \ldots, k_{S_n}
\]

од \(\mathfrak K_r\) множимо с права посрєдством \(u_S\), настава

\[
\text{(II)}\qquad k_{S_1}u_S, \ldots, k_{S_n}u_S.
\]

(Регуларну) матрицу \(A_S\), ктора прєводи (I) в (II),

\[
\text{(III)}\qquad (k_{S_1}u_S, \ldots, k_{S_n}u_S) = (k_{S_1}, \ldots, k_{S_n})A_S
\]

называмо матрицоју, ктора прєдставја \(u_S\). Прєдставјење се односи на все елементы од \(\mathfrak K_{r}\), кторе како \(u_S\) стварјајут автоморфизм од \(\mathfrak Z/P\), т.~ј. на все елементы \(u_Sr_S\), кде \(r_S \neq 0\) јест избрано либовољно из \(\mathfrak Z\). Вси ти елементи творет групу \(\mathfrak G^*\), ктора како нормалну подгрупу содрживаје мултипликативну групу \(\mathfrak Z^*\) од \(\mathfrak Z\); факторна група \(\mathfrak G^*/\mathfrak Z^*\) јест изоморфна с Галоисовој групоју \(\mathfrak G\) од \(\mathfrak Z/P\) тако, же всакы елемент једного побочного класа од \(\mathfrak G^*\) по одношењу к \(\mathfrak Z^*\) стварјаје автоморфизм од \(\mathfrak Z\), приреджены тому класу. Та група \(\mathfrak G^*\) прєдставја се управо указаным начином

\(u_S \mapsto A_S\)

Але то прєдставјење не јест хомоморфно, ни реципрочно хомоморфно; из

\(u_S \mapsto A_S\), \(u_T \mapsto A_T\)

не слєдује \(u_Su_T \mapsto A_SA_T\), але заради

\[
\begin{aligned}
(k_{S_1}, \ldots, k_{S_n})u_Su_T
&=(k_{S_1}, \ldots, k_{S_n})A_Su_T
=(k_{S_1}, \ldots, k_{S_n})u_TA_S^T\\
&=(k_{S_1}, \ldots, k_{S_n})A_TA_S^T:\ u_Su_T\mapsto A_TA_S^T.
\end{aligned}
\]

То прєдставјење од \(\mathfrak{G}^*\) называмо прєкриженым прєдставјењем. Ако \(h_{S_1}, \ldots, h_{S_n}\) јест друга база од \(\mathfrak{K}_r\) по одношењу к \(\mathfrak{Z}\), и

\[
\text{(IV)}\qquad (h_{S_1}, \ldots, h_{S_n}) = (k_{S_1}, \ldots, k_{S_n})M
\]

то из (III) и (IV) добыва се:

\[(h_{S_1},\ldots,h_{S_n})\,u_S=(h_{S_1},\ldots,h_{S_n})\,M^{-1}A_S\,M^S,\]

т.~ј.: Ако \(u_S\) јест при \(k\)-базе прєдставјено посрєдством \(A_S\), то при \(h\)-базе јест прєдставјено посрєдством \(M^{-1}A_SM^S\), ако \(M\) јест матрица, ктора трансформује \(k\)-базу в \(h\)-базу.

Формула (III), ктора дефинујеме прєкрижено прєдставјење, показује, же прєкрижено прєдставјење стоји в подобном одношењу к прстену \(\mathfrak K_r\), како обычно реципрочно прєдставјење (але записано с права вмєсто с лєва) к својему модулу прєдставјења. Разлика настава, зато же вмєсто релације комутативности

\[(m^*\tau^*)c = (m^*c)\tau^*\]
(види стр.~5)

при реципрочном модулу прєдставјења ту

\[
(kz)u_S=(ku_S)z^S\qquad k\text{ в }\mathfrak K_r,\quad z\text{ в }\mathfrak Z
\]

наступаје.

Прєкрижено прєдставјење од \(\mathfrak{G}^*\) можно такоже образовати посрєдством базы некого правого идеала \(\mathfrak r\) од \(\mathfrak{K}_r\). Всакы правы идеал од \(\mathfrak{K}_r\) стварјаје на тој начин класу еквивалентных прєдставјениј; при том прєдставјења \(u_S \mapsto A_S, u_S \mapsto B_S\) называјут се еквивалентными, ако обстаја матрица \(M\) с \(B_S = M^{-1}A_SM^S\).

Вмєсто ограничења на прєдставјења од \(\mathfrak G^*\) в \(\mathfrak Z\), створјене правыми идеалами од \(\mathfrak K_r\), нынє можно обче дефинирати прєкрижено прєдставјење \(m\)-того ступења од \(\mathfrak G^*\) в \(\mathfrak Z\) тако:

\begin{enumerate}
\item Всакому \(u_S\) одповєдаје матрица \(A_S\) реда \(m\) над \(\mathfrak Z\).
\item Из \(u_S \mapsto A_S\), \(u_T \mapsto A_T\) слєдује \(u_Su_T \mapsto A_TA_S^T\).
\end{enumerate}

Изберемо правы \(\mathfrak Z\)-модул ранга \(m\)

\[\mathfrak{M}=x_1\mathfrak{Z}+\cdots+x_m\mathfrak{Z}\]

и установјењем

\[(x_1u_S,\ldots,x_mu_S)=(x_1,\ldots,x_m)A_S\]

чинимо го правым \(\mathfrak G^*\)-модулом с релацијоју комутативности

\[(mz) u_S = (mu_S) z^S\]

цєлком аналогно конструкцији модула прєдставјења к заданому прєдставјењу на стр.~5. Нынє показујемо, како на стр.~20, же \(\mathfrak{M}\) поставаје конечным правым \(\mathfrak K_r\)-модулом посрєдством установјења

\[m \cdot z u_S = mz \cdot u_S = mu_S \cdot z^S\]

обче

\[m \cdot \sum z_S u_S = \sum m z_S \cdot u_S = \sum m u_S \cdot z_S^S.\]

При том никако не јест потрєбно рачунање аналогно рачунању на стр.~20, зато же при нашем установјењу употрєбили јесмо установјену базу.

Але по M.З. § 18 всакы конечны просты \(\mathfrak K_r\)-модул \(\mathfrak{M}\) јест операторно изоморфны с простым правым идеалом од \(\mathfrak K_r\); из тога аналогно теорему 2 на стр.~22 закључујемо, же обстаја само \emph{једна} класа иредуцибилных прєкриженых прєдставјениј од \(\mathfrak{G}^*\) в \(\mathfrak Z\) --- јменно она, ктора приналежи к простым правым идеалам.

Ступењ тога иредуцибилного прєдставјења јест \(t\), кде \(t\) јест индекс (абсолутно число експонента) од \(\mathfrak K\). То се добыва простым считањем рангов. Зато же ступењ јест равны рангу \((\mathfrak r:\mathfrak Z)\) простого правого идеала над \(\mathfrak Z\). Нынє на стр.~22 имаје мєсто:

\[
(\mathfrak Z:P)=n=rt;\qquad (\mathfrak K_r:P)=n^2;\qquad (\mathfrak r:P)=n\cdot t
\]

послєдње заради \(\mathfrak K_r = \mathfrak{r}_1 + \cdots + \mathfrak{r}_r\).

Из
\[(\mathfrak{r}: P) = (\mathfrak{r}: \mathfrak Z) \cdot (\mathfrak Z: P)\]
или \(n \cdot t = (\mathfrak{r}: \mathfrak Z) \cdot n\) слєдује зато \((\mathfrak{r}: \mathfrak Z) = t\).
% END BOOK_S25

% BEGIN BOOK_S26
\subsection*{§ 26. Множење прєкриженых прєдставјениј}\label{multiplikation-der-verschraenkten-darstellungen}

Досели јесмо всегда разгледали прєкрижено прєдставјење како прєдставјење од \(\mathfrak{G}^*\). Але јего можно разгледати такоже како прєдставјење Галоисовој групы \(\mathfrak{G}\) од \(\mathfrak Z/P\), и то тако, же всакому елементу \(S\) од \(\mathfrak{G}\) приреджајемо матрицу \(A_S\), ктора прєдставја елемент \(u_S\) из \(\mathfrak K_r\), стварјајучи автоморфизм \(S\) посрєдством трансформације -- \(u_S^{-1}zu_S=z^S\) --.

За означење такого прєдставјења \(S\mapsto A_S\)\srcfn{1}{Тако прєдставјење (независно од хиперкомплексных системов)
\[
S\mapsto A_S;\quad T\mapsto A_T;\quad ST\mapsto A_TA_S^T=A_{ST}\cdot a_{S,T}^{-1}
\]
(или транспоновано прєдставјење) было одправным пунктом Спеисера (Math. Зсцхр. 5) при опрєдєљењу факторных системов \(a_{S,T}\). Те факторне системы показывајут се зато идентичными с ту из другој основы уведеными.} трєба сказати, в ктором \(\mathfrak K_r\) оно јест взето, кторе \(u_S\) и кторе \(k_{S_1},\ldots,k_{S_n}\) од \(\mathfrak K\) над \(\mathfrak Z\) сут употрєбјене. Замєна базы даваје еквивалентно прєдставјење, а замєна од \(u\) даваје асоциоване факторне системы \(a_{S,\,T}\); сучствене сут то јест \(\mathfrak K_r\), в кторых прєдставјење јест взето, т.~ј. класа \(\{a_{S,\,T}\}\) асоциованых факторных системов.

Нехај нынє будут дане два така прєдставјења од \(\mathfrak G\).

\[
\begin{array}{lll}
S\mapsto u_S\mapsto A_S & \text{належечи мали факторны систем} & a_{R,T},\\
S\mapsto v_S\mapsto B_S & \text{належечи мали факторны систем} & b_{R,T}.
\end{array}
\]

Под Кронекеровым продуктом \(A\times B\) двох матриц \(A=(a_{ik})\), \(B=(b_{\mu\nu})\) разумєјемо матрицу \(A\times B=(a_{ik}b_{\mu\nu})\).

\noindent\textbf{Тврджење.} \emph{Такоже \(S\mapsto A_S\times B_S\) јест прєкрижено прєдставјење од \(\mathfrak G\), и оно приналежи к продукту \(\mathfrak K_r\times\mathfrak S_s\), ако \(A_S\) приналежи к \(\mathfrak K_r\), а \(B_S\) к \(\mathfrak S_s\).}

\noindent\emph{Доказ.} Прєдставјење \(S\mapsto A_S\times B_S\) јест прєкрижено прєдставјење. Нехај буде
\[
S\mapsto A_S\times B_S,\qquad T\mapsto A_T\times B_T
\]
и
\[
ST\mapsto A_{ST}\times B_{ST}
=(A_TA_S^T\times B_TB_S^T)a_{S,T}^{-1}b_{S,T}^{-1},
\]
зато трєба доказати:
\[
(A_T\times B_T)(A_S\times B_S)^T a_{S,T}^{-1}b_{S,T}^{-1}
=(A_TA_S^T\times B_TB_S^T)a_{S,T}^{-1}b_{S,T}^{-1}.
\]
То але слєдује из слєдујучего разважања:

Ако матрице \(A\) и \(B\) прєдставимо посрєдством матричных јединиц \(c_{ik}\) и \(d_{\mu\nu}\) (кторе не имајут ничто обче с матричными јединицами од \(\mathfrak K_r\) или \(\mathfrak S_s\))
\[
A=\sum c_{ik}\alpha_{ik};\qquad
B=\sum d_{\mu\nu}\beta_{\mu\nu},\qquad\text{тогда}
\]
\[
A\times B=B\times A
=\sum c_{ik}d_{\mu\nu}\cdot\alpha_{ik}\cdot\beta_{\mu\nu},
\]
т.~ј. оно јест равно матрици в новых матричных јединицах
\(c_{ik}d_{\mu\nu}=d_{\mu\nu}c_{ik}\). Множење \(A\)-матриц и \(B\)-матриц зато једнозначно одповєдаје множењу елементов в прємом продукту двох матричных прстенов:
\[
\sum c_{ik}\mathfrak Z\times\sum d_{\mu\nu}\mathfrak Z.
\]
Зато всака \(A\)-матрица јест комутативна с всакој \(B\)-матрицој (але \(A\)-матрице ни \(B\)-матрице не сут меджу собоју комутативне); и та комутативност всакој \(A\)-матрице с всакој \(B\)-матрицоју даваје управо --- по множењу с \(a_{S,T}^{-1}b_{S,T}^{-1}\) --- релацију, ктору трєба доказати.

\section*{Глава VI. Теорија прєкриженых продуктов}\label{kapitel-vi.-theorie-der-verschraenkten-produkte}
% END BOOK_S26

% BEGIN BOOK_S27
\subsection*{§ 27. Прєдставјење од \(\mathfrak K_r\) како прєкриженых продуктов}

Теорија прєкриженых продуктов, развијена в § 24 из факторных системов, нынє буде поновно независно основана. Значење тых прєкриженых продуктов јест в том, же Галоисово поље и јего групы сут \emph{заједно} описане в \(\mathfrak K_r\).

\noindent\textbf{Теорем 1.} \emph{Ако \(\mathfrak Z\) јест Галоисово разпадно поље од \(\mathfrak K_r\), то јест максимално комутативно подпоље, и ако \(\mathfrak G\) јест јего Галоисова група, то \(\mathfrak K_r\) јест равно прєкриженому продукту \(\mathfrak Z\cdot\mathfrak G\). Обратно, всакы прєкрижены продукт \(\mathfrak Z\cdot\mathfrak G\) стварјаје \(\mathfrak K_r\) такого вида.}

\noindent\textbf{Доказ и дефиниције. Дефиниција 1.} Нехај \(\mathfrak G=\{E,S,\ldots,T\}\) буде Галоисова група од \(\mathfrak Z\); равност од \(\mathfrak K_r\) прєкриженому продукту значи:
\[
\mathfrak K_r=\mathfrak Zu_E+\cdots+\mathfrak Zu_T
\]
с \(zu_S=u_Sz^S\) и \(u_Su_T=a_{S,T}u_{ST}\) (\(a_{S,T}\) в \(\mathfrak Z\) сут мале факторне системы).

\noindent\textbf{Дефиниција 2.} Стварјајуча група \(\mathfrak G^*\) јест опрєдєљена како множство регуларных елементов \(g\) из \(\mathfrak K_r\), кторе посрєдством трансформације прєводет \(\mathfrak Z\) како цєлост в себе. \(g^{-1}\mathfrak Zg=\mathfrak Z\), то јест \(g^{-1}zg=z^S\), с \(S\) в \(\mathfrak G\). зато \(\mathfrak G\) јест групово-хомоморфны образ од \(\mathfrak G^*\), зато же всакы автоморфизм од \(\mathfrak Z\) јест створјен некым \(g\), \(\mathfrak G^*\to\mathfrak G\), и даље
\[
\mathfrak G\cong\mathfrak G^*/\mathfrak T^*
\quad\text{с}\quad
\mathfrak T^*\supseteq\mathfrak Z^*,
\]
(же \(\mathfrak T^*=\mathfrak Z^*\), буде доказано поздєје --- јешче не было доказано, же \(\mathfrak Z\) было такоже максимално комутативно подпрстен. То можно доказати и независно).

\noindent\textbf{1.} \emph{Доказ прєдставјења како прєкриженого продукта.} Нехај \(u_E,u_S,\ldots,u_T\) будут елементи из \(\mathfrak K_r\), кторе стварјајут автоморфизмы \(E,\ldots,T\), то јест \(zu_S=u_Sz^S\), особливо \(u_E=e\) заради 2. Трєба показати, же сума
\[
\mathfrak M=\{\mathfrak Zu_E,\ldots,\mathfrak Zu_T\}
\]
јест прєма и равна \(\mathfrak Zu_E+\cdots+\mathfrak Zu_T\); тогда она јест равна \(\mathfrak K_r\) заради једнакости рангов. То але слєдује тако (\foreign{Schwarz}):

\(\mathfrak M\) јест лєвы и правы \(\mathfrak Z\)-модул, то јест модул прєдставјења од \(\mathfrak Z\) в \(\mathfrak Z\), и посрєдством \(\mathfrak M\) стварјајут се \(n\) различне прєдставјења; тым самым --- зато же \(\mathfrak Z\) јест комутативно поље првого вида --- \(\mathfrak M\) имаје ранг принајмење \(n\); иначе доказ тече управо тако.

\noindent\textbf{2.} \(\mathfrak G\cong\mathfrak G^*/\mathfrak Z^*\), т.~ј. \(\mathfrak T^*=\mathfrak Z^*\); само елементи из \(\mathfrak Z\) индуцирајут идентитет на \(\mathfrak Z\) посрєдством трансформације.

Нехај именно буде \(tz=zt\) за всако \(z\) из \(\mathfrak Z\); положимо
\[
t=z_0+z_1u_S+\cdots+z_{n-1}u_T,
\]
зато
\[
zt=zz_0+\cdots+zz_{n-1}u_T
\]
и
\[
tz=z_0z+\cdots+z_{n-1}u_Tz
=z_0z+z^{S^{-1}}z_1u_S+\cdots+z^{T^{-1}}z_{n-1}u_T,
\]
зато заради прємој сумы:
\[
(z-z^{S^{-1}})z_1u_S=0,\ldots,
(z-z^{T^{-1}})z_{n-1}u_T=0
\]
и из тога, зато же елементи \(u_S,\ldots\) сут регуларне:
\[
(z-z^{S^{-1}})z_1=0,\ldots,
(z-z^{T^{-1}})z_{n-1}=0,
\qquad\text{за всако }z.
\]
(или зато же ранг јест равны \(n\)). Але по дефиницији за \(S,\ldots\), и тако даље (все неравне \(E\)), обстаја принајмење једно \(z\), тако же \(z-z^{S^{-1}}\neq0\); и неко \(\bar z\), тако же \(\bar z-\bar z^{T^{-1}}\neq0\); зато добыва се \(z_i=0\), \(i=1,\ldots,(n-1)\), и зато \(t=z_0\) в \(\mathfrak Z^*\), тым јест доказано.

\noindent\textbf{3.} Елементи \(u\) изполњајут законы множења: \(u_Su_T=a_{S,T}u_{ST}\), при чем \(a_{S,T}\) изполњајут прєкрижены асоциативны закон
\[
a_{R,S}\cdot a_{RS,T}=a_{S,T}^{R^{-1}}\cdot a_{R,ST}
\]
(\(a_{S,T}\) сут мали факторны систем).

Зато же по 2. добыва се:
\[
\mathfrak Z^*u_S\cdot\mathfrak Z^*u_T=\mathfrak Z^*u_{ST};
\]
зато
\[
u_Su_T=a_{S,T}u_{ST};
\]
с \(a_{S,T}\) из \(\mathfrak Z^*\), то јест из \(\mathfrak Z\) и \(\neq0\). Заради \(u_R\cdot u_Su_T=u_Ru_S\cdot u_T\) даље добыва се
\[
\begin{aligned}
u_Ra_{S,T}u_{ST}
&=a_{S,T}^{R^{-1}}\cdot u_Ru_{ST}
 =a_{S,T}^{R^{-1}}a_{R,ST}u_{RST}\\
&=a_{R,S}u_{RS}\cdot u_T
 =a_{R,S}a_{RS,T}u_{RST},
\end{aligned}
\]
Тым сут подане законы множења.

\noindent\textbf{4. Обрат:} прєкрижены продукт \(=\mathfrak K_r\). Прєма сума:
\[
\mathfrak M=\mathfrak Zu_E+\cdots+\mathfrak Zu_T
\]
буде опрєдєљена како прстен посрєдством релациј: \(zu_S=u_Sz^S\), \(u_Su_T=a_{S,T}u_{ST}\), (кде \(a_{S,T}\) изполњајут прєкрижене асоциативне релације; то јест \(u_Ru_Su_T\) јест једнозначно опрєдєљено):
\[
\begin{array}{ll}
1.&z_1\cdot z_2u_S=z_1z_2\cdot u_S;\\[2pt]
2.&zu_S\cdot u_T=z\cdot u_Su_T
\end{array}
\]
тым \(\mathfrak M\) поставаје асоциативным прстеном.

\noindent\textbf{Обратны теорем 2.} \emph{Прєкрижены продукт \(\mathfrak Z\cdot\mathfrak G\) јест двустранно просты, с \(P\) како центром (по приступах Р. Брауера).} Доказ основываје се на слєдујучих помочных лемах:

\noindent\textbf{Помочна лема 1.} \emph{Всакы елемент, кторы јест по елементах комутативны с \(\mathfrak Z\), приналежи к \(\mathfrak Z\).}

\noindent\emph{Доказ.} Нехај \(w=\sum_S c_{\lambda_S}u_S\), (\(c_{\lambda_S}\) в \(\mathfrak Z\)) и нехај буде прєдположено: \(wz=zw\) за всако \(z\) из \(\mathfrak Z\). зато
\[
\sum_S zc_{\lambda_S}u_S
 =\sum_S c_{\lambda_S}u_Sz
 =\sum_S c_{\lambda_S}z^{S^{-1}}u_S
\]
зато \(zc_{\lambda_S}=c_{\lambda_S}z^{S^{-1}}\), или \(c_{\lambda_S}(z-z^{S^{-1}})=0\), за всако \(z\) и \(S\) (линеарна независност од \(u\)!). Из тога але слєдује: \(c_{\lambda_S}=0\) за \(S\neq E\), зато же за всако тако \(S\) обстаја принајмење \emph{једно} \(z\), тако же \(z-z^{S^{-1}}\neq0\); зато \(w\) јест в \(\mathfrak Z\).

\noindent\textbf{Помочна лема 2.} \emph{Всакы \(\mathfrak Z\)-двојны модул \(\mathfrak A\), кторы јест двојно изоморфны с \(\mathfrak Zu_S=u_S\mathfrak Z\) (за \(\mathfrak Z\) како лєвы и правы оператор), јест идентичны с \(\mathfrak Zu_S\).}

\noindent\emph{Доказ 2.} \(\mathfrak A\) јест изоморфны с \(\mathfrak Zu_S\) како лєвы модул; ако зато \(u_S\mapsto a\) из \(\mathfrak A\), то \(\mathfrak Zu_S\mapsto\mathfrak Za\), то \(\mathfrak A=\mathfrak Za\). Из двојного изоморфизма слєдује (1) \(za=az^S\) (заради \(u_Sz\mapsto az\)); и из тога \(\mathfrak A=\mathfrak Zu_S\). Ако именно положимо: \(a=u_Su_S^{-1}a=u_Sw\), добыва се: \(za=zu_Sw=u_Sz^Sw\), але заради (1): \(za=u_Swz^S\), зато (по множењу с \(u_S\)): \(z^Sw=wz^S\) за всако \(z\) из \(\mathfrak Z\). Зато такоже \(z^S\) пробєгаје цєло \(\mathfrak Z\), \(w\) јест по елементах комутативны с \(\mathfrak Z\), зато по 1. приналежи к \(\mathfrak Z\).

\noindent\textbf{Помочна лема 3.} \emph{Всакы \(\mathfrak Z\)-двојны модул из \(\mathfrak Z\cdot\mathfrak G\) имаје форму: \(\mathfrak Zu_{S_1}+\cdots+\mathfrak Zu_{S_k}\), кде \(u_{S_1}\) означаје ныкоје \(k\) различне из \(u_S,\ldots,u_T\).}

\noindent\emph{Доказ.} Прєкрижены продукт \(\mathfrak Z\cdot\mathfrak G\) јест полно редуцибилны како \(\mathfrak Z\)-двојны модул, зато же всакы суманд \(\mathfrak Zu_S\) имаје ранг 1. Како двојне модулы, вси суманды такоже приналежет к различным класам, зато не сут двојно изоморфне. Заради полној редуцибилности всакы \(\mathfrak Z\)-двојны модул \(\mathfrak M\) јест прємы суманд и сам полно редуцибилны. \(\mathfrak M=\sum\mathfrak A_{S_i}\), кде \(\mathfrak A_{S_i}\) јест изоморфны с \(\mathfrak Zu_{S_i}\); и зато по 2. јест идентичны с \(\mathfrak Zu_{S_i}\). (Вси \(\mathfrak A_{S_i}\) приналежет к различным класам.)

\noindent\textbf{Помочна лема 4.} \emph{Всако прстен из \(\mathfrak Z\cdot\mathfrak G\), кторо содрживаје \(\mathfrak Z\), особливо само \(\mathfrak Z\cdot\mathfrak G\), јест двустранно просто. --- Подпрстены, кторе содрживајут \(\mathfrak Z\), једнозначно одповєдајут подгрупам од \(\mathfrak G\).}

\noindent\emph{Доказ.} Всако прстен \(\mathfrak o\), кторо содрживаје \(\mathfrak Z\), јест \(\mathfrak Z\)-двојны модул, зато по 3. имаје форму \(\sum\mathfrak Z\cdot u_{S_i}\); зато же \(\mathfrak o\) јест прстен, продукт двох \(u_{S_i}\) поновно приналежи к \(\mathfrak o\), зато такоже приналежечи \(\mathfrak Z\)-модул. Зато \(u_{S_i}\) творет групу до факторных системов, зато же јих јест конечно много. зато \(\mathfrak o=\mathfrak Z\cdot\mathfrak H\), кде \(\mathfrak H\) јест подгрупа од \(\mathfrak G\). (Точније, \(\mathfrak Z\cdot u_{S_i}\) творет мултипликативны систем.) Ако нынє \(\mathfrak a\) јест либовољны двустранны идеал \(\neq0\) из \(\mathfrak o\), то он јест \(\mathfrak o\)-двојны модул, и тым \(\mathfrak Z\)-двојны модул и подмодул од \(\mathfrak o\). зато вмєстє с \(\mathfrak Zu_{S_i}\) содрживаје цєлу групу \(\mathfrak H\), и зато поставаје \(\mathfrak o=\mathfrak Z\cdot\mathfrak H\).

\noindent\textbf{Помочна лема 5.} \emph{\(P\) јест центар од \(\mathfrak Z\cdot\mathfrak G\).}

\noindent\emph{Доказ.} Ако \(w\) јест по елементах комутативны с \(\mathfrak Z\cdot\mathfrak G\), то особливо с \(\mathfrak Z\), зато по 1. приналежи к \(\mathfrak Z\). Зато из \(wu_S=u_Sw^S\) слєдује \(w=w^S\), ако јест комутативны. зато \(w\) приналежи к \(P\).

\noindent\emph{Примєтка 1.} Доказ показује, же теорем имаје мєсто такоже, когда \(\mathfrak Z\) јест бесконечно алгебраично Галоисово поље над \(P\). Зато же теоремы о полној редуцибилности оставајут в силє (Krull Захлринге, Math. Ann. 99). При том \(\mathfrak Z\cdot\mathfrak G\) јест систем всих конечных сум \(\sum_{S_\lambda}u_{S_\lambda}\cdot c_{S_\lambda}\).

Само при 4. трєба за \(\mathfrak o\) прєдположити \(\mathfrak o=\mathfrak Z\cdot\mathfrak H\), а то всегда имаје мєсто за полны продукт \(\mathfrak Z\cdot\mathfrak G\), зато же из належења од \(u_{S_i}u_{S_k}\) нынє не можно закључити групово својство система, ако \(\mathfrak o\) јест бесконечно над \(\mathfrak Z\).

\noindent\emph{Примєтка 2.} При прємом основању ту трєба присојединити теорију прєкриженого матричного прєдставјења, развијену в §§ 25/26. В слєдујучем, на концу § 28, употрєбја се, же иредуцибилно прєкрижено прєдставјење имаје ступењ \(t\), кде \(t\) јест абсолутно число компонентов (индекс) од \(\mathfrak K\). Поновно трєба указати, же та теорија прєдставјења даваје идентитет факторных системов, уведеных ту како константы множења, с појмом обычным в литературє.
% END BOOK_S27

% BEGIN BOOK_S28
\subsection*{§ 28. Продуктны теорем за факторне системы}\label{produktsatz-fuer-faktorensysteme}

Нехај буде
\[
\mathfrak K_r=\mathfrak Z\cdot\mathfrak G
=\mathfrak Zu_E+\cdots+\mathfrak Zu_T,
\qquad u_Su_T=u_{ST}a_{S,T},
\]
\[
\overline{\mathfrak K}_r
=\overline{\mathfrak Z}\,\overline{\mathfrak G}
=\overline{\mathfrak Z}\bar u_E+\cdots+\overline{\mathfrak Z}\bar u_T,
\qquad
\bar u_S\bar u_T=\bar u_{ST}\bar b_{S,T},
\]
(\(\mathfrak Z\) и \(\overline{\mathfrak Z}\) сут изоморфне Галоисове поља, а \(\mathfrak G\), \(\overline{\mathfrak G}\) њих групы), зато добыва се
\[
\mathfrak K_r\times\overline{\mathfrak K}_r
=\widetilde{\mathfrak Z}\,\widetilde{\mathfrak G}\times P_n
=\mathfrak A_f\times P_n
=\mathfrak A_j
\]
с \(\tilde u_S\tilde u_T=\tilde u_{ST}\tilde c_{S,T}\) (\(\widetilde{\mathfrak Z}\) јест изоморфно с \(\mathfrak Z\) и \(\overline{\mathfrak Z}\)). То слєдује из рачунања рангов и из тога, же \(\mathfrak Z\) јест такоже разпадно поље).

\noindent\textbf{Тврджење.} \emph{Добыва се \(\tilde c_{S,T}=\tilde a_{S,T}\tilde b_{S,T}\), кде \(\tilde a,\tilde b\) сут с \(a\) и \(\bar b\) изоморфне елементи в \(\widetilde{\mathfrak Z}\).}

Доказ основываје се на том, же прстен автоморфизмов либовољного лєвого идеала из \(\mathfrak A_j\), кторы имаје абсолутну должину \(n\), јест изоморфно с \(\mathfrak A_f=\widetilde{\mathfrak Z}\widetilde{\mathfrak G}\), и же то прстен автоморфизмов можно прємо опрєдєлити при погодном лєвом идеалу из \(\mathfrak A_j\). Обче имаје мєсто слєдујучи теорем, именно с \(s\) вмєсто \(n\), але заради удобности изречены за \(\mathfrak K_r\):

\noindent\textbf{1. Теорем 1.} \emph{Прстен автоморфизмов лєвого идеала \(\mathfrak l\) абсолутној должины \(n\) из \(\mathfrak K_{rs}=\mathfrak K_r\times P_s\) јест изоморфно с \(\mathfrak K_r\).}

Доказ слєдује из слєдујучих помочных лем.

\noindent\textbf{Помочна лема 1.} \emph{Вси лєве идеалы абсолутној должины \(n\) из \(\mathfrak K_{rs}\) сут операторно изоморфне. Њих прстены автоморфизмов сут зато прстеново изоморфне. Особливо \(\mathfrak l=\mathfrak K_r\times\mathfrak b\) имаје абсолутну должину \(n\), при чем под \(\mathfrak b\) разумєјемо просты идеал из \(P_s\).}

\noindent\emph{Доказ.} То, же лєве идеалы имајут абсолутну должину \(n\), т.~ј. разпадајут се на \(n\) абсолутно неразложимых простых идеалов, имаје за послєдицу једнаку должину в \(\mathfrak K_{rs}\), т.~ј. в \(\mathfrak K_{rs}\) наступаје равнако много простых идеалных сумандов. Из тога слєдује њих операторны изоморфизм, а тым и прстеновы изоморфизм прстенов автоморфизмов. Даље \(\mathfrak K_r=\mathfrak Z\cdot\mathfrak G\) имаје абсолутну должину \(n\), т.~ј. ранг јест \(n^2\), зато \(\mathfrak K_r\times P_s\) имаје абсолутну должину \(n\cdot s\), и зато \(\mathfrak K_r\times\mathfrak b\) имаје абсолутну должину \(n\).

\noindent\textbf{Помочна лема 2.} \emph{Ако \(\mathfrak o\) јест прстен с јединицоју, \(\mathfrak l=\mathfrak o e_1\) јест лєвы идеал и прємы суманд, то јест \(e_1\) јест идемпотент, то прстен автоморфизмов од \(\mathfrak l\) јест изоморфно с \(e_1\mathfrak o e_1\).}

\noindent\emph{Доказ.} Множење с елементом \(e_1re_1\) даваје операторны хомоморфизм од \(\mathfrak l\) в себе; различне елементи \(e_1re_1\) давајут различне хомоморфизмы, зато же при том к \(e_1\) сут приреджене различне елементи, и всакы хомоморфизм јест створјен тако: \(e_1\mapsto re_1\) и \(re_1=e_1re_1\) заради \(e_1e_1=e_1\). Сума и продукт в \(e_1\mathfrak o e_1\) але одповєдајут суми и продукту приредженых хомоморфизмов.

\noindent\textbf{Помочна лема 3.} \emph{Ако положимо \(P_s=\sum_1^s c_{ik}P\) с матричными јединицами \(c_{ik}\) и \(\mathfrak b=c_{11}P+\cdots+c_{s1}P\), то како прстен автоморфизмов од \(\mathfrak K_r\times\mathfrak b\) добыва се прстен \(\mathfrak K_rc_{11}\), кторо јест изоморфно с \(\mathfrak K_r\).}

\noindent\emph{Доказ.} Зато же \(c_{11}\) јест по елементах комутативны с \(\mathfrak K_r\) и ниједин елемент из \(\mathfrak K_r\) не јест анулованы, зато же \(c_{11}\) јест стварјајучи идемпотент од \(\mathfrak K_r\times P_s\), и јест
\[
c_{11}\cdot\mathfrak K_r\times P_s\cdot c_{11}
=\mathfrak K_r\times c_{11}P_sc_{11}
=\mathfrak K_rc_{11}.
\]
Из тых помочных лем прємо слєдује теорем 1.

\noindent\textbf{2. Опрєдєљење и разклад од \(\mathfrak K_r\times\overline{\mathfrak K}_r\) на лєве идеалы абсолутној должины \(n\).} По дефиницији добыва се
\[
\mathfrak A_f=\mathfrak K_r\times\overline{\mathfrak K}_r
=\mathfrak Z\cdot\mathfrak G\times\overline{\mathfrak Z}\cdot\overline{\mathfrak G},
\]
при чем елемент с чртоју јест по елементах комутативны с елементом без чрты, зато
\[
\mathfrak A_f
=(\mathfrak G\times\overline{\mathfrak G})\cdot(\mathfrak Z\times\overline{\mathfrak Z})
=\mathfrak G\times\overline{\mathfrak G}\cdot(\mathfrak Ze_1+\cdots+\mathfrak Ze_n)
\]
с \(\mathfrak Ze_i=\overline{\mathfrak Z}e_i\), (множење Галоисового поља самого с собоју). Зато же вси суманды имајут једнакы ранг по одношењу к \(\mathfrak Z\), а зато такоже једнаку абсолутну должину --- абсолутна должина од \(\mathfrak A_f\) але јест \(=n^2\) --- всакы суманд имаје абсолутну должину \(n\). В слєдујучем за основу буде взето
\[
\mathfrak l=(\mathfrak G\times\overline{\mathfrak G})\mathfrak Ze_1,
\]
при чем релација \(ze_1=e_1\bar z\) установјаје опрєдєљены изоморфизм меджу \(\mathfrak Z\) и \(\overline{\mathfrak Z}\).

\noindent\textbf{3. Прєкрижено прєдставјење од \(e_1\mathfrak A_fe_1\), и тым доказ продуктного теорема.} --- По 1. и 2. факторне системы продукта \(\mathfrak K_r\times\overline{\mathfrak K}_r\) сут изоморфне с факторными системами прєкриженого прєдставјења од \(e_1\mathfrak A_fe_1=\mathfrak A\).

\noindent\textbf{Теорем 2.} \emph{\(\mathfrak A=\widetilde{\mathfrak Z}\cdot\widetilde{\mathfrak G}\) с \(\widetilde{\mathfrak Z}=\mathfrak Ze_1=\overline{\mathfrak Z}e_1\), \(\widetilde{\mathfrak G}=\{\ldots,e_1u_S\bar u_Se_1,\ldots\}\), при чем
\[
e_1u_S\bar u_Se_1\cdot e_1u_T\bar u_Te_1
=e_1u_{ST}\bar u_{ST}e_1\cdot a_{S,T}e_1\bar b_{S,T}e_1,
\]
тако же с \(\tilde a_{S,T}=a_{S,T}e_1\) и \(\tilde b_{S,T}=\bar b_{S,T}e_1\) продуктны теорем јест доказан; (\(\tilde a\) и \(\tilde b\) сут в том самом пољу \(\widetilde{\mathfrak Z}\)).}

\noindent\emph{Доказ.} \(e_1\) поставаје јединица прстена; елементи коњуговане с \(\tilde z=ze_1\) сут зато дане посрєдством \(\tilde z^Se_1\); трєба показати, же елементи \(u_S\bar u_Se_1\) индуцирајут те субституције, чиме буде доказано прєкрижено прєдставјење.

\noindent\textbf{Помочна лема 1.} \(e_1u_S\bar u_S=u_S\bar u_Se_1\), зато \(=e_1u_S\bar u_Se_1\).

\noindent\emph{Доказ.} Положимо \(u_S^{-1}e_1u_S=e^S\)\srcfn{1}{Тым установјењем добыва се \(e_1=e^E\).} или \(e_1u_S=u_Se^S\); тогда \(e^S\) пробєгајут коњуговане \(e_1,e_2,\ldots,e_n\). Тогда имаје мєсто (1)\srcfn{2}{заради \(\sum a_i\bar A_i\bar u_R=\sum a_i\bar u_R\bar A_i^R=\bar u_R\sum(a_i^R)^{R^{-1}}\bar A_i^{-R}\).}
\[
e_1\bar u_R=\bar u_Re^{R^{-1}}.
\]
Зато же \(e_1=\sum a_i\bar A_i=\sum a_i^R\bar A_i^R\), кде \(a_i\) пробєгаје либовољну \(\mathfrak Z\)-базу, а \(\bar A_i\) приналежечу комплементарну \(\overline{\mathfrak Z}\)-базу, зато то само имаје мєсто за \(a_i^R\) и \(\bar A_i^R\). зато \(e_1\bar u_R=\bar u_Re^{R^{-1}}\), чиме (1) јест доказано. Из (1) слєдује
\[
e_1u_S\bar u_S=u_Se^S\bar u_S=u_S\bar u_S\cdot e^{SS^{-1}}=u_S\bar u_Se_1,
\]
чиме јест доказана помочна лема 1.

\noindent\textbf{Помочна лема 2.} \(ze_1\cdot e_1u_S\bar u_Se_1=e_1u_S\bar u_Se_1\cdot z^Se_1\).

\noindent\emph{Доказ.} 
\[
ze_1\cdot e_1u_S\bar u_Se_1
=zu_S\bar u_Se_1\quad(\text{Помочна лема 1})
=u_Sz^S\bar u_Se_1
=u_S\bar u_Se_1\cdot z^Se_1
=b=e_1b,
\]
(\(z\) јест комутативны с \(\bar u\) и \(e_1\)). --- Нехај \(u'=u_S\bar u_Se_1\). Помочна лема 2 изрека, же \(u'\) индицираје субституције од \(\mathfrak Ze_1\), тако же имаје мєсто прєкрижено прєдставјење подано в теорему.

\noindent\textbf{Помочна лема 3.} \emph{Факторны систем поставаје равны \(a_{S,T}e_1\cdot\bar b_{S,T}e_1\).}

\noindent\emph{Доказ.}
\[
\begin{aligned}
e_1u_S\bar u_Se_1\cdot e_1u_T\bar u_Te_1
&=u_S\bar u_S\cdot u_T\bar u_Te_1
 =u_{ST}a_{S,T}\cdot\bar u_{ST}\bar b_{S,T}e_1\\
&=u_{ST}\bar u_{ST}\cdot a_{S,T}\bar b_{S,T}e_1
 =e_1u_{ST}\bar u_{ST}e_1\cdot a_{S,T}e_1\cdot\bar b_{S,T}e_1.
\end{aligned}
\]
Тым јест доказан теорем из 2., а зато и продуктны теорем за факторне системы.

\noindent\textbf{4. Продуктны теорем за класы асоциованых факторных системов:}
\[
\{a_{S,T}\}\cdot\{b_{S,T}\}=\{a_{S,T}\cdot b_{S,T}\}.
\]
То прємо слєдује из 3., зато же прєход к новым стварјајучим елементам \(v\) и \(\bar v\) означаје такоже такы прєход к \(v\bar ve_1\), зато же \(u=vc\) и \(\bar u=\bar v\bar d\) давајут
\[
u\bar ue_1=vc\cdot\bar v\bar de_1=v\bar v\cdot c\bar de_1=v\bar ve_1\cdot c\bar de_1.
\]

\noindent\textbf{5.} Ако \(\{a\}\) јест факторны систем од \(\{\mathfrak K\}\), а \(t\) јест индекс од \(\{\mathfrak K\}\), то \(\{a\}^t\) јест равны јединчному класу. То слєдује из тога, же иредуцибилно прєкрижено прєдставјење од \(\mathfrak G\) имаје ступењ \(t\), управо како Шуров доказ по прєдставјењу од \(p_{ik}\) (§ 24, стр. 37).

\noindent\textbf{6.} Изоморфизм меджу групоју \(\mathscr K_{\mathfrak Z}\) од \(\{\mathfrak K\}\) к \(\mathfrak Z\), т.~ј. с \(\mathfrak Z\) како разпадным пољем, и групоју клас \(\{a_{S,T}\}\) асоциованых факторных системов. По конструкцији в § 27 обстаја једнозначна кореспонденција меджу \(\{\mathfrak K\}\) и \(\{a_{S,T}\}\); по 4. приреджење јест хомоморфно, зато имамо изоморфизм. Из тога особливо слєдује, же јединчне класы одповєдајут једне другим:

\noindent\emph{Прєкрижены продукт јест тогда и само тогда равны \(P\), то јест матричному прстену в центру, когда факторны систем јест асоциованы с јединчным системом.}
% END BOOK_S28

% BEGIN BOOK_S29
\subsection*{§ 29. Теорем о главном роду в минималном}\label{hauptgeschlechtssatz-im-minimalen}

Нехај \(\mathfrak K_r=\mathfrak Z\cdot\mathfrak G\) буде прєкрижены продукт, а \(\mathfrak G^*\) група всих елементов, кторе посрєдством внутрных автоморфизмов прєводет \(\mathfrak Z\) како цєлост в себе (\(\mathfrak Z^*u_E,\mathfrak Z^*u_S\), --- записано в побочных класах).

\noindent\textbf{Дефиниција.} „\emph{Главны род в минималном}“ \(\mathscr H\) состоји се из всих груповых автоморфизмов од \(\mathfrak G^*\), кторе сут продолжење идентитета на \(\mathfrak Z^*\). Разлог за то назвање слєдује из специализације на цикличне поља. Види § 30.

\noindent\textbf{Теорем.} \emph{Всакы автоморфизм главного рода имаје форму \(u_S\mapsto u_S\cdot b^Sb^{-1}=u_Sb^{S-1}\), кде \(b\) јест в \(\mathfrak Z\), и обратно, всако тако приреджење стварјаје автоморфизм главного рода.}

\noindent\emph{Доказ.} Нехај при груповом автоморфизму буде \(u_S\mapsto v_S\), \(z\mapsto z\), за всако \(z\) из \(\mathfrak Z\), и нехај при том продукты одповєдајут једни другим; тогда автоморфизм можно дополнити до автоморфизма прстена \(\mathfrak K_r\) посрєдством приреджења
\[
\sum u_Sc_{\lambda_S}\mapsto\sum v_Sc_{\lambda_S},
\]
при чем \(c_{\lambda_S}\) сут в \(\mathfrak Z\). Але всакы такы автоморфизм јест, како знано, внутрны.

Нехај \(b\) буде стварјајучи елемент, зато \(v_S=bu_Sb^{-1}\), \(z=bzb^{-1}\); зато же \(b\) јест комутативны со всими \(z\), он лежи в \(\mathfrak Z\), и зато добыва се
\[
v_S=bu_Sb^{-1}=u_Sb^Sb^{-1}.
\]
Обратно, ако положимо \(v_S=u_Sb^Sb^{-1}\) с \(b\) в \(\mathfrak Z\), то добыва се \(v_S=bu_Sb^{-1}\); иде зато о автоморфизм од \(\mathfrak K_r\), кторы прєводи \(\mathfrak Z\) по елементах в себе, а тым по дефиницији прєводи такоже \(\mathfrak G^*\) како цєлост в себе; он јест зато автоморфизм главного рода.
% END BOOK_S29

% BEGIN BOOK_S30
\subsection*{§ 30. Специализација на цикличне разпадне поља}\label{spezialisierung-auf-zyklische-zerfaellungskoerper}

Ако прєкрижены продукт образујемо особливо с цикличным пољем \(\mathfrak Z\), приходимо к толковању знаных нормных теоремов. Најпрво имаје мєсто

\noindent\textbf{Теорем 1.} \emph{Група \(\mathscr K_{\mathfrak Z}\) од \(\{\mathfrak K\}\) с установјеным цикличным разпадным пољем \(\mathfrak Z\) над \(P\) јест изоморфна с групоју \(P^*/N\mathfrak Z^*\), кде \(N\mathfrak Z^*\) пробєгаје групу всих норм \(Nz\) елементов \(z\neq0\) из \(\mathfrak Z\).}

\noindent\emph{Доказ.} Доказ основываје се на нормовању прєкриженого прєдставјења и на нєколико помочных лемах о том.

\noindent\textbf{Дефиниција.} Циклично прєкрижено прєдставјење называје се \emph{нормованым}, ако имаје форму
\[
\mathfrak K_r=\mathfrak Z+\mathfrak Zu+\cdots+\mathfrak Zu^{n-1},
\]
с \(zu=uz^S\), \(u^n=a\), кде \(S\) јест стварјајучи елемент из \(\mathfrak G\) (ако зато положимо \(u_{S^r}=u^r\), \(r=0,\ldots,n-1\), вмєсто обчего \(u_{S^r}=b_ru^r\) с \(b_r\) в \(\mathfrak Z\)).

\noindent\textbf{Помочна лема 1.} \emph{Циклично прєкрижено прєдставјење можно всегда взети нормованым. При том \(a\) лежи в \(P\), и всако \(a\) из \(P\) даваје тако прєдставјење.}

\noindent\emph{Доказ.} Прва чест јест јасна, зато же из \(u^{-1}zu=z^S\) слєдује \(u^{-r}zu^r=z^{S^r}\), зато \(u^r\) стварјаје субституцију \(z\mapsto z^{S^r}\). Даље, зато же \(u^n\) индуцираје идентитет \(z\mapsto z^{S^n}\), добыва се \(u^n=a\) с \(a\) в \(\mathfrak Z\). Але \(u^n\) јест комутативны со всими \(u^r\), то само имаје мєсто за \(a\), и зато \(a^{S^r}=a\); зато \(a\) лежи в \(P\).

зато \(a\) јест једина сучствена константа множења. Зато условја асоциативности не имајут содржање; всако \(a\) из \(P\) даваје асоциативны систем, а тым и неко \(\mathfrak K_r\).

\noindent\textbf{Помочна лема 2.} \emph{Вмєсто полного класа \(\{a_{S,T}\}\) асоциованых факторных системов достаточно јест разгледати нормоване факторне системы, кторе сут содрживане в класу --- т.~ј. кторе наставајут из нормованых прєдставјениј. Ту подмножину класа называмо нормованым класом. Група полных клас \(\{a_{S,T}\}\) јест изоморфна с групоју нормованых клас.}

\noindent\emph{Доказ.} Зато же по помочној леми 1 ниједин нормованы клас не јест праздны, приреджење јест зато једнозначно. Операција оставаје иста, зато же јест опрєдєљена посрєдством либовољного прєдставитеља.

\noindent\textbf{Помочна лема 3.} \emph{Група нормованых факторных системов јест изоморфна с мултипликативној групоју \(P^*\), а група нормованых клас јест изоморфна с \(P^*/N\mathfrak Z^*\).}

\noindent\emph{Доказ.} Из помочној лемы 1 слєдује, же всакы нормованы факторны систем имаје вид \([1,1,\ldots,a,a,\ldots]\), зато же
\[
u^\nu\cdot u^\mu=u^{\nu+\mu}\quad(\nu+\mu<n),
\qquad
u^\nu\cdot u^\mu=au^\varrho\quad(\nu+\mu=n+\varrho,\ 0\leq\varrho\leq n-2).
\]
Але приреджење
\[
[1,\ldots,a,a,\ldots]\mapsto a,
\qquad
[1,1,\ldots,b,b,\ldots]\mapsto b,
\]
\[
[1,\ldots,ab,\ldots]\mapsto ab,
\]
јест груповы изоморфизм на \(P^*\), зато же, поновно по помочној леми 1, \(a\) пробєгаје все елементы из \(P^*\). --- Два асоциоване нормоване факторне системы наставајут једен из другого посрєдством трансформације \(v=uc\) (\(c\) в \(\mathfrak Z\)). То але даваје
\[
v=uc,\qquad v^2=u^2\cdot c^S\cdot c;
\]
\[
v^r=u^r c^{S^{r-1}}c^{S^{r-2}}\cdots c;
\qquad
v^n=u^nN(c).
\]
Из \(u^n=a\) зато слєдује \(v^n=aN(c)\). Нормованы клас состоји се зато из факторных системов \([1,\ldots,aN(c),aN(c),\ldots]\), кде \(c\) пробєгаје все елементы \(\neq0\) из \(\mathfrak Z\). Тым јест доказана друга чест помочној лемы.

Из помочных лем 2 и 3 слєдује горњи теорем 1, с обзиром на изоморфизм тамошњих \(\{\mathfrak K\}\) с класами \(\{a_{S,T}\}\).

\noindent\textbf{Теорем 2.} \emph{Група главного рода в минималном јест изоморфна с групоју тых елементов \(c\) из \(\mathfrak Z\), за кторе \(N(c)=1\). Из \(N(c)=1\) зато слєдује обстајање некого \(b\neq0\) из \(\mathfrak Z\), тако же \(c=b^S\cdot b^{-1}=b^{S-1}\) (симболична \(S-1\)-та потенција).}

\noindent\emph{Доказ.} При доказу помочној лемы 3 достало се: \(N(c)=1\) јест идентично с тым, же трансформација \(v=uc\) прєдставја автоморфизм главного рода; различным \(c\) одповєдајут различне автоморфизмы и обратно, а продукту одповєдаје продукт.

\noindent\textbf{Примєтка.} Теорем 2 оправдава назвање главны род. По смыслу факторгрупу
\[
\mathfrak Z^*/\mathscr H
\quad(\mathscr H=\text{главны род})
\cong\mathfrak Z^*/\mathfrak Z^{*S-1}
\]
можно назвати групоју родов в минималном; за ту групу имаје мєсто, же она јест изоморфна с \(N(\mathfrak Z^*)\), зато в свези с теоремом 1,
\[
\begin{aligned}
\text{група родов в минималном}
&=\mathfrak Z^*/\mathscr H
 \cong\mathfrak Z^*/\mathfrak Z^{*S-1}
 \cong N(\mathfrak Z^*)\\
&=\text{група од }\{\mathfrak K\}\text{ к }\mathfrak Z
 =\mathscr K_{\mathfrak Z}
 \cong P^*/N(\mathfrak Z^*).
\end{aligned}
\]
% END BOOK_S30

% BEGIN BOOK_S31
\subsection*{§ 31. Примєњења цикличного специалного случаја}\label{anwendungen-des-zyklischen-spezialfalles}

Циклично прєкрижено прєдставјење даваје нове доказы за теоремы доказане в § 20: же всако конечно тєло јест комутативно и же кромє кватернионов не обстаја никако право некомутативно разширјење поља реалных чисел. --- Те доказы можно бы было назвати доказами теорије класовых пољ в минималном.

\noindent\textbf{1. Всако тєло из конечного числа елементов јест комутативно (теорем 10 § 20).}

Доказ основываје се на двох знаных фактах о комутативных конечных пољах (Галоисовых пољах).

\noindent\textbf{1.} \emph{Галоисове поља сут цикличне над всакым својим подпољем.}

\noindent\textbf{2.} \emph{Ако \(\mathfrak Z\) јест Галоисово поље, то \(\mathfrak Z\) јест циклично над \(P\), то всакы елемент из \(P\) јест норма некого елемента од \(\mathfrak Z\). \(P^*=N(\mathfrak Z^*)\).}

Нехај нынє \(\mathfrak K\) буде конечно, \(P\) јего центар, а \(\mathfrak Z\) неко циклично разпадно поље; тогда (§ 30 конец), меджу другым, група \(\mathscr K_{\mathfrak Z}\) од \(\{\mathfrak K\}\) к \(\mathfrak Z\) јест изоморфна с \(P^*/N\mathfrak Z^*\); зато по 2. она јест равна јединици \(\{P\}\). То але значи, же не обстаја од \(P\) различно тєло с \(P\) како центром, \(\mathfrak K=P\).

\noindent\textbf{2. Једино право некомутативно разширјење над пољем \(P\) реалных чисел сут кватернионы (теорем 9 § 20).}

\noindent\emph{Доказ.} \(P\) мора быти центар, зато же \(P\) имаје како комутативно разширјење само поље \(\mathfrak Z\) комплексных чисел, над кторым како центром, зато же \(\mathfrak Z\) јест алгебраично затворјено, не обстаја некомутативно разширјење тєла. Једночасно само \(\mathfrak Z\) приходи в обзир како разпадно поље. И \(\mathfrak Z\) јест циклично над \(P\). Даље, всако позитивно число јест норма, зато група \(P^*/N\mathfrak Z^*\) имаје ступењ 2; зато кромє \(\{P\}\) обстаја само једен клас \(\{\mathfrak K\}\), а то сут кватернионы. Једночасно добыва се нормална форма, ако взмемо \(-1\) како нормованы факторны систем в класу асоциованых системов. Зато же тогда добыва се \(\mathfrak K=\mathfrak Z+\mathfrak Zu\), с \(u^2=-1\). зато \(\mathfrak K=P+Pi+(P+Pi)u\) с \(i^2=-1\), \(u^2=-1\).

\noindent\textbf{Примєтка к 1.} Такоже доказ знаного факта 2 можно на основє конца § 30 водити тако: Добыва се \(\mathscr H=\mathfrak Z^{*S-1}\cong\mathfrak Z^*/P^*\), зато же все и само елементи од \(P^*\) стварјајут идентичны автоморфизм главного рода, или такоже при примєњењу од \(S-1\) прєходет в јединицу. зато при конечном \(\mathfrak Z\) число елементов од \(\mathscr H\) јест равно индексу од \(P^*\) в \(\mathfrak Z^*\). Зато число елементов од \(\mathfrak Z^*/\mathscr H\) јест равно числу елементов од \(P^*\); то само зато имаје мєсто за \(N(\mathfrak Z^*)=\mathfrak Z^*/\mathscr H\), и зато \(P^*=N(\mathfrak Z^*)\). (Сравни: число родов = число амбигвных клас. Тому числово-теоретичному начину изражања але всуде трєба додати „в минималном“.)

\noindent\textbf{3.} \emph{Ако \(P\) јест \(p\)-адично основно поље, а \(\mathfrak Z\) јест циклично \(n\)-того ступења над \(P\), то група \(\mathscr K_{\mathfrak Z}\) од \(\{\mathfrak K\}\) к \(\mathfrak Z\) јест изоморфна с Галоисовој групоју од \(\mathfrak Z\).}

То јест по концу § 30 прєма послєдица „теорије класовых пољ в малом“, ктора изрека
\[
P^*/N\mathfrak Z^*\cong\mathfrak G.
\]

\noindent\textbf{Примєтка к 3.} Слєдује, же груповы ред елементов из \(\mathscr K_{\mathfrak Z}\), то јест експонент од \(\{\mathfrak K\}\), поставаје равны јего индексу. Зато же експонент стварјајучего елемента од \(\mathscr K_{\mathfrak Z}\) поставаје \(=n\), зато такоже индекс како дєлитељ од \(n\) и многократник експонента; група \(\mathscr K_{\mathfrak Z}\) содрживаје зато само класы, кторых индекс јест дєлитељ од \(n\), и за кторе \(\mathfrak Z\) јест разпадно поље. Надто H. Hasse показал, Math. Ann. 104, же \(\mathscr K_{\mathfrak Z}\) состоји се из всих клас \(\{\mathfrak K\}\), за кторе индекс јест дєлитељ од \(n\). Всако циклично разпадно поље \(\mathfrak Z\) \(n\)-того ступења над \(p\)-адичным \(P\) води зато к тој самој групє \(\mathscr K_{\mathfrak Z}\) клас тєл \(\{\mathfrak K\}\).

Два факта:

\noindent\textbf{1.} група всих \(\{\mathfrak K\}\) над \(P\), кторых индекс јест дєлитељ од \(n\), јест циклична \(n\)-того ступења.

\noindent\textbf{2.} Всако циклично разпадно поље \(n\)-того ступења над \(P\) стварјаје полну групу од \(\{\mathfrak K\}\).

можно разгледати сучственым содржањем обратного теорема теорије класовых пољ в малом.


\clearpage
% END BOOK_S31

\end{document}
